% La salud y la prevención de las enfermedades — Espagnol 1ère A — Cours signé OnBuch+
% Programme officiel MINESEC. Compiler avec : tectonic EA06-salud-prevencion-enfermedades.tex
\newcommand{\DOCMATIERE}{Espagnol}
\newcommand{\DOCNIVEAU}{1ère A}
\newcommand{\DOCMODULE}{Módulo 2 : Environnement, santé et bien-être}
\newcommand{\DOCLECON}{EA06}
\newcommand{\DOCTITRE}{La salud y la prevención de las enfermedades}
% OnBuch+ — préambule « cours signés » ENRICHI (compilé avec tectonic / XeLaTeX)
% Système visuel « Pop » d'OnBuch+ : fond crème (jamais blanc), bordures encre
% épaisses, ombres « dures » décalées, titres Archivo Black en capitales.
% Version MATHÉMATIQUES : graphes (pgfplots/tikz), boîtes pédagogiques
% (définition, propriété, méthode, attention, le savais-tu, exercice résolu,
% exercice, TP, bilan), page de garde et signature OnBuch+ ancrée en bas.
%
% Macros attendues dans main.tex AVANT \input{preamble} :
%   \DOCMATIERE  \DOCNIVEAU  \DOCMODULE  \DOCLECON (numéro)  \DOCTITRE
\documentclass[11pt]{article}

\usepackage{fontspec}
\usepackage[spanish,french]{babel} % français = langue principale ; l'espagnol via \esp{..} / espagnol
\usepackage[a4paper,margin=2cm,top=3.4cm,bottom=2.8cm,headheight=1.4cm,footskip=1.3cm]{geometry}
\usepackage{amsmath,amssymb,amsfonts}
\usepackage{mathtools} % \xmapsto, \coloneqq... souvent utilisés par les modèles
\usepackage{cancel} % simplifications barrées (bilans de forces)
% \abs{z} (module, valeur absolue) et \norm{u} (norme), utilisés par les modèles
\providecommand{\abs}{}\let\abs\relax\DeclarePairedDelimiter{\abs}{\lvert}{\rvert}
\providecommand{\norm}{}\let\norm\relax\DeclarePairedDelimiter{\norm}{\lVert}{\rVert}
\usepackage[table]{xcolor} % [table] : fournit aussi \rowcolors (alternance de lignes)
\usepackage{pagecolor}
\usepackage{tikz}
\usetikzlibrary{arrows.meta,positioning,calc,shapes.geometric,shapes.misc,shapes.arrows,decorations.pathmorphing,decorations.pathreplacing,decorations.markings,patterns,shadows,mindmap,trees,fit,backgrounds,matrix,intersections,angles,quotes}
\usepackage{pgfplots}
\usepackage[version=4]{mhchem} % équations nucléaires \ce{^{235}_{92}U}
\usepackage[european,siunitx]{circuitikz} % schémas électriques
\pgfplotsset{compat=1.18}
\usepgfplotslibrary{fillbetween,groupplots} % aires entre courbes (calcul intégral)
\usepackage{enumitem}
\usepackage{fancyhdr}
% [most] charge toutes les bibliothèques tcolorbox (listings, minted, raster,
% documentation...) et gonfle inutilement la mémoire TeX sur un document long
% (~300 boîtes) : on ne charge que ce qui est réellement utilisé.
\usepackage[skins,breakable]{tcolorbox}
\usepackage{graphicx}
\usepackage{colortbl}
\usepackage{booktabs}
\usepackage{tabularx}
\usepackage{diagbox} % case d'en-tête diagonale des tableaux à double entrée
% Colonnes extensibles centrée / gauche / droite (C, L, R), souvent utilisées par les modèles
\newcolumntype{C}{>{\centering\arraybackslash}X}
\newcolumntype{L}{>{\raggedright\arraybackslash}X}
\newcolumntype{R}{>{\raggedleft\arraybackslash}X}
\usepackage{array}
\usepackage{multicol}
\usepackage{siunitx}
% parse-numbers=false : accepte aussi \SI{2,5 \times 10^{-3}}{...}
\sisetup{locale=FR,output-decimal-marker={,},group-minimum-digits=4,parse-numbers=false}
\DeclareSIUnit\ohm{\ensuremath{\symup{\Omega}}} % U+2126 absent de la police
\DeclareSIUnit\division{div} % réglages d'oscilloscope (ms/div, V/div)
\DeclareSIUnit\tr{tr} % tours (vitesse de rotation, ex. \SI{1500}{\tr\per\minute})
\DeclareSIUnit\molar{M} % concentration molaire (ex. \SI{3}{\molar}, \SI{1}{\micro\molar})
\DeclareSIUnit\an{an} % année (le modèle confond parfois avec \year, un compteur TeX) — ex. \SI{5}{\kilo\meter\per\an}
\DeclareSIUnit\FCFA{FCFA} % franc CFA (ex. \SI{500}{\FCFA\per\kilo\gram})
\DeclareSIUnit\degreeBrix{\textdegree Brix} % teneur en sucre d'un sirop/jus (réfractométrie)
\DeclareSIUnit\month{mois} % le modèle utilise parfois \month, un compteur TeX, comme unité de durée
\usepackage{qrcode}
% \degree : commande fréquemment utilisée par les modèles pour un angle ou une
% température en dehors de \SI{}{} (non fournie par les paquets chargés).
\providecommand{\degree}{^{\circ}}
% \qed (fin de démonstration), utilisé par les modèles sans amsthm
\providecommand{\qed}{\hfill$\square$}
% \barre : double barre des valeurs interdites dans un tableau de variations (tkz-tab)
\providecommand{\barre}{\ensuremath{\|}}
% environnement proof (amsthm non chargé) utilisé par les modèles
\ifdefined\proof\else\newenvironment{proof}[1][Démonstration]{\par\noindent\textit{#1.}\ }{\qed\par}\fi
\usepackage{lastpage}
\usepackage{hyperref}
\hypersetup{hidelinks}

% ============================================================
% Palette « Pop » OnBuch+ — mêmes hex que la classe P (plus_ui.dart)
% ============================================================
\definecolor{poporange}{HTML}{F59321}
\definecolor{popdark}{HTML}{E07A0C}
\definecolor{popcream}{HTML}{FFF6E8}
\definecolor{popink}{HTML}{1C1714}
\definecolor{poppurple}{HTML}{7A5AE0}
\definecolor{popgreen}{HTML}{2E9E6A}
\definecolor{poppink}{HTML}{E0457B}
\definecolor{popblue}{HTML}{2F7DE1}
\definecolor{popgold}{HTML}{C98A12}
\definecolor{popmuted}{HTML}{8A7F76}
\definecolor{popline}{HTML}{EADFD2}
\definecolor{popgray}{HTML}{8A7F76}
\colorlet{popgrey}{popgray}
% Alias tolérants (le modèle nomme parfois les couleurs différemment)
\colorlet{popwhite}{white}
\colorlet{popblack}{popink}
\colorlet{popred}{poppink}
\colorlet{popyellow}{popgold}
% Teintes claires (fonds de tableaux/encadrés) — le modèle nomme parfois
% les couleurs avec un suffixe L (ex. popblueL) sans les avoir définies.
\colorlet{poporangeL}{poporange!20}
\colorlet{popdarkL}{popdark!20}
\colorlet{poppurpleL}{poppurple!20}
\colorlet{popgreenL}{popgreen!20}
\colorlet{poppinkL}{poppink!20}
\colorlet{popblueL}{popblue!20}
\colorlet{popgoldL}{popgold!20}
\colorlet{popredL}{popred!20}
\colorlet{popgrayL}{popgray!20}
% Teintes claires pour fonds de tableaux / zones
\colorlet{popblueL}{popblue!10!white}
\colorlet{poporangeL}{poporange!14!white}
\colorlet{popgreenL}{popgreen!12!white}
\colorlet{poppurpleL}{poppurple!12!white}
\colorlet{poppinkL}{poppink!12!white}
\colorlet{popgoldL}{popgold!14!white}

\pagecolor{popcream}

% ============================================================
% Typographie
% ============================================================
\defaultfontfeatures{Ligatures=TeX}
\setmainfont{PlusJakartaSans-Medium.ttf}[Path=../../../fonts/,BoldFont=PlusJakartaSans-Bold.ttf]
% Maths en sans-serif (Fira Math) avec chiffres et lettres droites de Plus
% Jakarta, pour que les formules se fondent dans le texte.
\usepackage[math-style=ISO,bold-style=ISO]{unicode-math}
\setmathfont{FiraMath-Regular.otf}[Path=../../../fonts/]
\setmathfont{PlusJakartaSans-Medium.ttf}[Path=../../../fonts/,range={up/num,up/{Latin,latin},\mathup}]
\usepackage{icomma} % virgule décimale sans espace en mode math
% Caractères mathématiques Unicode absents des polices de texte (Plus Jakarta,
% Archivo Black) : rendus par leur équivalent mathématique, en texte comme en
% formule — sinon ils s'affichent en carré vide (ex. « Arithmétique dans ℤ »).
\usepackage{newunicodechar}
\newunicodechar{ℝ}{\ensuremath{\mathbb{R}}}
\newunicodechar{ℤ}{\ensuremath{\mathbb{Z}}}
\newunicodechar{ℂ}{\ensuremath{\mathbb{C}}}
\newunicodechar{ℕ}{\ensuremath{\mathbb{N}}}
\newunicodechar{ℚ}{\ensuremath{\mathbb{Q}}}
\newunicodechar{𝔻}{\ensuremath{\mathbb{D}}}
\newunicodechar{α}{\ensuremath{\alpha}}
\newunicodechar{β}{\ensuremath{\beta}}
\newunicodechar{γ}{\ensuremath{\gamma}}
\newunicodechar{δ}{\ensuremath{\delta}}
\newunicodechar{ε}{\ensuremath{\varepsilon}}
\newunicodechar{θ}{\ensuremath{\theta}}
\newunicodechar{λ}{\ensuremath{\lambda}}
\newunicodechar{μ}{\ensuremath{\mu}}
\newunicodechar{σ}{\ensuremath{\sigma}}
\newunicodechar{τ}{\ensuremath{\tau}}
\newunicodechar{φ}{\ensuremath{\varphi}}
\newunicodechar{ψ}{\ensuremath{\psi}}
\newunicodechar{ω}{\ensuremath{\omega}}
\newunicodechar{Σ}{\ensuremath{\Sigma}}
% majuscules produites par \MakeUppercase dans les titres (α-aminés -> Α)
\newunicodechar{Α}{\ensuremath{\alpha}}
\newunicodechar{Β}{\ensuremath{\beta}}
\newunicodechar{Γ}{\ensuremath{\Gamma}}
\newunicodechar{Δ}{\ensuremath{\Delta}}
\newunicodechar{Ε}{\ensuremath{\varepsilon}}
\newunicodechar{Θ}{\ensuremath{\Theta}}
\newunicodechar{Λ}{\ensuremath{\Lambda}}
\newunicodechar{Μ}{\ensuremath{\mu}}
\newunicodechar{Τ}{\ensuremath{\tau}}
\newunicodechar{Φ}{\ensuremath{\Phi}}
\newunicodechar{Ψ}{\ensuremath{\Psi}}
\newunicodechar{Ω}{\ensuremath{\Omega}}
\newunicodechar{↦}{\ensuremath{\mapsto}}
\newunicodechar{∈}{\ensuremath{\in}}
\newunicodechar{∉}{\ensuremath{\notin}}
\newunicodechar{∧}{\ensuremath{\wedge}}
\newunicodechar{⇔}{\ensuremath{\Leftrightarrow}}
\newunicodechar{⟺}{\ensuremath{\Longleftrightarrow}}
\newunicodechar{⇒}{\ensuremath{\Rightarrow}}
\newunicodechar{⟹}{\ensuremath{\Longrightarrow}}
\newunicodechar{∘}{\ensuremath{\circ}}
\newunicodechar{∪}{\ensuremath{\cup}}
\newunicodechar{∩}{\ensuremath{\cap}}
\newunicodechar{≡}{\ensuremath{\equiv}}
\newunicodechar{⊂}{\ensuremath{\subset}}
\newunicodechar{⊥}{\ensuremath{\perp}}
\newunicodechar{∃}{\ensuremath{\exists}}
\newunicodechar{∀}{\ensuremath{\forall}}
\newunicodechar{∛}{\ensuremath{\sqrt[3]{\,}}}
\newunicodechar{∜}{\ensuremath{\sqrt[4]{\,}}}
\newunicodechar{⃗}{\ensuremath{{}^{\to}}}
\newunicodechar{ⁿ}{\ifmmode^{n}\else\textsuperscript{\ensuremath{n}}\fi}
\newunicodechar{ˣ}{\ifmmode^{x}\else\textsuperscript{\ensuremath{x}}\fi}
\newunicodechar{ᵇ}{\ifmmode^{b}\else\textsuperscript{\ensuremath{b}}\fi}
\newunicodechar{ʳ}{\ifmmode^{r}\else\textsuperscript{\ensuremath{r}}\fi}
\newunicodechar{ᵘ}{\ifmmode^{u}\else\textsuperscript{\ensuremath{u}}\fi}
\newunicodechar{ʸ}{\ifmmode^{y}\else\textsuperscript{\ensuremath{y}}\fi}
\newunicodechar{ᵃ}{\ifmmode^{a}\else\textsuperscript{\ensuremath{a}}\fi}
\newunicodechar{ᵉ}{\ifmmode^{e}\else\textsuperscript{\ensuremath{e}}\fi}
\newunicodechar{ᵏ}{\ifmmode^{k}\else\textsuperscript{\ensuremath{k}}\fi}
\newunicodechar{ᵖ}{\ifmmode^{p}\else\textsuperscript{\ensuremath{p}}\fi}
\newunicodechar{ᴺ}{\ifmmode^{N}\else\textsuperscript{\ensuremath{N}}\fi}
\newunicodechar{ᶜ}{\ifmmode^{c}\else\textsuperscript{\ensuremath{c}}\fi}
\newunicodechar{ʰ}{\ifmmode^{h}\else\textsuperscript{\ensuremath{h}}\fi}
\newunicodechar{ᵠ}{\ifmmode^{q}\else\textsuperscript{\ensuremath{q}}\fi}
\newunicodechar{ₙ}{\ifmmode_{n}\else\textsubscript{\ensuremath{n}}\fi}
\newunicodechar{ₐ}{\ifmmode_{a}\else\textsubscript{\ensuremath{a}}\fi}
\newunicodechar{ᵢ}{\ifmmode_{i}\else\textsubscript{\ensuremath{i}}\fi}
\newunicodechar{ᵣ}{\ifmmode_{r}\else\textsubscript{\ensuremath{r}}\fi}
\newunicodechar{ₖ}{\ifmmode_{k}\else\textsubscript{\ensuremath{k}}\fi}
\newunicodechar{ᵦ}{\ifmmode_{\beta}\else\textsubscript{\ensuremath{\beta}}\fi}
\newunicodechar{ₓ}{\ifmmode_{x}\else\textsubscript{\ensuremath{x}}\fi}
\newfontfamily\popdisplay{ArchivoBlack-Regular.ttf}[Path=../../../fonts/]
\newfontfamily\poplogo{SpaceGrotesk-Bold.ttf}[Path=../../../fonts/]
\newfontfamily\popsemi{PlusJakartaSans-SemiBold.ttf}[Path=../../../fonts/]
\newfontfamily\popxbold{PlusJakartaSans-ExtraBold.ttf}[Path=../../../fonts/]
\newfontfamily\popxboldls{PlusJakartaSans-ExtraBold.ttf}[Path=../../../fonts/,LetterSpace=3.0]

\setlist[enumerate]{leftmargin=1.7em,itemsep=5pt,topsep=4pt}
\setlist[itemize]{leftmargin=1.7em,itemsep=4pt,topsep=4pt,label={\color{poporange}\textbullet}}
\setlength{\parindent}{0pt}
\setlength{\parskip}{5pt}
\renewcommand{\arraystretch}{1.25}

% pgfplots : style Pop par défaut
\pgfplotsset{
  every axis/.append style={axis line style={popink,line width=1pt},tick style={popink},
    grid style={popline,line width=0.5pt},label style={font=\small},tick label style={font=\footnotesize}},
  popcurve/.style={poporange,line width=1.8pt,no markers,smooth},
}
% Même style disponible dans un \draw TikZ simple (hors pgfplots)
\tikzset{popcurve/.style={poporange,line width=1.8pt}}
\tikzset{popgrid/.style={popline,very thin}}
\colorlet{popcurve}{poporange} % le nom du style est parfois utilisé comme couleur

% ============================================================
% Pill / squiggle
% ============================================================
\newcommand{\poppill}[3]{%
  \tikz[baseline=-0.55ex]\node[draw=popink,line width=1.2pt,rounded corners=2.6pt,%
    fill=#2,inner xsep=6.5pt,inner ysep=3pt]{%
    \popxboldls\fontsize{7}{7}\selectfont\color{#3}\MakeUppercase{#1}};%
}
\newcommand{\popsquiggle}[1]{%
  \tikz[baseline=-0.4ex]{
    \draw[#1,line width=2.6pt,line cap=round]
      plot [smooth] coordinates {(0,0.09) (0.34,0.01) (0.68,0.21) (1.02,0.01) (1.36,0.10)};
  }%
}
% Mot-clé surligné (marqueur orange sous le texte)
\newcommand{\cle}[1]{{\popxbold\color{popdark}#1}}

% ============================================================
% En-tête / pied
% ============================================================
\pagestyle{fancy}
\fancyhf{}
\renewcommand{\headrulewidth}{0pt}
\renewcommand{\footrulewidth}{0pt}
\lhead{\raisebox{-0.15ex}{\poplogo\fontsize{13}{13}\selectfont\color{popink}OnBuch\color{poporange}+}}
\rhead{\poppill{\DOCMATIERE\ \textbullet\ \DOCNIVEAU\ \textbullet\ Leçon \DOCLECON}{white}{popink}}
\lfoot{\popxbold\fontsize{7.6}{7.6}\selectfont\color{popmuted}\MakeUppercase{Cours signé OnBuch+}\ \textbullet\ {\popsemi\DOCTITRE}}
\rfoot{\tikz[baseline=-0.6ex]\node[rounded corners=8.5pt,fill=popink,minimum height=17pt,minimum width=17pt,inner xsep=5pt,inner sep=0]{\popxbold\fontsize{8}{8}\selectfont\color{popcream}\thepage\,/\,\pageref*{LastPage}};}

% ============================================================
% Titres
% ============================================================
\newcommand{\chaptitre}[2]{%
  \begin{tcolorbox}[enhanced,colback=poporange,colframe=popink,boxrule=2.6pt,arc=11pt,
      drop shadow=popink,left=15pt,right=15pt,top=12pt,bottom=11pt,before skip=3pt,after skip=18pt]
    \poppill{#2}{popink}{popcream}\par\vspace{9pt}
    {\raggedright\hyphenpenalty=10000\exhyphenpenalty=10000\popdisplay\fontsize{22}{24}\selectfont\color{popcream}\MakeUppercase{#1}\par}
  \end{tcolorbox}
}
% Section numérotée : pastille orange avec le numéro + titre Archivo
\newcounter{popsec}
\newcommand{\coursec}[1]{%
  \stepcounter{popsec}\setcounter{popsub}{0}%
  \par\vspace{16pt}\noindent
  \begin{minipage}{\linewidth}
  \tikz[baseline=-0.7ex]\node[draw=popink,line width=1.6pt,fill=poporange,rounded corners=4pt,minimum size=22pt,inner sep=0]{\popdisplay\fontsize{12}{12}\selectfont\color{popcream}\thepopsec};%
  \hspace{8pt}{\popdisplay\fontsize{15}{17}\selectfont\color{popink}\MakeUppercase{#1}}\par
  \vspace{4pt}\noindent{\color{poporange}\rule{36pt}{3.6pt}}
  \end{minipage}\par\nopagebreak\vspace{8pt}
}
\newcounter{popsub}
\newcommand{\courssub}[1]{%
  \stepcounter{popsub}%
  \par\vspace{9pt}\noindent{\popxbold\fontsize{11.5}{13}\selectfont\color{popink}{\color{poporange}\thepopsec.\thepopsub}\hspace{6pt}#1}\par\nopagebreak\vspace{4pt}
}

% ============================================================
% Boîtes pédagogiques — même recette : fond blanc, bordure épaisse colorée,
% ombre dure encre, badge plein. Titre optionnel [#1].
% ============================================================
\tcbset{popbase/.style={breakable,enhanced,colback=white,boxrule=2.3pt,arc=9pt,
  shadow={3pt}{-3pt}{0pt}{popink},left=14pt,right=14pt,top=11pt,bottom=11pt,before skip=11pt,after skip=15pt,
  fontupper=\popsemi\fontsize{10.6}{14.5}\selectfont\color{popink},
  fontlower=\popsemi\fontsize{10.6}{14.5}\selectfont\color{popink}}}
% Coche dessinée (le glyphe ✓ n'existe pas dans les polices du document)
\newcommand{\popcheck}[1]{\tikz[baseline=-0.5ex]\draw[#1,line width=1.6pt,line cap=round,line join=round](-0.13,0.0)--(-0.04,-0.09)--(0.13,0.1);}
\providecommand{\coche}{\popcheck{popgreen}}
\providecommand{\resultat}[1]{\par\smallskip\noindent\fcolorbox{popgreen}{popgreenL}{\parbox{\dimexpr\linewidth-2\fboxsep-2\fboxrule\relax}{\textbf{Résultat.} #1}}\par\smallskip}
\newcommand{\popboxhead}[3]{\poppill{#1}{#2}{white}\ifx\relax#3\relax\else\hspace{6pt}{\popxbold\fontsize{10.6}{12}\selectfont\color{popink}#3}\fi\par\vspace{7pt}}

\newtcolorbox{aretenir}[1][]{popbase,colframe=popgreen,before upper={\popboxhead{À retenir}{popgreen}{#1}}}
% Texte en espagnol (document de lecture, dialogue, poème, extrait) : cadre
% d'étude. Un titre optionnel (source, auteur) s'affiche dans l'en-tête.
\newtcolorbox{textebox}[1][]{popbase,colframe=popblue,before upper={\popboxhead{Texto}{popblue}{#1}\begin{otherlanguage}{spanish}},after upper={\end{otherlanguage}}}
% Marques ✓ / ✗ (forme correcte / fautive) souvent utilisées par les modèles
\providecommand{\cross}{\textcolor{poppink}{\ensuremath{\times}}}
\providecommand{\xmark}{\cross}\providecommand{\crossmark}{\cross}\providecommand{\cmark}{\ensuremath{\checkmark}}
% Ligne de réponse / justification à compléter (inventée par les modèles)
\providecommand{\justification}[1]{\par\noindent\textit{Justification~:} \dotfill\par}
% \textipa (package tipa non chargé) : la police ne couvre pas l'API -> simple passage
\providecommand{\textipa}[1]{#1}
% Soulignement (ulem non chargé) : \uline, \ul
\providecommand{\uline}[1]{\underline{#1}}\providecommand{\ul}[1]{\underline{#1}}
% Mot ou phrase en espagnol à l'intérieur d'un paragraphe français
\newcommand{\esp}[1]{\ifmmode\text{\textit{#1}}\else\begingroup\selectlanguage{spanish}\itshape #1\endgroup\fi}
\newtcolorbox{exemplebox}[1][]{popbase,colframe=poporange,before upper={\popboxhead{Exemple}{poporange}{#1}}}
\newtcolorbox{definition}[1][]{popbase,colframe=popblue,before upper={\popboxhead{Définition}{popblue}{#1}}}
\newtcolorbox{propriete}[1][]{popbase,colframe=poppurple,before upper={\popboxhead{Propriété}{poppurple}{#1}}}
\newtcolorbox{methode}[1][]{popbase,colframe=popgold,before upper={\popboxhead{Méthode}{popgold}{#1}}}
\newtcolorbox{attention}[1][]{popbase,colframe=poppink,before upper={\popboxhead{Attention}{poppink}{#1}}}
\newtcolorbox{experience}[1][]{popbase,colframe=popblue,colback=popblueL,before upper={\popboxhead{Expérience}{popblue}{#1}}}
% « Le savais-tu ? » : carte violette pleine (contexte, culture, Cameroun)
\newtcolorbox{savaistu}[1][]{popbase,colback=poppurple,colframe=popink,
  fontupper=\popsemi\fontsize{10.4}{14.2}\selectfont\color{white},
  before upper={\poppill{Le savais-tu\,?}{white}{poppurple}\ifx\relax#1\relax\else\hspace{6pt}{\popxbold\color{white}#1}\fi\par\vspace{7pt}}}
% Exercice résolu : énoncé en haut, \tcblower puis la solution sur fond crème
\newcounter{popexo}\newcounter{popexr}
\newtcolorbox{exoresolu}[1][]{popbase,colframe=poporange,colbacklower=popcream,
  segmentation style={popink,line width=1.2pt,dashed},
  before upper={\stepcounter{popexr}\popboxhead{Exercice résolu \thepopexr}{poporange}{#1}},
  before lower={\poppill{Solution}{popink}{popcream}\par\vspace{6pt}}}
% Exercice (non résolu en place) — la correction est dans la partie Corrigés
\newtcolorbox{exercice}[1][]{popbase,colframe=popink,
  before upper={\stepcounter{popexo}\popboxhead{Exercice \thepopexo}{popink}{#1}}}
% Corrigé
\newtcolorbox{corrige}[1][]{popbase,colframe=popgreen,colback=popgreenL,
  before upper={\popboxhead{Corrigé}{popgreen}{#1}}}
% Objectifs / prérequis / bilan
\newtcolorbox{objectifs}{popbase,colframe=popink,colback=poporangeL,
  before upper={\popboxhead{Objectifs de la leçon}{popink}{}}}
\newtcolorbox{prerequis}{popbase,colframe=popink,colback=popblueL,
  before upper={\popboxhead{Prérequis}{popblue}{}}}
\newtcolorbox{bilan}{popbase,colframe=popink,colback=white,boxrule=2.8pt,
  before upper={\popboxhead{Fiche bilan}{poporange}{L'essentiel à connaître pour l'examen}}}
% Figure encadrée avec légende
\newtcolorbox{popfigure}[1][]{enhanced,breakable=false,colback=white,colframe=popline,boxrule=1.4pt,arc=8pt,
  left=10pt,right=10pt,top=10pt,bottom=8pt,before skip=10pt,after skip=12pt,halign=center,
  lower separated=false,halign lower=center,
  fontlower=\popsemi\fontsize{9}{11.5}\selectfont\color{popmuted},
  #1}
\newcommand{\legende}[1]{\par\vspace{6pt}{\popsemi\fontsize{9}{11.5}\selectfont\color{popmuted}#1}}

% ============================================================
% Signature OnBuch+
% ============================================================
\newcommand{\poplogoinline}[1]{{\poplogo\fontsize{#1}{#1}\selectfont\color{popink}OnBuch\color{poporange}+}}
\newcommand{\signatureonbuch}{%
  \begin{tcolorbox}[enhanced,colback=popink,colframe=popink,boxrule=2.2pt,arc=10pt,
      drop shadow=poporange,left=14pt,right=14pt,top=10pt,bottom=10pt,before skip=0pt,after skip=0pt]
    \tikz[baseline=-0.7ex]\node[circle,fill=popgreen,draw=popcream,line width=1.3pt,minimum size=22pt,inner sep=0]{\popcheck{white}};%
    \hspace{9pt}\begin{minipage}[c]{0.62\linewidth}
      {\popxbold\fontsize{12}{13}\selectfont\color{popcream}Signé\ \textbullet\ {\poplogo OnBuch\color{poporange}+}}\par\vspace{2pt}
      {\raggedright\popsemi\fontsize{8}{10.5}\selectfont\color{popline}\MakeUppercase{Équipe pédagogique OnBuch+}\\\MakeUppercase{Conforme au programme officiel MINESEC}\par}
    \end{minipage}\hfill
    {\poplogo\fontsize{22}{22}\selectfont\color{popcream}OnBuch\color{poporange}+}
  \end{tcolorbox}%
}

% ============================================================
% Page de garde
% #1 titre · #2 matière · #3 niveau · #4 module · #5 n° leçon · #6 durée ·
% #7 description / objectifs (texte ou itemize)
% ============================================================
\newcommand{\pagegarde}[7]{%
  \thispagestyle{empty}
  \newgeometry{a4paper,margin=1.7cm,top=1.6cm,bottom=1.4cm}%
  \begin{tikzpicture}[remember picture,overlay]
    \fill[poporange!18] ([xshift=-3.2cm,yshift=-3.6cm]current page.north east) circle (4.6cm);
    \fill[poppurple!14] ([xshift=2.4cm,yshift=7.6cm]current page.south west) circle (3.8cm);
  \end{tikzpicture}%
  {\poplogo\fontsize{30}{30}\selectfont\color{popink}OnBuch\color{poporange}+}\hfill
  \raisebox{0.6ex}{\poppill{Cours signé \textbullet\ Premium}{popink}{popcream}}\par
  \vspace{1pt}\popsquiggle{poppurple}\par
  \vspace{8pt}
  \begin{tcolorbox}[enhanced,colback=poporange,colframe=popink,boxrule=2.8pt,arc=13pt,
      drop shadow=popink,left=18pt,right=18pt,top=12pt,bottom=13pt,after skip=10pt]
    \poppill{#2 \textbullet\ #3}{popink}{popcream}\hspace{5pt}\poppill{#4}{white}{popink}\par\vspace{8pt}
    {\popdisplay\fontsize{14}{14}\selectfont\color{popink}LEÇON #5}\par\vspace{4pt}
    {\raggedright\hyphenpenalty=10000\exhyphenpenalty=10000\popdisplay\fontsize{25}{27}\selectfont\color{popcream}\MakeUppercase{#1}\par}
    \vspace{8pt}
    {\popxbold\fontsize{9.5}{9.5}\selectfont\color{popink}\MakeUppercase{Durée conseillée : #6}}
  \end{tcolorbox}
  \begin{tcolorbox}[enhanced,colback=white,colframe=popink,boxrule=2.2pt,arc=10pt,
      drop shadow=popink,left=16pt,right=16pt,top=10pt,bottom=10pt,after skip=8pt]
    \poppill{Dans cette leçon}{poppurple}{white}\par\vspace{5pt}
    \popsemi\fontsize{9.3}{12.6}\selectfont\color{popink}#7
  \end{tcolorbox}
  \noindent\begin{minipage}[t]{0.487\linewidth}
    \begin{tcolorbox}[enhanced,colback=popgreen,colframe=popink,boxrule=2.2pt,arc=10pt,
        drop shadow=popink,left=11pt,right=11pt,top=10pt,bottom=10pt,height=3.1cm,valign=center]
      \tcbox[colback=white,colframe=popink,boxrule=1.3pt,arc=5pt,boxsep=0pt,nobeforeafter,
        left=3pt,right=3pt,top=3pt,bottom=3pt,valign=center]{\qrcode[height=1.9cm]{https://play.google.com/store/apps/details?id=com.luvvix.onbuch&hl=fr}}%
      \hspace{8pt}\begin{minipage}[c]{0.5\linewidth}
        {\popdisplay\fontsize{11}{12.5}\selectfont\color{white}\MakeUppercase{Télécharge OnBuch}}\par\vspace{4pt}
        {\raggedright\popsemi\fontsize{8.4}{11}\selectfont\color{white}Gratuit \textbullet\ Google Play\\Tuteur IA, annales, concours, résultats\par}
      \end{minipage}
    \end{tcolorbox}
  \end{minipage}\hfill
  \begin{minipage}[t]{0.487\linewidth}
    \begin{tcolorbox}[enhanced,colback=poppurple,colframe=popink,boxrule=2.2pt,arc=10pt,
        drop shadow=popink,left=11pt,right=11pt,top=10pt,bottom=10pt,height=3.1cm,valign=center]
      \tikz[baseline=-0.7ex]\node[circle,fill=white,draw=popink,line width=1.3pt,minimum size=1.6cm,inner sep=0]{\popdisplay\fontsize{24}{24}\selectfont\color{poppurple}+};%
      \hspace{8pt}\begin{minipage}[c]{0.6\linewidth}
        {\popdisplay\fontsize{11}{12.5}\selectfont\color{white}\MakeUppercase{Passe à OnBuch+}}\par\vspace{4pt}
        {\raggedright\popsemi\fontsize{8.4}{11}\selectfont\color{white}Tous les cours signés, Tuteur IA illimité, fiches PDF — onglet OnBuch+ de l'app\par}
      \end{minipage}
    \end{tcolorbox}
  \end{minipage}
  \vspace{8pt}
  \popcontenu
  \vfill
  \signatureonbuch
  \restoregeometry
  \newpage
}

% Rangée « Ce cours contient » (page de garde)
\newcommand{\poptile}[3]{% #1 couleur #2 gros texte #3 légende
  \begin{tcolorbox}[enhanced,colback=#1,colframe=popink,boxrule=2pt,arc=9pt,shadow={2.5pt}{-2.5pt}{0pt}{popink},
      left=6pt,right=6pt,top=9pt,bottom=9pt,halign=center,valign=center,height=1.75cm,width=\linewidth,nobeforeafter]
    {\popdisplay\fontsize{12.5}{13}\selectfont\color{white}\MakeUppercase{#2}}\par\vspace{3pt}
    {\centering\popsemi\fontsize{7.8}{9.5}\selectfont\color{white}#3\par}
  \end{tcolorbox}}
\newcommand{\popcontenu}{%
  \par\noindent{\popxboldls\fontsize{7.5}{7.5}\selectfont\color{popmuted}CE COURS SIGNÉ CONTIENT}\par\vspace{7pt}
  \noindent\begin{minipage}[t]{0.235\linewidth}\poptile{popblue}{Cours}{complet, illustré, pas à pas}\end{minipage}\hfill
  \begin{minipage}[t]{0.235\linewidth}\poptile{poporange}{Méthodes}{et exercices résolus}\end{minipage}\hfill
  \begin{minipage}[t]{0.235\linewidth}\poptile{poppink}{Exercices}{type examen + corrigés}\end{minipage}\hfill
  \begin{minipage}[t]{0.235\linewidth}\poptile{popgreen}{Fiche bilan}{l'essentiel à retenir}\end{minipage}\par}

% Sommaire
\newtcolorbox{sommaire}{enhanced,colback=white,colframe=popink,boxrule=2.2pt,arc=10pt,
  drop shadow=popink,left=18pt,right=18pt,top=13pt,bottom=13pt,before skip=6pt,after skip=20pt,
  before upper={\poppill{Sommaire}{poppurple}{white}\par\vspace{9pt}\popsemi\fontsize{10.6}{14.5}\selectfont\color{popink}}}

% Signature de fin de cours — poussée tout en bas de la dernière page
\newcommand{\findecours}{%
  \par\vfill\nopagebreak
  \begin{center}\popsquiggle{poporange}\end{center}\vspace{4pt}
  \signatureonbuch
}
\AtBeginDocument{\def\llbracket{\mathopen{[\mkern-3.5mu[}}\def\rrbracket{\mathclose{]\mkern-3.5mu]}}} % intervalles d'entiers (glyphe ⟦ absent de la police)
% Glyphes absents de Fira Math : redéfinis à partir de glyphes disponibles
\AtBeginDocument{%
  \def\setminus{\mathbin{\backslash}}\let\smallsetminus\setminus
  \def\vdots{\mathord{\vbox{\baselineskip3.2pt\lineskiplimit0pt\kern4pt\hbox{.}\hbox{.}\hbox{.}}}}%
  \def\ddots{\mathinner{\mkern1mu\raise6.5pt\hbox{.}\mkern2mu\raise3.6pt\hbox{.}\mkern2mu\raise0.7pt\hbox{.}\mkern1mu}}%
  \def\triangle{\mathord{\scalebox{0.9}{$\symup{\Delta}$}}}%
  \def\nabla{\mathord{\raisebox{\depth}{\scalebox{1}[-1]{$\symup{\Delta}$}}}}%
  \def\checkmark{\popcheck{popdark}}%
  \def\star{\mathbin{\ast}}\def\bigstar{\mathord{\ast}}%
  \def\diamond{\mathbin{\scalebox{0.6}{$\Diamond$}}}%
  \def\bigcup{\mathop{\mathchoice{\raisebox{-0.3em}{\scalebox{1.7}{$\cup$}}}{\scalebox{1.25}{$\cup$}}{\cup}{\cup}}\displaylimits}%
  \def\bigcap{\mathop{\mathchoice{\raisebox{-0.3em}{\scalebox{1.7}{$\cap$}}}{\scalebox{1.25}{$\cap$}}{\cap}{\cap}}\displaylimits}%
  \def\lesseqqgtr{\mathrel{\lessgtr}}\def\bigoplus{\mathop{\oplus}\displaylimits}%
}
\providecommand{\ssi}{si et seulement si} % abréviation fréquente
\AtBeginDocument{\providecommand{\wideparen}{\overparen}} % arc de cercle

%% FIGFIX-BEGIN v3 — correctifs globaux des figures (docs/onbuch-generation/tools/figfix)
\usepackage{adjustbox}\usepackage{etoolbox}
% toute figure TikZ/pgfplots plus large que la ligne est réduite à la largeur disponible
\BeforeBeginEnvironment{tikzpicture}{\begin{adjustbox}{max width=\linewidth}}
\AfterEndEnvironment{tikzpicture}{\end{adjustbox}}
% légendes pgfplots lisibles par-dessus les courbes
\pgfplotsset{every axis/.append style={legend style={fill=white,fill opacity=0.92,text opacity=1,draw=black!15,inner sep=3pt}}}
% étiquettes d'arêtes (syntaxe edge["..."]) avec fond blanc
\tikzset{every edge quotes/.append style={fill=white,inner sep=1.2pt}}
% bulles de cartes mentales : texte à la ligne dans la bulle, bulle un peu plus grande
\tikzset{every concept/.append style={text width=2.0cm,align=center,minimum size=2.3cm,inner sep=2pt}}
%% FIGFIX-END
\begin{document}

\pagegarde{La salud y la prevención de las enfermedades}{Espagnol}{1ère A}{Módulo 2 : Environnement, santé et bien-être}{EA06}{2 h}{Cette leçon mixte du module « Environnement, santé et bien-être » permet de lire et commenter un texte sur les habitudes de vie favorables à la santé, d'acquérir le lexique de la santé et de la maladie, et de maîtriser l'impératif et le subjonctif pour formuler des conseils et des messages de prévention contre le paludisme, le choléra et le VIH/sida.
\vspace{4pt}
\begin{multicols}{2}\small
\begin{itemize}
\item À la fin de cette leçon, je sais lire et comprendre un texte espagnol sur la santé et les habitudes de vie
\item À la fin de cette leçon, je sais identifier et utiliser le lexique de la santé, de la maladie et de la prévention
\item À la fin de cette leçon, je sais analyser des habitudes de vie favorables ou défavorables à la santé
\item À la fin de cette leçon, je sais conjuguer et employer l'impératif affirmatif et négatif pour donner des conseils
\item À la fin de cette leçon, je sais employer le subjonctif présent après « es necesario que » et « conviene que »
\item À la fin de cette leçon, je sais distinguer les structures à infinitif et les structures à subjonctif selon le sujet
\end{itemize}\end{multicols}}

\begin{sommaire}
\begin{multicols}{2}
\begin{enumerate}
\item Texte support : Vivir sano
\item Lexique thématique de la santé
\item Grammaire 1 : l'impératif pour conseiller
\item Grammaire 2 : le subjonctif avec « es necesario que » et « conviene que »
\item Méthode du commentaire et commentaire modèle
\item Production guidée : messages de prévention
\item Activité d'intégration
\item Méthodes et exercices résolus
\item Exercices
\item Corrigés des exercices
\item Fiche bilan
\end{enumerate}
\end{multicols}
\end{sommaire}

\begin{objectifs}
\begin{itemize}
\item À la fin de cette leçon, je sais lire et comprendre un texte espagnol sur la santé et les habitudes de vie
\item À la fin de cette leçon, je sais identifier et utiliser le lexique de la santé, de la maladie et de la prévention
\item À la fin de cette leçon, je sais analyser des habitudes de vie favorables ou défavorables à la santé
\item À la fin de cette leçon, je sais conjuguer et employer l'impératif affirmatif et négatif pour donner des conseils
\item À la fin de cette leçon, je sais employer le subjonctif présent après « es necesario que » et « conviene que »
\item À la fin de cette leçon, je sais distinguer les structures à infinitif et les structures à subjonctif selon le sujet
\item À la fin de cette leçon, je sais rédiger un message de prévention clair et correct en espagnol
\item À la fin de cette leçon, je sais traduire des phrases de conseil du français vers l'espagnol et inversement
\end{itemize}
\end{objectifs}

\begin{prerequis}
\begin{itemize}
\item Présent de l'indicatif des verbes réguliers et irréguliers courants (Seconde)
\item Lexique de base du corps humain et de la vie quotidienne (Seconde)
\item Notion de mode (indicatif / impératif) vue en Seconde
\item Pronoms personnels sujets et accentuation (Seconde)
\item Structures impersonnelles simples : hay que, es importante (début d'année)
\end{itemize}
\end{prerequis}

\newpage

\coursec{Texte support : Vivir sano}

C'est la saison des pluies à Yaoundé. L'eau stagne dans les caniveaux du quartier, les moustiques pullulent dès la tombée de la nuit, et le constat est inquiétant : plusieurs élèves de ta classe sont absents, victimes du paludisme. À l'infirmerie du lycée, l'infirmière déborde de travail. Alors le délégué de santé de ton établissement te sollicite : il faut rédiger des affiches de prévention... en espagnol ! Pourquoi en espagnol ? Parce qu'une délégation d'élèves camerounais hispanisants part bientôt en échange scolaire en Guinée équatoriale, pays hispanophone voisin, et que les messages de santé publique doivent être compris par tous. Comment conseiller, interdire, recommander en espagnol ? Comment parler des symptômes, des vaccins et des bonnes habitudes de vie ? Cette leçon te donne tous les outils pour devenir un vrai relais de santé publique. Et tout commence par la lecture attentive d'un texte.

\textbf{Problématique :} quelles habitudes de vie permettent de prévenir les maladies, et comment un texte en espagnol peut-il nous convaincre d'adopter ces habitudes ?

\courssub{Découverte : lecture du texte}

\begin{methode}[Comment lire un texte de santé en espagnol]
\begin{enumerate}
\item \textbf{Lis le titre et anticipe} : le titre \esp{Vivir sano : un reto para los jóvenes cameruneses} annonce le sujet (la santé) et le public visé (les jeunes Camerounais). Avant de lire, demande-toi : quelles habitudes le texte va-t-il probablement citer ?
\item \textbf{Repère les mots transparents} : ce sont les mots qui ressemblent au français, comme \esp{prevención}, \esp{higiene}, \esp{vacuna}, \esp{alimentación}, \esp{ejercicio}. Ils te donnent le sens général sans dictionnaire.
\item \textbf{Souligne les conseils} : dans un texte de prévention, les conseils sont souvent à l'impératif (\esp{lávate}, \esp{duerme}) ou introduits par \esp{es necesario que} et \esp{conviene que}.
\item \textbf{Résume en une phrase} : après la lecture, formule l'idée principale en une seule phrase, en espagnol si possible.
\end{enumerate}
\end{methode}

Lis maintenant le texte à voix haute, en respectant la prononciation espagnole : le \esp{j} se prononce « r » grasseyé (\esp{jóvenes} : « rô-bé-nès »), le \esp{ñ} se prononce « gn » (\esp{niños} : « ni-gnos »), et le \esp{h} est toujours muet (\esp{hábitos} : « a-bi-tos »).

\begin{textebox}[Vivir sano : un reto para los jóvenes cameruneses]
En Camerún, muchas enfermedades como el paludismo o el cólera se pueden prevenir con hábitos sencillos. Sin embargo, cada año, durante la temporada de lluvias, muchos jóvenes caen enfermos y faltan a la escuela. ¿Por qué? A menudo, la causa son las malas costumbres : dormir sin mosquitera, beber agua no potable o comer siempre lo mismo.

Sin embargo, vivir sano no es un sueño imposible : es un esfuerzo de todos los días. Dormir bajo una mosquitera, lavarse las manos con jabón, beber agua potable y comer de forma equilibrada son gestos que salvan vidas. En el mercado de Yaoundé, las frutas y las verduras no son caras : ¡aprovecha! Una alimentación rica en frutas, en pescado y en cereales fortalece el cuerpo contra las infecciones.

Los médicos repiten : es necesario que los jóvenes hagan ejercicio al menos treinta minutos al día y conviene que duerman ocho horas cada noche. El deporte no es solo para los campeones de fútbol : caminar, correr o bailar ya es suficiente. Además, es importante que las familias vacunen a los niños contra las enfermedades graves.

La higiene también es fundamental : lávate las manos antes de cada comida, mantén limpia tu casa y no tires la basura en la calle, porque el agua sucia transmite el cólera.

Entonces, ¿qué esperas? ¡Cuida tu salud, es tu primer tesoro! Prevenir es más fácil y más barato que curar. Un joven sano es un joven que aprende mejor, que juega mejor y que construye el futuro de su país.
\end{textebox}

\textbf{Traduction du texte :} « Au Cameroun, de nombreuses maladies comme le paludisme ou le choléra peuvent être prévenues par des habitudes simples. Pourtant, chaque année, pendant la saison des pluies, beaucoup de jeunes tombent malades et manquent l'école. Pourquoi ? Souvent, la cause, ce sont les mauvaises habitudes : dormir sans moustiquaire, boire de l'eau non potable ou manger toujours la même chose.

Pourtant, vivre sainement n'est pas un rêve impossible : c'est un effort de tous les jours. Dormir sous une moustiquaire, se laver les mains avec du savon, boire de l'eau potable et manger de façon équilibrée sont des gestes qui sauvent des vies. Au marché de Yaoundé, les fruits et les légumes ne sont pas chers : profites-en ! Une alimentation riche en fruits, en poisson et en céréales renforce le corps contre les infections.

Les médecins le répètent : il est nécessaire que les jeunes fassent de l'exercice au moins trente minutes par jour et il convient qu'ils dorment huit heures chaque nuit. Le sport n'est pas réservé aux champions de football : marcher, courir ou danser suffit déjà. De plus, il est important que les familles fassent vacciner les enfants contre les maladies graves.

L'hygiène aussi est fondamentale : lave-toi les mains avant chaque repas, garde ta maison propre et ne jette pas les ordures dans la rue, car l'eau sale transmet le choléra.

Alors, qu'attends-tu ? Prends soin de ta santé, c'est ton premier trésor ! Prévenir est plus facile et moins cher que guérir. Un jeune en bonne santé est un jeune qui apprend mieux, qui joue mieux et qui construit l'avenir de son pays. »

\textbf{Glossaire (mots difficiles) :}

\begin{center}
\begin{tabularx}{\linewidth}{|X|X|X|X|}
\hline
\rowcolor{popblueL} Español & Français & Español & Français \\
\hline
el reto & le défi & la mosquitera & la moustiquaire \\
\hline
el sueño & le sommeil, le rêve & el esfuerzo & l'effort \\
\hline
prevenir & prévenir & caer enfermo & tomber malade \\
\hline
la temporada de lluvias & la saison des pluies & faltar a la escuela & manquer l'école \\
\hline
el agua potable & l'eau potable & el jabón & le savon \\
\hline
la basura & les ordures & fortalecer & renforcer \\
\hline
salvar vidas & sauver des vies & el tesoro & le trésor \\
\hline
\end{tabularx}
\end{center}

\begin{exemplebox}[Deux phrases clés traduites]
\esp{Dormir bajo una mosquitera previene el paludismo.} = « Dormir sous une moustiquaire prévient le paludisme. » Remarque : en espagnol, l'infinitif \esp{dormir} peut être sujet de la phrase, comme en français.

\esp{... son gestos que salvan vidas.} = « ... sont des gestes qui sauvent des vies. » Le verbe \esp{salvar} (sauver) est un faux ami partiel : il ne signifie jamais « saluer » (qui se dit \esp{saludar}).
\end{exemplebox}

\begin{attention}[Faux amis à ne pas confondre]
\begin{itemize}
\item \esp{salvar} = sauver ; « saluer » se dit \esp{saludar}. \esp{El médico salva vidas.} (Le médecin sauve des vies.) / \esp{Saludo a mi profesora.} (Je salue ma professeure.)
\item \esp{la salud} = la santé ; « le salut » se dit \esp{el saludo}.
\item \esp{el sueño} = le sommeil (et aussi le rêve) ; ne pas confondre avec \esp{el suegro} (le beau-père) !
\end{itemize}
\end{attention}

\begin{savaistu}[L'espagnol, langue de santé publique en Afrique]
La Guinée équatoriale, voisine du Cameroun, est le seul pays d'Afrique subsaharienne où l'espagnol est langue officielle. Les messages de santé publique y sont rédigés en espagnol : affiches contre le paludisme, campagnes de vaccination, conseils d'hygiène... Maîtriser l'espagnol de la santé, c'est donc pouvoir communiquer avec tout un pays frère, de Malabo à Bata !
\end{savaistu}

\courssub{Repérage des idées principales}

Un texte de prévention se construit presque toujours selon le même schéma : un \textbf{problème}, ses \textbf{causes}, puis des \textbf{solutions}. Vérifions-le avec notre texte.

\begin{center}
\begin{tabularx}{\linewidth}{|X|X|X|}
\hline
\rowcolor{popblueL} Élément & Contenu dans le texte & Phrase repère \\
\hline
Problème & Des maladies évitables (paludisme, choléra) rendent les jeunes malades et les empêchent d'aller à l'école. & \esp{muchos jóvenes caen enfermos y faltan a la escuela} \\
\hline
Causes & Les mauvaises habitudes de vie. & \esp{dormir sin mosquitera, beber agua no potable} \\
\hline
Solutions & Des gestes simples : moustiquaire, hygiène, alimentation équilibrée, sport, sommeil, vaccins. & \esp{son gestos que salvan vidas} \\
\hline
\end{tabularx}
\end{center}

\begin{popfigure}
\begin{tikzpicture}[mindmap, grow cyclic, every node/.style={concept, align=center}, concept color=popblue!60,
  root concept/.append style={concept color=popblue, font=\Large\bfseries, text=white},
  level 1/.append style={level distance=3.4cm, sibling angle=90},
  level 1 concept/.append style={font=\small\bfseries}]
\node[root concept]{Vivir sano}
  child[concept color=poporange!70]{ node[align=center]{alimentación\\ equilibrada} }
  child[concept color=popgreen!70]{ node {higiene} }
  child[concept color=poppurple!70]{ node[align=center]{ejercicio\\ y sueño} }
  child[concept color=poppink!70]{ node[align=center]{prevención\\ (mosquitera, vacunas)} };
\end{tikzpicture}
\legende{Carte mentale du texte : les quatre piliers d'une vie saine selon l'auteur.}
\end{popfigure}

\courssub{Compréhension globale}

Réponds d'abord à ces questions générales, à l'oral ou par écrit, en phrases complètes.

\begin{exercice}[Compréhension globale]
\begin{enumerate}
\item \esp{¿De qué trata el texto?} (De quoi traite le texte ?)
\item \esp{¿Qué enfermedades se mencionan?} (Quelles maladies sont mentionnées ?)
\item \esp{¿A quién se dirige el texto?} (À qui le texte s'adresse-t-il ?)
\item \esp{¿Cuál es la idea principal del último párrafo?} (Quelle est l'idée principale du dernier paragraphe ?)
\end{enumerate}
\end{exercice}

\begin{corrige}[Exercice : compréhension globale]
\begin{enumerate}
\item \esp{El texto trata de los hábitos de vida saludables para prevenir las enfermedades en Camerún.} (Le texte traite des habitudes de vie saines pour prévenir les maladies au Cameroun.)
\item \esp{Se mencionan el paludismo, el cólera y las infecciones.} (On mentionne le paludisme, le choléra et les infections.)
\item \esp{Se dirige a los jóvenes cameruneses.} (Il s'adresse aux jeunes Camerounais.)
\item \esp{La idea principal es que prevenir es más fácil y más barato que curar.} (L'idée principale est que prévenir est plus facile et moins cher que guérir.)
\end{enumerate}
\end{corrige}

\courssub{Compréhension fine}

\begin{exoresolu}[Vrai ou faux, avec justification]
\textbf{Énoncé.} Dis si les affirmations suivantes sont \esp{verdadero} (V) ou \esp{falso} (F), et justifie ta réponse par une phrase du texte.

a) \esp{El paludismo no se puede prevenir.}

b) \esp{Las frutas y las verduras son muy caras en Yaoundé.}

c) \esp{Los médicos recomiendan dormir ocho horas.}

\tcblower
\textbf{Solution détaillée.}

a) \textbf{Falso.} Le texte dit : \esp{muchas enfermedades como el paludismo ... se pueden prevenir con hábitos sencillos}. Le paludisme peut donc être prévenu.

b) \textbf{Falso.} Le texte affirme le contraire : \esp{las frutas y las verduras no son caras : ¡aprovecha!}

c) \textbf{Verdadero.} Le texte dit : \esp{conviene que duerman ocho horas cada noche}.
\end{exoresolu}

\begin{exercice}[Retrouve les conseils et les champs lexicaux]
\begin{enumerate}
\item Relève dans le texte quatre conseils donnés aux jeunes et traduis-les en français.
\item Classe les mots suivants dans le bon champ lexical : \esp{paludismo, jabón, vacuna, frutas, cólera, manos, pescado, síntoma, verduras, higiene}.
\begin{itemize}
\item Champ lexical de la maladie : ...
\item Champ lexical de l'hygiène : ...
\item Champ lexical de l'alimentation : ...
\end{itemize}
\end{enumerate}
\end{exercice}

\begin{corrige}[Exercice : conseils et champs lexicaux]
\begin{enumerate}
\item Exemples de conseils : \esp{dormir bajo una mosquitera} (dormir sous une moustiquaire) ; \esp{lavarse las manos con jabón} (se laver les mains avec du savon) ; \esp{beber agua potable} (boire de l'eau potable) ; \esp{hacer ejercicio al menos treinta minutos al día} (faire de l'exercice au moins trente minutes par jour) ; \esp{no tires la basura en la calle} (ne jette pas les ordures dans la rue).
\item Champ lexical de la maladie : \esp{paludismo, cólera, síntoma, vacuna} (la \esp{vacuna} appartient au champ médical de la maladie et de sa prévention). Champ lexical de l'hygiène : \esp{jabón, manos, higiene}. Champ lexical de l'alimentation : \esp{frutas, pescado, verduras}.
\end{enumerate}
\end{corrige}

\courssub{Premier relevé grammatical : impératif et subjonctif}

Avant même d'étudier la grammaire en détail (sections 3 et 4 de cette leçon), observons comment l'auteur \textbf{conseille} et \textbf{recommande}. Deux procédés apparaissent dans le texte.

\begin{propriete}[Deux façons de conseiller observées dans le texte]
\begin{itemize}
\item \textbf{L'impératif} sert à donner un ordre ou un conseil direct : \esp{¡aprovecha!} (profite !), \esp{lávate las manos} (lave-toi les mains), \esp{mantén limpia tu casa} (garde ta maison propre), \esp{no tires la basura} (ne jette pas les ordures), \esp{¡cuida tu salud!} (prends soin de ta santé !).
\item \textbf{Le subjonctif} apparaît après les expressions impersonnelles de nécessité ou de conseil : \esp{es necesario que los jóvenes \textbf{hagan} ejercicio} (il est nécessaire que les jeunes \textbf{fassent} de l'exercice) ; \esp{conviene que \textbf{duerman} ocho horas} (il convient qu'ils \textbf{dorment} huit heures) ; \esp{es importante que las familias \textbf{vacunen} a los niños} (il est important que les familles \textbf{vaccinent} les enfants).
\end{itemize}
\end{propriete}

\begin{attention}[Piège pour les francophones]
En français, on utilise souvent l'infinitif après « il faut » : « il faut dormir huit heures ». En espagnol, quand le sujet de l'action est précisé, on ne peut \textbf{pas} dire \esp{es necesario que los jóvenes hacer ejercicio} : il faut obligatoirement le \textbf{subjonctif} : \esp{es necesario que los jóvenes hagan ejercicio}. L'infinitif n'est possible que si le sujet n'est pas exprimé : \esp{es necesario hacer ejercicio} (il est nécessaire de faire de l'exercice).
\end{attention}

\begin{exercice}[Relevé dans le texte]
\begin{enumerate}
\item Relève tous les verbes à l'impératif du texte et indique leur infinitif.
\item Relève toutes les constructions avec \esp{es necesario que}, \esp{conviene que}, \esp{es importante que} et souligne le verbe au subjonctif.
\end{enumerate}
\end{exercice}

\begin{corrige}[Exercice : relevé grammatical]
\begin{enumerate}
\item Impératifs : \esp{aprovecha} (infinitif \esp{aprovechar}) ; \esp{lávate} (\esp{lavarse}, verbe pronominal) ; \esp{mantén} (\esp{mantener}, irrégulier) ; \esp{no tires} (\esp{tirar}, impératif négatif = subjonctif présent) ; \esp{cuida} (\esp{cuidar}).
\item Subjonctifs : \esp{es necesario que los jóvenes \textbf{hagan} ejercicio} ; \esp{conviene que \textbf{duerman} ocho horas} ; \esp{es importante que las familias \textbf{vacunen} a los niños}.
\end{enumerate}
\end{corrige}

\courssub{Appropriation : et toi ?}

La lecture d'un texte ne s'arrête pas à sa compréhension : il faut le mettre en relation avec ta propre vie. C'est l'étape de l'\textbf{appropriation}.

\begin{experience}[¿Y tú, tienes hábitos de vida saludables?]
Par groupes de deux, posez-vous les questions suivantes et répondez en espagnol, avec des phrases simples :
\begin{enumerate}
\item \esp{¿Duermes bajo una mosquitera?} (Dors-tu sous une moustiquaire ?)
\item \esp{¿Cuántas horas duermes cada noche?} (Combien d'heures dors-tu chaque nuit ?)
\item \esp{¿Haces ejercicio? ¿Qué deporte practicas?} (Fais-tu de l'exercice ? Quel sport pratiques-tu ?)
\item \esp{¿Te lavas las manos antes de comer?} (Te laves-tu les mains avant de manger ?)
\item \esp{¿Qué comes normalmente en el recreo?} (Que manges-tu normalement à la récréation ?)
\end{enumerate}
Modèle de réponse : \esp{Sí, duermo bajo una mosquitera, pero solo duermo seis horas. Juego al fútbol los sábados.} (Oui, je dors sous une moustiquaire, mais je ne dors que six heures. Je joue au football le samedi.)
\end{experience}

\begin{aretenir}[L'essentiel de cette section]
\begin{itemize}
\item Un texte de prévention suit le schéma \textbf{problème -- causes -- solutions} ; le repérer aide à comprendre et à résumer.
\item Les \textbf{mots transparents} (\esp{prevención, higiene, vacuna, alimentación}) permettent de deviner le sens général sans dictionnaire.
\item Pour conseiller, l'espagnol utilise l'\textbf{impératif} (\esp{¡cuida tu salud!}) ou le \textbf{subjonctif} après \esp{es necesario que}, \esp{conviene que}, \esp{es importante que}.
\item Vocabulaire clé : \esp{el reto} (défi), \esp{la mosquitera} (moustiquaire), \esp{el sueño} (sommeil), \esp{el esfuerzo} (effort), \esp{prevenir} (prévenir), \esp{salvar vidas} (sauver des vies).
\end{itemize}
\end{aretenir}

Tu sais maintenant lire et analyser un texte de santé. La prochaine étape : enrichir ton vocabulaire avec le lexique thématique de la santé, avant d'apprendre à formuler toi-même des conseils à l'impératif et au subjonctif.

\coursec{Lexique thématique de la santé}

Dans la section précédente, nous avons lu et commenté le texte \esp{Vivir sano}, qui décrivait les habitudes de vie favorables à la santé : manger équilibré, dormir suffisamment, se laver les mains, faire de l'exercice. Ce texte contenait de nombreux mots nouveaux : \esp{salud}, \esp{enfermedad}, \esp{vacuna}, \esp{higiene}\ldots{} Le moment est venu de les organiser, de les mémoriser et de savoir les employer correctement. Un bon lexique n'est pas une liste apprise par cœur : c'est un réseau de mots reliés entre eux par le sens. Nous allons donc classer le vocabulaire de la santé en cinq champs lexicaux, puis apprendre à construire avec eux de vrais messages de prévention.

\begin{definition}[Champ lexical]
Un \cle{champ lexical} est l'ensemble des mots qui se rapportent à une même notion ou à un même domaine de la réalité. Par exemple, \esp{fiebre}, \esp{tos} et \esp{dolor de cabeza} appartiennent au champ lexical des symptômes. Classer les mots par champs facilite la mémorisation et permet de rédiger un texte cohérent.
\end{definition}

\courssub{Champ 1 : la santé et le corps}

Observons d'abord ces phrases tirées de situations de la vie quotidienne :

\begin{exemplebox}[Observer avant d'apprendre]
\esp{La salud es lo más importante.} « La santé est le plus important. »

\esp{Mi abuelo tiene ochenta años, pero está muy sano.} « Mon grand-père a quatre-vingts ans, mais il est en très bonne santé. »

\esp{El deporte es bueno para el cuerpo.} « Le sport est bon pour le corps. »

\esp{Para llevar una vida sana, hay que dormir bien.} « Pour mener une vie saine, il faut bien dormir. »

\esp{María hace deporte todos los días: está en forma.} « María fait du sport tous les jours : elle est en forme. »
\end{exemplebox}

\begin{center}
\begin{tabularx}{\linewidth}{|X|X|X|}
\hline
\rowcolor{popblueL} Español & Français & Remarque \\
\hline
la salud & la santé & mot féminin ; sans article, on dit \esp{¡Salud!} « Santé ! / À tes souhaits ! » \\
\hline
sano, sana & sain, en bonne santé & adjectif variable : \esp{un niño sano}, \esp{una niña sana} \\
\hline
el cuerpo & le corps & masculin, comme la plupart des mots en -o : \esp{el cuerpo humano} \\
\hline
la vida sana & la vie saine & \esp{llevar una vida sana} = mener une vie saine \\
\hline
estar en forma & être en forme & toujours avec \esp{estar} (état), jamais avec \esp{ser} \\
\hline
la cabeza & la tête & \esp{me duele la cabeza} = j'ai mal à la tête \\
\hline
las manos & les mains & féminin malgré le -o : \esp{la mano}, \esp{las manos} \\
\hline
\end{tabularx}
\end{center}

\begin{attention}[Sano : ser ou estar ?]
En français, « être en bonne santé » ne pose aucun problème. En espagnol, on distingue : \esp{ser sano} (être d'une constitution saine, être bon pour la santé : \esp{La fruta es sana} = le fruit est sain, bon pour la santé) et \esp{estar sano} (être en bonne santé à un moment donné : \esp{Estoy sano, gracias} = je suis en bonne santé, merci). De même, \esp{estar en forma} utilise toujours \esp{estar}, car c'est un état qui peut changer.
\end{attention}

\courssub{Champ 2 : les maladies}

Voici le champ le plus important pour nos messages de prévention. Observons :

\esp{El paludismo es una enfermedad frecuente en Camerún.} « Le paludisme est une maladie fréquente au Cameroun. »

\esp{El cólera se transmite por el agua contaminada.} « Le choléra se transmet par l'eau contaminée. »

\esp{La gripe se contagia fácilmente en las escuelas.} « La grippe se contagie facilement dans les écoles. »

\begin{center}
\begin{tabularx}{\linewidth}{|X|X|X|}
\hline
\rowcolor{popblueL} Español & Français & Remarque \\
\hline
la enfermedad & la maladie & mot féminin en -dad ; jamais *\esp{maladia} \\
\hline
el paludismo & le paludisme & aussi \esp{la malaria} ; transmis par le moustique \\
\hline
el cólera & le choléra & masculin : \esp{el cólera} ; accent sur la première syllabe (mot \emph{esdrújula}) \\
\hline
el VIH/sida & le VIH/sida & \esp{VIH} se prononce « uve-i-hache » ; \esp{sida} est masculin : \esp{el sida} \\
\hline
la gripe & la grippe & féminin ; \esp{tener gripe} = avoir la grippe \\
\hline
enfermo, enferma & malade & adjectif et nom : \esp{un enfermo} = un malade \\
\hline
contagiar & contaminer, transmettre & \esp{contagiarse} = être contaminé ; verbe régulier en -ar \\
\hline
el mosquito & le moustique & \esp{el mosquito transmite el paludismo} \\
\hline
\end{tabularx}
\end{center}

\begin{attention}[Les deux erreurs classiques des francophones]
1. « Je suis malade » se dit \esp{Estoy enfermo} avec \esp{estar} (état passager), jamais *\esp{Soy enfermo}, qui suggérerait une caractéristique permanente.

2. « Maladie » ne se traduit jamais par *\esp{maladia} : ce mot n'existe pas en espagnol. Le bon mot est \esp{la enfermedad}. Attention aussi au genre : on dit \esp{una enfermedad grave} (une maladie grave).
\end{attention}

\courssub{Champ 3 : les symptômes et la douleur}

Pour décrire ce que l'on ressent, l'espagnol utilise très souvent le verbe \esp{tener} (avoir), comme le français, mais pas toujours de la même façon.

\esp{Tiene fiebre y dolor de cabeza.} « Il a de la fièvre et mal à la tête. »

\esp{El niño tiene tos desde ayer.} « L'enfant tousse depuis hier. »

\esp{Me siento mal, tengo que ir al médico.} « Je me sens mal, je dois aller chez le médecin. »

\begin{center}
\begin{tabularx}{\linewidth}{|X|X|X|}
\hline
\rowcolor{popblueL} Español & Français & Remarque \\
\hline
el dolor & la douleur & \esp{tener dolor de cabeza} = avoir mal à la tête \\
\hline
el síntoma & le symptôme & masculin malgré le -a : \esp{el síntoma} \\
\hline
la fiebre & la fièvre & \esp{tener fiebre} = avoir de la fièvre \\
\hline
la tos & la toux & \esp{tener tos} = tousser, avoir la toux \\
\hline
el dolor de cabeza & le mal de tête & littéralement « douleur de tête » \\
\hline
sentirse mal & se sentir mal & verbe pronominal : \esp{me siento, te sientes, se siente\ldots} \\
\hline
cansado/a & fatigué & \esp{estar cansado} = être fatigué \\
\hline
\end{tabularx}
\end{center}

\begin{propriete}[Dire « j'ai mal à\ldots » : la construction avec doler]
Pour exprimer une douleur localisée, l'espagnol emploie le verbe \esp{doler} (faire mal) construit comme \esp{gustar} : \esp{Me duele la cabeza} « J'ai mal à la tête » (littéralement « la tête me fait mal ») ; \esp{Me duelen las manos} « J'ai mal aux mains ». Le verbe s'accorde avec la partie du corps : \esp{duele} au singulier, \esp{duelen} au pluriel. On peut aussi dire simplement : \esp{Tengo dolor de cabeza}.
\end{propriete}

\begin{popfigure}
\begin{center}
\begin{tikzpicture}[scale=1]
% silhouette
\draw[thick, popblue, fill=popblueL] (0,2.1) circle (0.55);
\draw[thick, popblue, fill=popblueL, rounded corners=8pt] (-0.55,1.4) rectangle (0.55,-0.6);
\draw[thick, popblue, fill=popblueL, rounded corners=4pt] (-1.15,1.3) rectangle (-0.6,-0.1);
\draw[thick, popblue, fill=popblueL, rounded corners=4pt] (0.6,1.3) rectangle (1.15,-0.1);
\draw[thick, popblue, fill=popblueL, rounded corners=4pt] (-0.5,-0.6) rectangle (-0.1,-2.0);
\draw[thick, popblue, fill=popblueL, rounded corners=4pt] (0.1,-0.6) rectangle (0.5,-2.0);
% étiquettes parties du corps
\node[align=center, fill=popgreenL, rounded corners=3pt, draw=popgreen] (cab) at (-3.2,2.4) {\textbf{la cabeza}\\ la tête};
\draw[-{Stealth}, popgreen, thick] (cab.east) -- (-0.55,2.25);
\node[align=center, fill=popgreenL, rounded corners=3pt, draw=popgreen] (cue) at (3.4,0.5) {\textbf{el cuerpo}\\ le corps};
\draw[-{Stealth}, popgreen, thick] (cue.west) -- (0.55,0.5);
\node[align=center, fill=popgreenL, rounded corners=3pt, draw=popgreen] (man) at (-3.2,-0.2) {\textbf{las manos}\\ les mains};
\draw[-{Stealth}, popgreen, thick] (man.east) -- (-1.15,-0.1);
% bulles symptômes
\node[align=center, fill=poporangeL, rounded corners=3pt, draw=poporange] (dol) at (3.4,2.4) {\textbf{el dolor de cabeza}\\ le mal de tête};
\draw[-{Stealth}, poporange, thick, dashed] (dol.west) -- (0.55,2.3);
\node[align=center, fill=poporangeL, rounded corners=3pt, draw=poporange] (fie) at (3.4,-1.2) {\textbf{la fiebre}\\ la fièvre};
\draw[-{Stealth}, poporange, thick, dashed] (fie.west) -- (0.5,-0.9);
\node[align=center, fill=poporangeL, rounded corners=3pt, draw=poporange] (tos) at (-3.2,-1.6) {\textbf{la tos}\\ la toux};
\draw[-{Stealth}, poporange, thick, dashed] (tos.east) -- (-0.4,-1.1);
\end{tikzpicture}
\end{center}
\legende{Le corps et les symptômes : en vert, les parties du corps ; en orange, les symptômes.}
\end{popfigure}

\courssub{Champ 4 : la prévention et les soins}

\esp{La vacuna protege contra la enfermedad.} « Le vaccin protège contre la maladie. »

\esp{Los niños se vacunan en el hospital.} « Les enfants se font vacciner à l'hôpital. »

\esp{La higiene es fundamental para prevenir el cólera.} « L'hygiène est fondamentale pour prévenir le choléra. »

\begin{center}
\begin{tabularx}{\linewidth}{|X|X|X|}
\hline
\rowcolor{popblueL} Español & Français & Remarque \\
\hline
la vacuna & le vaccin & \esp{poner una vacuna} = faire un vaccin \\
\hline
vacunar(se) & vacciner / se faire vacciner & \esp{Me vacuno contra la gripe} = je me fais vacciner contre la grippe \\
\hline
la higiene & l'hygiène & prononcer « i-hié-né » ; le h est muet, le g devant e se prononce comme le j espagnol \\
\hline
la prevención & la prévention & un seul \emph{v}, accent sur la \emph{o} finale \\
\hline
prevenir & prévenir, empêcher & irrégulier comme \esp{venir} : \esp{prevengo, previenes\ldots} \\
\hline
el médico / la médica & le médecin / la femme médecin & aussi \esp{el doctor / la doctora} \\
\hline
el hospital & l'hôpital & le h est muet : « ospi-tal » \\
\hline
el tratamiento & le traitement & \esp{seguir un tratamiento} = suivre un traitement \\
\hline
la mosquitera & la moustiquaire & \esp{dormir bajo mosquitera} = dormir sous moustiquaire \\
\hline
\end{tabularx}
\end{center}

\begin{attention}[Orthographe : prevención et prevenir]
Le français écrit « prévention » avec \emph{pré-} ; l'espagnol écrit \esp{prevención} sans accent sur le e initial, mais avec un accent obligatoire sur la dernière syllabe. Le verbe \esp{prevenir} suit la conjugaison irrégulière de \esp{venir} : \esp{yo prevengo}, \esp{tú previenes}, \esp{él previene}, \esp{nosotros prevenimos}, \esp{vosotros prevenís}, \esp{ellos previenen}. Ne confondez pas non plus \esp{prevenir} (empêcher) et \esp{avisar} (avertir quelqu'un).
\end{attention}

\begin{definition}[Alimentación equilibrada]
\esp{La alimentación equilibrada} es una alimentación que contiene todos los nutrientes necesarios para el cuerpo: frutas, verduras, proteínas, cereales y agua. « L'alimentation équilibrée est une alimentation qui contient tous les nutriments nécessaires au corps : fruits, légumes, protéines, céréales et eau. » Exemple : \esp{Una alimentación equilibrada incluye frutas y verduras todos los días.} « Une alimentation équilibrée comprend des fruits et des légumes tous les jours. »
\end{definition}

\courssub{Champ 5 : les habitudes de vie}

C'est le champ lexical de l'action quotidienne, celui qui servira directement à rédiger vos messages de prévention.

\begin{center}
\begin{tabularx}{\linewidth}{|X|X|X|}
\hline
\rowcolor{popblueL} Español & Français & Remarque \\
\hline
la alimentación equilibrada & l'alimentation équilibrée & \esp{comer de forma equilibrada} = manger équilibré \\
\hline
el ejercicio & l'exercice (physique) & \esp{hacer ejercicio} = faire de l'exercice \\
\hline
dormir & dormir & irrégulier : \esp{duermo, duermes, duerme\ldots} \\
\hline
beber agua potable & boire de l'eau potable & \esp{agua} est féminin malgré \esp{el agua} \\
\hline
lavarse las manos & se laver les mains & verbe pronominal + article défini : \esp{me lavo las manos} \\
\hline
fumar / no fumar & fumer / ne pas fumer & \esp{Fumar es malo para la salud} = fumer est mauvais pour la santé \\
\hline
descansar & se reposer & \esp{Hay que descansar ocho horas} = il faut se reposer huit heures \\
\hline
\end{tabularx}
\end{center}

\begin{propriete}[Le corps et l'article défini]
En français on dit « je me lave \emph{mes} mains » ; en espagnol, avec un verbe pronominal, la partie du corps prend l'\cle{article défini} et non l'adjectif possessif : \esp{Me lavo las manos} (jamais *\esp{me lavo mis manos}). De même : \esp{Me cepillo los dientes} « je me brosse les dents ».
\end{propriete}

\begin{exemplebox}[Expressions utiles pour les messages de prévention]
\esp{Protege tu salud.} « Protège ta santé. »

\esp{Lava tus manos con agua y jabón.} « Lave tes mains avec de l'eau et du savon. »

\esp{La prevención es la mejor medicina.} « La prévention est le meilleur médicament. »

\esp{Duerme bajo una mosquitera.} « Dors sous une moustiquaire. »

\esp{Bebe agua potable.} « Bois de l'eau potable. »

\esp{No fumes: el tabaco mata.} « Ne fume pas : le tabac tue. »

Ces phrases sont construites à l'impératif : nous étudierons cette grammaire dans la section suivante.
\end{exemplebox}

\courssub{Mots transparents et faux amis}

De nombreux mots de la santé sont \cle{transparents} : ils se ressemblent en français et en espagnol. C'est une chance, mais attention aux pièges.

\begin{center}
\begin{tabularx}{\linewidth}{|X|X|X|}
\hline
\rowcolor{popblueL} Français & Español & Remarque \\
\hline
la prévention & la prevención & orthographe : un seul \emph{v}, accent final \\
\hline
le symptôme & el síntoma & masculin malgré le -a \\
\hline
l'hygiène & la higiene & h muet, g devant e se prononce « hié » \\
\hline
le traitement & el tratamiento & suffixe -miento = -ment \\
\hline
être enrhumé & estar constipado & FAUX AMI : \esp{constipado} = enrhumé ; « constipé » (digestion) = \esp{estreñido} \\
\hline
la maladie & la enfermedad & jamais *\esp{maladia} \\
\hline
\end{tabularx}
\end{center}

\begin{attention}[Le faux ami célèbre : constipado]
\esp{Estoy constipado} signifie « je suis enrhumé, j'ai un rhume », et non « je suis constipé ». Pour dire « constipé » au sens digestif, l'espagnol utilise \esp{estreñido}. Retenez aussi \esp{un resfriado} = un rhume. Cette confusion est un classique des sujets de version à l'examen.
\end{attention}

\begin{savaistu}[Le vocabulaire des vraies affiches]
Au Cameroun, les campagnes de vaccination et la distribution de moustiquaires imprégnées ont fait reculer le paludisme, première cause de consultation dans les hôpitaux. Les affiches du ministère de la Santé publique utilisent exactement le vocabulaire de cette leçon : \esp{vacuna}, \esp{mosquitera}, \esp{agua potable}, \esp{lavarse las manos}. En Guinée équatoriale, pays hispanophone voisin du Cameroun, ces mêmes mots circulent dans les messages de prévention en espagnol : apprendre ce lexique, c'est apprendre la langue de la santé publique en Afrique centrale.
\end{savaistu}

\courssub{Entraînement : utiliser le lexique}

\begin{exoresolu}[Compléter un message de prévention]
Énoncé. Complétez ce message de prévention avec les mots suivants : \esp{vacuna}, \esp{manos}, \esp{paludismo}, \esp{agua potable}, \esp{mosquitera}.

\esp{Para prevenir las enfermedades, lava tus \ldots{} con jabón, bebe siempre \ldots{}, duerme bajo una \ldots{} para evitar el \ldots{} y pide la \ldots{} en el centro de salud.}

\tcblower

Solution détaillée. On raisonne par le sens de chaque phrase :

1. \esp{lava tus manos con jabón} : on lave les \emph{mains} avec du savon (hygiène).

2. \esp{bebe siempre agua potable} : on boit de l'\emph{eau potable} (prévention du choléra).

3. \esp{duerme bajo una mosquitera} : on dort sous une \emph{moustiquaire}.

4. \esp{para evitar el paludismo} : la moustiquaire protège du \emph{paludisme} (masculin : \esp{el paludisme}).

5. \esp{pide la vacuna en el centro de salud} : on demande le \emph{vaccin} au centre de santé.

Phrase complète traduite : « Pour prévenir les maladies, lave tes mains avec du savon, bois toujours de l'eau potable, dors sous une moustiquaire pour éviter le paludisme et demande le vaccin au centre de santé. »
\end{exoresolu}

\begin{exercice}[À vous de jouer]
1. Traduisez en espagnol : a) « Mon frère est malade : il a de la fièvre. » b) « Le vaccin protège contre la maladie. » c) « Je me lave les mains avant de manger. » d) « Pour être en forme, il faut dormir huit heures. »

2. Trouvez l'intrus dans chaque série et justifiez : a) \esp{fiebre / tos / vacuna / dolor de cabeza} ; b) \esp{médico / hospital / tratamiento / ejercicio}.

3. Rédigez trois slogans de prévention contre le choléra en utilisant : \esp{agua potable}, \esp{lavarse las manos}, \esp{higiene}.
\end{exercice}

\begin{corrige}[Exercice : utiliser le lexique]
1. a) \esp{Mi hermano está enfermo: tiene fiebre.} (état : \esp{estar} ; \esp{tener fiebre}) b) \esp{La vacuna protege contra la enfermedad.} c) \esp{Me lavo las manos antes de comer.} (article défini, pas de possessif) d) \esp{Para estar en forma, hay que dormir ocho horas.}

2. a) L'intrus est \esp{vacuna} : c'est un moyen de prévention, les trois autres sont des symptômes. b) L'intrus est \esp{ejercicio} : c'est une habitude de vie, les trois autres appartiennent au champ des soins.

3. Exemples de réponses : \esp{Bebe siempre agua potable.} ; \esp{Lávate las manos antes de comer.} ; \esp{La higiene protege a tu familia.}
\end{corrige}

\begin{aretenir}[L'essentiel du lexique de la santé]
\begin{itemize}
\item « Maladie » = \esp{la enfermedad} ; « être malade » = \esp{estar enfermo} (avec \esp{estar}).
\item Symptômes : \esp{tener fiebre}, \esp{tener tos}, \esp{tener dolor de cabeza} ou \esp{me duele la cabeza}.
\item Prévention : \esp{la vacuna}, \esp{vacunarse}, \esp{la higiene}, \esp{prevenir} (conjugué comme \esp{venir}).
\item Habitudes de vie : \esp{hacer ejercicio}, \esp{dormir}, \esp{beber agua potable}, \esp{lavarse las manos} (article défini avec les parties du corps).
\item Faux ami : \esp{constipado} = enrhumé ; « constipé » = \esp{estreñido}.
\end{itemize}
\end{aretenir}

Ce lexique est désormais votre boîte à outils. Dans les sections suivantes, nous apprendrons à le mettre en phrases grâce à l'impératif et au subjonctif, afin de rédiger de véritables messages de prévention comme ceux des campagnes de santé publique.

\coursec{Grammaire 1 : l'impératif pour conseiller}

Dans la section précédente, nous avons constitué notre boîte à outils lexicale : \esp{salud}, \esp{enfermedad}, \esp{paludismo}, \esp{vacuna}, \esp{higiene}, \esp{alimentación equilibrada}, \esp{ejercicio}, \esp{dolor}, \esp{síntoma}. Mais un vocabulaire ne suffit pas : pour \emph{conseiller}, \emph{recommander} ou \emph{interdire}, il faut un mode verbal particulier. Or, si l'on relit le texte \esp{Vivir sano}, on remarque que presque toutes ses phrases de conseil sont construites avec la même forme verbale : l'\cle{impératif}. C'est lui que nous allons maintenant observer, construire et manipuler, car c'est l'outil grammatical numéro un de tout message de prévention.

\courssub{Observation : les conseils du texte \esp{Vivir sano}}

Reprenons quatre phrases tirées du texte support et des affiches de la campagne de santé qu'il décrit :

\begin{exemplebox}[Phrases observées]
\esp{Cuida tu salud.} « Prends soin de ta santé. »

\esp{Lava tus manos con agua y jabón.} « Lave tes mains avec de l'eau et du savon. »

\esp{No fumes.} « Ne fume pas. »

\esp{Hagan ejercicio cada día.} « Faites du sport chaque jour. »
\end{exemplebox}

Observons attentivement :

\begin{itemize}
\item Dans \esp{Cuida} et \esp{Lava}, le verbe ressemble à la 3\textsuperscript{e} personne du singulier du présent de l'indicatif (\esp{él cuida}, \esp{él lava}), mais il n'y a \textbf{pas de sujet exprimé} : le verbe s'adresse directement à quelqu'un.
\item Dans \esp{No fumes}, la forme change : ce n'est plus \esp{fumas} mais \esp{fumes}, avec un \esp{-e-} caractéristique du subjonctif.
\item Dans \esp{Hagan}, on s'adresse à un groupe (\esp{ustedes}) et la forme est encore différente.
\end{itemize}

Conclusion de l'observation : l'espagnol possède un mode spécial pour donner des ordres et des conseils, l'impératif, dont les formes \emph{affirmatives} et \emph{négatives} ne se construisent pas de la même manière. C'est la grande différence avec le français, où « lave » et « ne lave pas » utilisent la même forme verbale.

\courssub{Formation de l'impératif affirmatif}

\begin{propriete}[Règle de formation de l'impératif affirmatif]
\begin{itemize}
\item \textbf{tú} : 3\textsuperscript{e} personne du singulier du présent de l'indicatif : \esp{lava}, \esp{come}, \esp{duerme}, \esp{cuida}.
\item \textbf{usted} : présent de subjonctif, 3\textsuperscript{e} personne du singulier : \esp{lave}, \esp{coma}, \esp{duerma}.
\item \textbf{nosotros} : présent de subjonctif, 1\textsuperscript{re} personne du pluriel : \esp{lavemos}, \esp{comamos} (« lavons-nous, mangeons »).
\item \textbf{vosotros} : infinitif dont on remplace le \esp{-r} final par un \esp{-d} : \esp{lavar} $\to$ \esp{lavad}, \esp{comer} $\to$ \esp{comed}, \esp{dormir} $\to$ \esp{dormid}.
\item \textbf{ustedes} : présent de subjonctif, 3\textsuperscript{e} personne du pluriel : \esp{laven}, \esp{coman}, \esp{duerman}.
\end{itemize}
\end{propriete}

\begin{center}
\begin{tabularx}{\linewidth}{|l|X|X|X|}
\hline
\rowcolor{popblueL} Personne & \esp{lavar} (laver) & \esp{comer} (manger) & \esp{dormir} (dormir) \\
\hline
\esp{tú} & \esp{lava} & \esp{come} & \esp{duerme} \\
\hline
\esp{usted} & \esp{lave} & \esp{coma} & \esp{duerma} \\
\hline
\esp{nosotros} & \esp{lavemos} & \esp{comamos} & \esp{durmamos} \\
\hline
\esp{vosotros} & \esp{lavad} & \esp{comed} & \esp{dormid} \\
\hline
\esp{ustedes} & \esp{laven} & \esp{coman} & \esp{duerman} \\
\hline
\end{tabularx}
\end{center}

\begin{aretenir}[Deux astuces de mémorisation]
\begin{itemize}
\item Pour \esp{tú} affirmatif : prends la forme de \esp{él/ella} au présent. \esp{Ella bebe agua} $\to$ \esp{¡Bebe agua!} « Bois de l'eau ! »
\item Pour \esp{vosotros} : infinitif moins \esp{r}, plus \esp{d}. \esp{correr} $\to$ \esp{corred} « courez ».
\end{itemize}
\end{aretenir}

Remarque importante pour le contexte camerounais : en Amérique latine et en Guinée équatoriale, on n'utilise pratiquement pas \esp{vosotros} ; on s'adresse à un groupe avec \esp{ustedes} : \esp{Hagan ejercicio}, \esp{Duerman bajo una mosquitera}. C'est la forme que l'on trouve sur les affiches de santé publique.

\courssub{L'impératif négatif : toujours le subjonctif}

\begin{propriete}[Règle de l'impératif négatif]
Pour interdire ou déconseiller, l'espagnol n'utilise jamais les formes affirmatives précédées de \esp{no}. Il emploie \textbf{no + présent de subjonctif} à toutes les personnes :

\esp{no laves}, \esp{no lave}, \esp{no lavemos}, \esp{no lavéis}, \esp{no laven}

\esp{no comas}, \esp{no coma}, \esp{no comamos}, \esp{no comáis}, \esp{no coman}
\end{propriete}

\begin{exemplebox}[Interdictions traduites]
\esp{No fumes.} « Ne fume pas. »

\esp{No bebas agua sucia.} « Ne bois pas d'eau sale. »

\esp{No coman alimentos grasos en exceso.} « Ne mangez pas trop d'aliments gras. »

\esp{No olvidéis la vacuna.} « N'oubliez pas le vaccin. »
\end{exemplebox}

\begin{attention}[Piège majeur pour les francophones]
En français, « tu manges » et « ne mange pas » utilisent la même voyelle. En espagnol, non : l'affirmatif de \esp{tú} est \esp{come} (avec \esp{-e}) et le négatif est \esp{no comas} (avec \esp{-a}). On ne dit jamais *\esp{no comes}. De même : \esp{lava} mais \esp{no laves} ; \esp{duerme} mais \esp{no duermas}. Astuce : la voyelle « s'inverse » entre l'affirmatif et le négatif pour les verbes en \esp{-ar} et \esp{-er/-ir}.
\end{attention}

\courssub{Les verbes irréguliers fréquents dans le domaine de la santé}

Huit verbes très fréquents ont un impératif \esp{tú} irrégulier et court. Ils doivent être appris par cœur, car ils reviennent dans presque tous les messages de prévention :

\begin{center}
\begin{tabularx}{\linewidth}{|l|l|X|}
\hline
\rowcolor{popgreenL} Infinitif & Impératif \esp{tú} & Exemple traduit \\
\hline
\esp{hacer} & \esp{haz} & \esp{Haz ejercicio.} « Fais du sport. » \\
\hline
\esp{ir} & \esp{ve} & \esp{Ve al médico.} « Va chez le médecin. » \\
\hline
\esp{poner} & \esp{pon} & \esp{Pon la mesa limpia.} « Mets une table propre. » \\
\hline
\esp{ser} & \esp{sé} & \esp{Sé prudente.} « Sois prudent. » \\
\hline
\esp{tener} & \esp{ten} & \esp{Ten cuidado.} « Fais attention. » \\
\hline
\esp{salir} & \esp{sal} & \esp{Sal a caminar.} « Sors marcher. » \\
\hline
\esp{decir} & \esp{di} & \esp{Di la verdad al doctor.} « Dis la vérité au docteur. » \\
\hline
\esp{venir} & \esp{ven} & \esp{Ven a vacunarte.} « Viens te faire vacciner. » \\
\hline
\end{tabularx}
\end{center}

\begin{attention}[Accent obligatoire]
\esp{sé} (impératif de \esp{ser}) prend un accent pour se distinguer de \esp{se} (pronom). De même \esp{ve} (impératif de \esp{ir}) ne prend pas d'accent, mais \esp{vé} n'existe pas : ne confonds pas avec \esp{ve} de \esp{ver} (« il voit »), qui s'écrit aussi sans accent.
\end{attention}

Les verbes à changement de voyelle gardent ce changement à l'impératif \esp{tú} : \esp{dormir} $\to$ \esp{duerme} « dors », \esp{pensar} $\to$ \esp{piensa} « pense », \esp{volver} $\to$ \esp{vuelve} « reviens ». Exemple : \esp{Duerme ocho horas cada noche.} « Dors huit heures chaque nuit. »

\courssub{Le pronom réfléchi : la grande différence avec le français}

\begin{propriete}[Place du pronom avec l'impératif]
\begin{itemize}
\item À l'\textbf{affirmatif}, le pronom se place \textbf{après} le verbe et s'y \textbf{accolle} en un seul mot : \esp{lávate}, \esp{lávese}, \esp{lavémonos}, \esp{lavaos}, \esp{lávense}.
\item Au \textbf{négatif}, le pronom se place \textbf{avant} le verbe, détaché : \esp{no te laves}, \esp{no se lave}, \esp{no nos lavemos}, \esp{no os lavéis}, \esp{no se laven}.
\end{itemize}
\end{propriete}

\begin{center}
\begin{tabularx}{\linewidth}{|X|X|X|}
\hline
\rowcolor{popblueL} Français & Español & Remarque \\
\hline
Lave-toi les mains. & \esp{Lávate las manos.} & pronom accolé, accent ajouté \\
\hline
Ne te lave pas les mains sales. & \esp{No te laves con agua sucia.} & pronom avant le verbe \\
\hline
Lavez-vous les mains. & \esp{Lávense las manos.} & forme \esp{ustedes} \\
\hline
Vaccinons-nous. & \esp{Vacunémonos.} & \esp{-mos} + \esp{nos} : le \esp{s} final tombe \\
\hline
\end{tabularx}
\end{center}

\begin{attention}[Deux erreurs typiques des francophones]
\begin{enumerate}
\item \textbf{Oublier l'accent} quand on accole le pronom : \esp{lava} + \esp{te} donne \esp{lávate}, avec accent sur la voyelle tonique qui reste la même. On n'écrit jamais *\esp{lavate} ni *\esp{lava te} en deux mots.
\item \textbf{Garder le possessif français} avec les parties du corps : le français dit « lave \emph{tes} mains », l'espagnol dit \esp{lávate \textbf{las} manos}, avec l'article défini, car le pronom réfléchi indique déjà à qui appartiennent les mains. On ne dit pas *\esp{lava tus manos} quand on se lave soi-même (la phrase du texte \esp{Lava tus manos} est acceptée familièrement, mais la norme est \esp{lávate las manos}).
\end{enumerate}
\end{attention}

\begin{popfigure}
\begin{center}
\begin{tikzpicture}[node distance=0.55cm and 1.1cm]
\node[draw=popblue, fill=popblueL, rounded corners, align=center, font=\bfseries] (base) {\esp{LAVARSE}\\ se laver};
\node[draw=popgreen, fill=popgreenL, rounded corners, align=center, above left=0.55cm and 0.2cm of base] (tu) {\esp{tú}\\ \esp{lávate}};
\node[draw=popgreen, fill=popgreenL, rounded corners, align=center, above right=0.55cm and 0.2cm of base] (ud) {\esp{usted}\\ \esp{lávese}};
\node[draw=poporange, fill=poporangeL, rounded corners, align=center, below left=0.55cm and 0.2cm of base] (nos) {\esp{nosotros}\\ \esp{lavémonos}};
\node[draw=poporange, fill=poporangeL, rounded corners, align=center, below=0.55cm of base] (vos) {\esp{vosotros}\\ \esp{lavaos}};
\node[draw=poporange, fill=poporangeL, rounded corners, align=center, below right=0.55cm and 0.2cm of base] (uds) {\esp{ustedes}\\ \esp{lávense}};
\draw[-{Stealth}, thick, popblue] (base) -- (tu);
\draw[-{Stealth}, thick, popblue] (base) -- (ud);
\draw[-{Stealth}, thick, popblue] (base) -- (nos);
\draw[-{Stealth}, thick, popblue] (base) -- (vos);
\draw[-{Stealth}, thick, popblue] (base) -- (uds);
\end{tikzpicture}
\end{center}
\legende{Les cinq formes de l'impératif affirmatif de \esp{lavarse} : le pronom s'accolle et l'accent se maintient sur la voyelle tonique.}
\end{popfigure}

\courssub{Emplois de l'impératif : conseiller, recommander, interdire}

L'impératif sert à :

\begin{itemize}
\item \textbf{Conseiller} : \esp{Come frutas y verduras.} « Mange des fruits et des légumes. »
\item \textbf{Recommander} (santé publique) : \esp{Duerman bajo una mosquitera.} « Dormez sous une moustiquaire. »
\item \textbf{Interdire} : \esp{No tiren basura en la calle.} « Ne jetez pas d'ordures dans la rue. »
\item \textbf{Rédiger des slogans} : \esp{¡Vacúnate contra el paludismo!} « Vaccine-toi contre le paludisme ! »
\end{itemize}

\begin{popfigure}
\begin{center}
\begin{tikzpicture}[node distance=0.7cm and 1.6cm]
\node[draw=popdark, fill=popcream, rounded corners, align=center, font=\bfseries] (start) {Je veux donner\\ un conseil};
\node[draw=popgreen, fill=popgreenL, rounded corners, align=center, below left=0.7cm and 1.6cm of start] (aff) {Conseil positif\\ \textbf{impératif affirmatif}\\ \esp{Bebe agua limpia.}};
\node[draw=poppink, fill=poppinkL, rounded corners, align=center, below right=0.7cm and 1.6cm of start] (neg) {Interdiction\\ \textbf{no + subjonctif}\\ \esp{No bebas agua sucia.}};
\draw[-{Stealth}, very thick, popgreen] (start) -- node[above, sloped, font=\small]{oui / fais} (aff);
\draw[-{Stealth}, very thick, poppink] (start) -- node[above, sloped, font=\small]{non / ne fais pas} (neg);
\end{tikzpicture}
\end{center}
\legende{Organigramme de choix : conseil affirmatif (impératif) ou interdiction (\esp{no} + subjonctif).}
\end{popfigure}

\courssub{Méthode : rédiger un slogan de prévention}

\begin{methode}[Construire un slogan de santé en trois étapes]
\begin{enumerate}
\item \textbf{Choisir le verbe d'action} et le mettre à l'impératif, à la personne du public visé : \esp{come}, \esp{lávate}, \esp{duerman}, \esp{vacúnate}.
\item \textbf{Ajouter un complément court et concret} : \esp{frutas y verduras}, \esp{las manos con jabón}, \esp{bajo una mosquitera}.
\item \textbf{Ajouter le bénéfice} pour motiver : \esp{tu cuerpo te lo agradece} « ton corps t'en remercie », \esp{la salud es vida} « la santé, c'est la vie ».
\end{enumerate}
Modèle complet : \esp{Come frutas y verduras : tu cuerpo te lo agradece.} « Mange des fruits et des légumes : ton corps t'en remercie. »
\end{methode}

\begin{exoresolu}[Transformer des conseils en slogans à l'impératif]
Énoncé. Le club de santé du lycée de Yaoundé prépare une affiche contre le choléra. Mets les conseils suivants à l'impératif, en t'adressant aux élèves avec \esp{ustedes} : 1. \esp{lavarse las manos antes de comer} ; 2. \esp{beber agua potable} ; 3. \esp{no comer alimentos en la calle} ; 4. \esp{hervir el agua}.
\tcblower
Solution détaillée.

1. \esp{lavarse} $\to$ \esp{ustedes} : subjonctif \esp{laven} + pronom accolé \esp{se} + accent : \esp{\textbf{Lávense} las manos antes de comer.} « Lavez-vous les mains avant de manger. »

2. \esp{beber} $\to$ \esp{ustedes} : \esp{\textbf{Beban} agua potable.} « Buvez de l'eau potable. »

3. Conseil négatif : \esp{no} + subjonctif : \esp{\textbf{No coman} alimentos en la calle.} « Ne mangez pas d'aliments dans la rue. »

4. \esp{hervir} (changement \esp{e} $\to$ \esp{ie}) $\to$ \esp{\textbf{Hiervan} el agua.} « Faites bouillir l'eau. »

Vérification : formes affirmatives au subjonctif pour \esp{ustedes}, pronom accolé avec accent en 1, \esp{no} + subjonctif en 3.
\end{exoresolu}

\begin{exercice}[À toi de jouer]
Mets les verbes entre parenthèses à l'impératif, à la personne indiquée, pour compléter ces messages de prévention :

1. (\esp{cuidar}, \esp{tú}) \_\_\_\_\_\_ tu salud : es tu tesoro.

2. (\esp{no fumar}, \esp{tú}) \_\_\_\_\_\_ : el tabaco destruye los pulmones.

3. (\esp{hacer}, \esp{ustedes}) \_\_\_\_\_\_ ejercicio treinta minutos al día.

4. (\esp{dormir}, \esp{ustedes}) \_\_\_\_\_\_ bajo una mosquitera contra el paludismo.

5. (\esp{lavarse}, \esp{tú}) \_\_\_\_\_\_ las manos con agua y jabón.

6. (\esp{no olvidarse}, \esp{ustedes}) \_\_\_\_\_\_ de la vacuna.
\end{exercice}

\begin{corrige}[Exercice « À toi de jouer »]
1. \esp{Cuida} (3\textsuperscript{e} pers. sg. du présent). 2. \esp{No fumes} (\esp{no} + subjonctif). 3. \esp{Hagan} (impératif irrégulier de \esp{hacer} à \esp{ustedes}, du subjonctif \esp{haga}). 4. \esp{Duerman} (subjonctif avec changement \esp{o} $\to$ \esp{ue}). 5. \esp{Lávate} (pronom accolé + accent). 6. \esp{No se olviden} (pronom placé avant au négatif).
\end{corrige}

\begin{aretenir}[L'essentiel de la section]
\begin{itemize}
\item Impératif affirmatif \esp{tú} = 3\textsuperscript{e} personne du singulier du présent : \esp{cuida}, \esp{bebe}, \esp{duerme}.
\item Impératif négatif = \esp{no} + subjonctif : \esp{no fumes}, \esp{no bebas agua sucia}.
\item Irréguliers à connaître : \esp{haz}, \esp{ve}, \esp{pon}, \esp{sé}, \esp{ten}, \esp{sal}, \esp{di}, \esp{ven}.
\item Pronom réfléchi : accolé avec accent à l'affirmatif (\esp{lávate}), placé avant au négatif (\esp{no te laves}).
\item Parties du corps : article défini, pas de possessif : \esp{lávate \textbf{las} manos}.
\end{itemize}
\end{aretenir}

Nous savons désormais donner un ordre ou un conseil direct. Mais dans un message de prévention plus élaboré, on veut souvent \emph{nuancer} : « il faut que tu te laves les mains », « il convient que la population se vaccine ». C'est exactement le rôle du subjonctif après \esp{es necesario que} et \esp{conviene que}, que nous étudierons dans la section suivante.

\coursec{Grammaire 2 : le subjonctif avec « es necesario que » et « conviene que »}

Dans la section précédente, nous avons appris à donner des conseils directs avec l'\cle{impératif} : \esp{Lávate las manos.} « Lave-toi les mains. » Mais l'impératif s'adresse toujours à quelqu'un en face de nous. Or, dans un message de prévention, on veut souvent formuler un conseil \emph{général} ou viser une \emph{autre personne} : « Il faut que les jeunes fassent du sport », « Il convient que nous buvions de l'eau potable ». Pour cela, l'espagnol utilise une autre arme grammaticale : le \cle{subjonctif} introduit par des expressions impersonnelles comme \esp{es necesario que} ou \esp{conviene que}. C'est l'objet de cette section.

\courssub{Observation : deux phrases, deux structures}

\begin{textebox}[Mini-dialogue : chez le médecin scolaire]
— Doctor, ¿qué debemos hacer para estar sanos?\\
— Es necesario dormir ocho horas cada noche. Y es necesario que los jóvenes hagan ejercicio todos los días.\\
— ¿Y con la comida?\\
— Conviene que comáis frutas y verduras. También hay que beber agua potable.\\
— ¿Es importante que nos vacunemos contra el paludismo?\\
— Sí, es muy importante que todos se vacunen.
\end{textebox}

\textbf{Glossaire :} \esp{sano} = en bonne santé ; \esp{agua potable} = eau potable ; \esp{vacunarse} = se faire vacciner ; \esp{todos los días} = tous les jours.

Observons attentivement deux phrases du médecin :

\begin{itemize}
\item \esp{Es necesario \textbf{dormir} ocho horas.} → « Il faut dormir huit heures. » Verbe à l'\textbf{infinitif}.
\item \esp{Es necesario que los jóvenes \textbf{hagan} ejercicio.} → « Il faut que les jeunes fassent du sport. » Verbe conjugué après \esp{que}.
\end{itemize}

\begin{exemplebox}[Repérage de la différence]
Dans la première phrase, le conseil est \textbf{général} : il vaut pour tout le monde, personne en particulier n'est visé. Dans la seconde, le médecin désigne un \textbf{sujet précis} (\esp{los jóvenes}) différent de celui qui parle : l'espagnol introduit alors \esp{que} et change de mode verbal. Retenez cette opposition, elle est la clé de toute la section :
\begin{center}
\begin{tabularx}{\linewidth}{|X|X|}\hline
\rowcolor{popblueL} Conseil général (pas de sujet précis) & Sujet précis et différent \\ \hline
\esp{Es necesario dormir bien.} (infinitif) & \esp{Es necesario que tú duermas bien.} (que + subjonctif) \\ \hline
\esp{Hay que lavarse las manos.} & \esp{Es necesario que te laves las manos.} \\ \hline
\end{tabularx}
\end{center}
\end{exemplebox}

\courssub{La règle fondamentale : infinitif ou « que » + subjonctif}

\begin{propriete}[Règle du choix infinitif / subjonctif]
Après une expression impersonnelle de conseil ou de nécessité (\esp{es necesario}, \esp{es importante}, \esp{conviene}, \esp{es bueno}) :
\begin{enumerate}
\item Si le conseil est \textbf{général} ou si le sujet est \textbf{le même} que celui de l'expression → on emploie l'\textbf{infinitif} : \esp{Es necesario hacer ejercicio.} « Il faut faire du sport. »
\item Si le verbe a un \textbf{sujet précis et différent} → on emploie \esp{que} + \textbf{subjonctif présent} : \esp{Es necesario que tú hagas ejercicio.} « Il faut que toi, tu fasses du sport. »
\end{enumerate}
\end{propriete}

\begin{propriete}[Pourquoi le subjonctif ?]
Après \esp{es necesario que}, \esp{conviene que}, \esp{es importante que}, \esp{es bueno que}, le verbe se met au \textbf{subjonctif présent} parce qu'on n'énonce pas un fait réel, mais une \textbf{obligation}, une \textbf{nécessité} ou un \textbf{souhait} : l'action n'est pas encore accomplie, elle est seulement désirée ou exigée. C'est exactement la même logique qu'en français : « Il faut que tu \emph{aies} » (subjonctif français), et non « Il faut que tu \emph{as} » (indicatif).
\end{propriete}

\begin{popfigure}
\begin{tikzpicture}[
  node distance=0.9cm and 1.6cm,
  q/.style={draw=popblue, fill=popblueL, rounded corners, align=center, text width=4.6cm, font=\small},
  r/.style={draw=popgreen, fill=popgreenL, rounded corners, align=center, text width=5.2cm, font=\small},
  arr/.style={-{Stealth[length=2.5mm]}, thick, popdark}]
\node[q] (start) {Le conseil vise-t-il un \textbf{sujet précis différent} ?};
\node[r, below left=of start, xshift=-0.6cm, align=center] (oui){\textbf{OUI} → \esp{que} + \textbf{subjonctif}\\ \esp{Es necesario que tú duermas bien.}};
\node[r, below right=of start, xshift=0.6cm, align=center] (non){\textbf{NON} (général) → \textbf{infinitif}\\ \esp{Es necesario dormir bien.}\\ \esp{Hay que dormir bien.}};
\draw[arr] (start) -- (oui) node[midway, above left, font=\small\bfseries, popgreen] {sujet différent};
\draw[arr] (start) -- (non) node[midway, above right, font=\small\bfseries, poporange] {conseil général};
\end{tikzpicture}
\legende{Arbre de décision : infinitif ou « que » + subjonctif}
\end{popfigure}

\courssub{Formation du subjonctif présent}

\textbf{La recette de formation.} On part du radical de la \textbf{1\iere{} personne du singulier du présent de l'indicatif} (\esp{yo}), on enlève \esp{-o}, puis on ajoute les \textbf{terminaisons inversées} :

\begin{itemize}
\item verbes en \esp{-ar} → terminaisons en \textbf{-e} : \esp{hable, hables, hable, hablemos, habléis, hablen}
\item verbes en \esp{-er} / \esp{-ir} → terminaisons en \textbf{-a} : \esp{coma, comas, coma...} ; \esp{viva, vivas, viva...}
\end{itemize}

\begin{popfigure}
\begin{tikzpicture}[font=\small]
\matrix[matrix of nodes, nodes={draw=popline, minimum height=0.75cm, align=center, text depth=0.1ex},
  row 1/.style={nodes={fill=popblue, text=white, font=\small\bfseries}},
  column 1/.style={nodes={fill=popblueL, font=\small\bfseries}},
  row sep=-\pgflinewidth, column sep=-\pgflinewidth] (m) {
Personne & \esp{hablar} (hable) & \esp{comer} (coma) & \esp{vivir} (viva) & \esp{hacer} (haga) & \esp{dormir} (duerma) \\
\esp{yo} & hable & coma & viva & haga & duerma \\
\esp{tú} & hables & comas & vivas & hagas & duermas \\
\esp{él/ella/usted} & hable & coma & viva & haga & duerma \\
\esp{nosotros} & hablemos & comamos & vivamos & hagamos & durmamos \\
\esp{vosotros} & habléis & comáis & viváis & hagáis & durmáis \\
\esp{ellos/ustedes} & hablen & coman & vivan & hagan & duerman \\
};
\end{tikzpicture}
\legende{Subjonctif présent : verbes réguliers, irrégulier \esp{hacer} et verbe à changement vocalique \esp{dormir}}
\end{popfigure}

\begin{attention}[L'astuce du radical de « yo »]
Le radical se prend sur la 1\iere{} personne de l'indicatif, \emph{avec ses irrégularités} : \esp{hacer} → \esp{hago} → \esp{hag-} → \esp{haga} ; \esp{tener} → \esp{tengo} → \esp{teng-} → \esp{tenga} ; \esp{poder} → \esp{puedo} → \esp{pueda} ; \esp{dormir} → \esp{duermo} → \esp{duerma}. C'est pourquoi il faut bien connaître le présent de l'indicatif avant d'aborder le subjonctif.
\end{attention}

\textbf{Les verbes irréguliers indispensables de la leçon :}

\begin{center}
\begin{tabularx}{\linewidth}{|X|X|X|}\hline
\rowcolor{popblueL} Infinitif & Subjonctif (3\iere{} pers. sg.) & Exemple traduit \\ \hline
\esp{hacer} (faire) & \esp{haga} & \esp{Es necesario que hagas ejercicio.} « Il faut que tu fasses du sport. » \\ \hline
\esp{ir} (aller) & \esp{vaya} & \esp{Conviene que vayas al médico.} « Il convient que tu ailles chez le médecin. » \\ \hline
\esp{ser} (être) & \esp{sea} & \esp{Es importante que el agua sea potable.} « Il est important que l'eau soit potable. » \\ \hline
\esp{tener} (avoir) & \esp{tenga} & \esp{Es necesario que tengas cuidado.} « Il faut que tu fasses attention. » \\ \hline
\esp{dormir} (dormir) & \esp{duerma} & \esp{Es bueno que el niño duerma bajo un mosquitero.} « Il est bon que l'enfant dorme sous une moustiquaire. » \\ \hline
\esp{poder} (pouvoir) & \esp{pueda} & \esp{Es importante que puedas descansar.} « Il est important que tu puisses te reposer. » \\ \hline
\end{tabularx}
\end{center}

\textbf{Changements vocaliques} (e → ie, o → ue) : ils se produisent aux mêmes personnes qu'à l'indicatif (les trois singuliers et la 3\iere{} personne du pluriel) : \esp{pensar} → \esp{piense, pienses, piense, pensemos, penséis, piensen} ; \esp{poder} → \esp{pueda, puedas, pueda, podamos, podáis, puedan}. Notez \esp{dormir} : \esp{duerma} mais \esp{durmamos} (o → u à \esp{nosotros/vosotros}).

\textbf{Verbes en -ger / -gir :} le \esp{g} devient \esp{j} devant \esp{a} : \esp{proteger} → \esp{proteja} (\esp{Es necesario que te protejas de los mosquitos.} « Il faut que tu te protèges des moustiques. »), \esp{elegir} → \esp{elija}.

\textbf{Verbes en -car / -gar / -zar} (complément, niveau Bac) : \esp{sacar} → \esp{saque} ; \esp{llegar} → \esp{llegue} ; \esp{empezar} → \esp{empiece}. Ces changements garantissent la prononciation.

\courssub{Les structures synonymes pour conseiller}

\begin{aretenir}[La boîte à outils du conseil]
\begin{itemize}
\item \esp{Es necesario que} + subjonctif = il faut que, il est nécessaire que
\item \esp{Conviene que} + subjonctif = il convient que, il est bon que
\item \esp{Es importante que} + subjonctif = il est important que
\item \esp{Es bueno que} + subjonctif = il est bon que
\item \esp{Hay que} + \textbf{infinitif} = il faut (conseil général, JAMAIS de subjonctif)
\end{itemize}
\end{aretenir}

\begin{exemplebox}[Exemples traduits]
\begin{itemize}
\item \esp{Es necesario que te vacunes.} « Il faut que tu te fasses vacciner. »
\item \esp{Conviene que bebamos agua potable.} « Il convient que nous buvions de l'eau potable. »
\item \esp{Es importante que los niños duerman ocho horas.} « Il est important que les enfants dorment huit heures. »
\item \esp{Es bueno que comamos verduras cada día.} « Il est bon que nous mangions des légumes chaque jour. »
\item \esp{Hay que lavarse las manos antes de comer.} « Il faut se laver les mains avant de manger. »
\item \esp{Hay que hervir el agua para prevenir el cólera.} « Il faut faire bouillir l'eau pour prévenir le choléra. »
\end{itemize}
\end{exemplebox}

\courssub{Comparaison français / espagnol}

Bonne nouvelle : le français et l'espagnol fonctionnent de la \textbf{même manière} ici. « Il faut que tu te laves les mains » se traduit mot à mot : \esp{Es necesario que te laves las manos.} Le subjonctif existe dans les deux langues après « il faut que ». Seule la \emph{formation} diffère.

\begin{center}
\begin{tabularx}{\linewidth}{|X|X|X|}\hline
\rowcolor{popblueL} Français & Español & Remarque \\ \hline
Il faut dormir tôt. & \esp{Hay que dormir temprano.} & Conseil général : infinitif des deux côtés. \\ \hline
Il faut que tu dormes tôt. & \esp{Es necesario que duermas temprano.} & Sujet précis : que + subjonctif des deux côtés. \\ \hline
Il est important que nous nous protégions. & \esp{Es importante que nos protejamos.} & \esp{proteger} → \esp{protejamos} (g → j). \\ \hline
Il faut se faire vacciner. & \esp{Hay que vacunarse.} & \esp{hay que} + infinitif, jamais de « que ». \\ \hline
\end{tabularx}
\end{center}

\begin{methode}[Choisir la bonne structure en 3 étapes]
\begin{enumerate}
\item \textbf{Repérez le sujet du conseil.} Y a-t-il une personne précise, différente de celui qui parle (\esp{tú}, \esp{los jóvenes}, \esp{nosotros}...) ?
\item \textbf{Si oui} → construisez : expression impersonnelle + \esp{que} + verbe au \textbf{subjonctif présent}. Formez le subjonctif : radical de \esp{yo} + terminaison inversée (-ar → e ; -er/-ir → a).
\item \textbf{Si non} (conseil valable pour tous) → utilisez \esp{hay que} + infinitif, ou \esp{es necesario} + infinitif. C'est plus simple et très naturel.
\end{enumerate}
\end{methode}

\begin{exoresolu}[Traduisez en espagnol]
Énoncé. Traduisez : 1) « Il faut que les élèves fassent du sport. » 2) « Il faut boire de l'eau propre. » 3) « Il convient que tu ailles à l'hôpital. »
\tcblower
\textbf{Solution détaillée.}
\begin{enumerate}
\item Sujet précis : \esp{los alumnos}, différent du locuteur → \esp{que} + subjonctif. \esp{hacer} → \esp{hago} → \esp{hagan}. Réponse : \esp{Es necesario que los alumnos hagan deporte.}
\item Conseil général, aucun sujet précis → \esp{hay que} + infinitif. Réponse : \esp{Hay que beber agua limpia.}
\item Sujet précis \esp{tú} → \esp{que} + subjonctif. \esp{ir} est irrégulier : \esp{vayas}. Réponse : \esp{Conviene que vayas al hospital.}
\end{enumerate}
\end{exoresolu}

\begin{attention}[Les deux erreurs fatales des francophones]
\begin{enumerate}
\item \textbf{Indicatif après « es necesario que »} : *\esp{es necesario que tienes cuidado} est FAUX. Dès qu'il y a \esp{que}, le verbe passe au subjonctif : \esp{es necesario que \textbf{tengas} cuidado}. Le français vous aide : on dit « il faut que tu \emph{aies} », pas « il faut que tu \emph{as} ».
\item \textbf{Subjonctif après « hay que »} : *\esp{hay que te laves las manos} est IMPOSSIBLE. \esp{Hay que} est toujours suivi de l'\textbf{infinitif} : \esp{hay que lavarse las manos}. Si vous voulez un sujet précis, changez de structure : \esp{es necesario que te laves las manos}.
\end{enumerate}
\end{attention}

\begin{exercice}[À vous de jouer]
Mettez le verbe entre parenthèses à la forme correcte (infinitif ou subjonctif) :
\begin{enumerate}
\item \esp{Es necesario que tú} (vacunarse) \esp{contra la fiebre amarilla.}
\item \esp{Hay que} (hervir) \esp{el agua antes de beberla.}
\item \esp{Conviene que nosotros} (comer) \esp{mucha fruta.}
\item \esp{Es importante} (dormir) \esp{bajo un mosquitero.}
\item \esp{Es bueno que los jóvenes} (hacer) \esp{ejercicio cada día.}
\end{enumerate}
\end{exercice}

\begin{corrige}[Exercice « À vous de jouer »]
\begin{enumerate}
\item \esp{te vacunes} : sujet précis \esp{tú} → subjonctif de \esp{vacunarse}.
\item \esp{hervir} : \esp{hay que} + infinitif, toujours.
\item \esp{comamos} : sujet précis \esp{nosotros} → subjonctif de \esp{comer} (terminaison -a).
\item \esp{dormir} : conseil général sans sujet → infinitif.
\item \esp{hagan} : sujet précis \esp{los jóvenes} → subjonctif de \esp{hacer} (\esp{hago} → \esp{hagan}).
\end{enumerate}
\end{corrige}

\begin{savaistu}[Un subjonctif bien vivant]
En français, le subjonctif s'affaiblit : beaucoup de Francophones disent « il faut que tu \emph{viens} » au lieu de « il faut que tu \emph{viennes} ». En espagnol, au contraire, le subjonctif est \textbf{très vivant} : les Hispanophones l'emploient spontanément et correctement dans la conversation quotidienne, en Espagne comme au Cameroun hispanophone voisin, la Guinée équatoriale. Une erreur de subjonctif y est immédiatement remarquée. Maîtriser \esp{es necesario que tengas}, c'est parler un espagnol authentique — et c'est aussi une compétence précieuse pour rédiger les campagnes de santé publique, par exemple les messages de prévention du paludisme diffusés à la radio.
\end{savaistu}

\coursec{Méthode du commentaire et commentaire modèle}

Dans la section précédente, nous avons appris à manier le subjonctif après \esp{es necesario que} et \esp{conviene que} pour formuler des conseils. Ces structures, nous les avons rencontrées dans le texte \esp{Vivir sano}. Le moment est venu de changer de posture : au lieu d'étudier la langue phrase par phrase, nous allons apprendre à \cle{commenter} le texte dans son ensemble, c'est-à-dire à expliquer \emph{comment} il est construit et \emph{pourquoi} il produit un effet sur le lecteur. C'est exactement l'exercice du \esp{comentario de texto} demandé à l'examen.

\courssub{Qu'est-ce qu'un commentaire de texte ?}

\begin{definition}[Le commentaire de texte]
Le \cle{commentaire de texte} (\esp{comentario de texto}) est un exercice d'analyse écrite : il ne s'agit ni de résumer le texte, ni de donner son opinion personnelle, mais d'\textbf{expliquer comment le texte fonctionne} — ses idées, sa structure, ses procédés linguistiques — et quel effet il cherche à produire sur le lecteur. En espagnol, on le rédige entièrement dans la langue cible.
\end{definition}

Observez d'abord la différence entre ces deux phrases, écrites à propos de \esp{Vivir sano} :

\esp{El autor dice que hay que lavarse las manos y dormir bajo mosquitera.} « L'auteur dit qu'il faut se laver les mains et dormir sous moustiquaire. » → Ceci est un \textbf{résumé} : on répète le contenu. Note insuffisante.

\esp{El autor utiliza el imperativo « Cuida tu salud » para dirigirse directamente al lector y crear una relación de confianza.} « L'auteur utilise l'impératif "Prends soin de ta santé" pour s'adresser directement au lecteur et créer une relation de confiance. » → Ceci est une \textbf{analyse} : on relie un procédé (l'impératif) à un effet (l'interpellation du lecteur). C'est cela qui est attendu.

\courssub{Les quatre étapes du commentaire}

Un bon commentaire suit toujours la même architecture, en quatre étapes. Mémorisez-la comme un escalier : chaque marche repose sur la précédente.

\begin{popfigure}
\begin{tikzpicture}[font=\small, node distance=1.2cm and 1.5cm]
\node[draw=popblue, fill=popblueL, rounded corners, minimum width=4cm, minimum height=1cm, align=center] (e1) {\textbf{1. Introducción}\\ autor, fuente, tema};
\node[draw=popgreen, fill=popgreenL, rounded corners, minimum width=4cm, minimum height=1cm, align=center, below right=of e1] (e2) {\textbf{2. Plan aparente}\\ partes del texto};
\node[draw=poporange, fill=poporangeL, rounded corners, minimum width=4cm, minimum height=1cm, align=center, below right=of e2] (e3) {\textbf{3. Análisis}\\ ideas + procedimientos};
\node[draw=poppurple, fill=poppurpleL, rounded corners, minimum width=4cm, minimum height=1cm, align=center, below right=of e3] (e4) {\textbf{4. Conclusión}\\ balance + apertura};
\draw[-{Stealth}, thick, popmuted] (e1) -- (e2);
\draw[-{Stealth}, thick, popmuted] (e2) -- (e3);
\draw[-{Stealth}, thick, popmuted] (e3) -- (e4);
\end{tikzpicture}
\legende{L'escalier des quatre étapes du commentaire de texte}
\end{popfigure}

\begin{methode}[Réussir le commentaire en quatre étapes]
\begin{enumerate}
\item \textbf{Introduction} (\esp{introducción}) : présentez le texte en trois ou quatre phrases — qui écrit (ou quel type d'auteur), d'où vient le texte (article, affiche, discours), de quoi il parle (le thème), et à qui il s'adresse. Annoncez ensuite votre angle d'analyse.
\item \textbf{Plan apparent} (\esp{plan aparente}) : découpez le texte en deux ou trois grandes parties et donnez un titre à chacune. Montrez la progression : information → conseils → appel à l'action, par exemple.
\item \textbf{Analyse} (\esp{análisis}) : c'est le cœur du devoir. Pour chaque partie, procédez ainsi : \textbf{citez} le texte en espagnol entre guillemets, \textbf{traduisez} brièvement la citation si nécessaire, puis \textbf{analysez} le procédé et son effet. Jamais de paraphrase vide : chaque idée doit être reliée à un procédé linguistique (impératif, champ lexical, énumération, structure impersonnelle...).
\item \textbf{Conclusion} (\esp{conclusión}) : faites le bilan (quel est le but réel du texte ? a-t-il atteint son objectif ?) puis ouvrez sur une question plus large (la santé des jeunes au Cameroun, le rôle des campagnes de prévention...).
\end{enumerate}
\end{methode}

\courssub{Repérage des procédés dans « Vivir sano »}

Avant de rédiger, il faut \emph{observer}. Reprenons notre texte support et classons ses procédés. Ce travail de repérage est la matière première de votre analyse.

\begin{center}
\begin{tabularx}{\linewidth}{|X|X|X|}
\hline
\rowcolor{popblueL} \textbf{Procédé} & \textbf{Exemple cité du texte} & \textbf{Effet produit} \\
\hline
Impératif (ton injonctif) & \esp{Cuida tu salud}, \esp{Lava las manos} & S'adresse directement au lecteur, crée l'urgence et la proximité \\
\hline
Structures impersonnelles + subjonctif & \esp{es necesario que duermas bajo mosquitera} & Donne au conseil une valeur générale, objective, presque scientifique \\
\hline
Champ lexical de la santé & \esp{salud, enfermedad, vacuna, higiene, síntoma} & Construit le thème et crédibilise le discours \\
\hline
Énumérations & \esp{frutas, verduras, agua limpia...} & Donne une impression de richesse et de simplicité : les solutions existent \\
\hline
Exclamation finale & \esp{¡Tu salud es tu futuro!} & Marque les esprits, termine sur une émotion positive \\
\hline
\end{tabularx}
\end{center}

\begin{exemplebox}[De la citation à l'analyse : trois modèles]
\esp{El autor utiliza el imperativo « Cuida tu salud » para dirigirse directamente al lector.} « L'auteur utilise l'impératif "Prends soin de ta santé" pour s'adresser directement au lecteur. »

\esp{Con la estructura « es necesario que », el consejo se presenta como una verdad general, no como una simple opinión.} « Avec la structure "il faut que", le conseil se présente comme une vérité générale, et non comme une simple opinion. »

\esp{La exclamación final « ¡Tu salud es tu futuro! » resume el mensaje y motiva al lector joven.} « L'exclamation finale "Ta santé, c'est ton avenir !" résume le message et motive le jeune lecteur. »
\end{exemplebox}

\begin{attention}[L'erreur qui coûte le plus de points]
L'erreur la plus fréquente des candidats francophones est de \textbf{résumer le texte au lieu de l'analyser}. Écrire \esp{el texto habla del paludismo y del cólera} (« le texte parle du paludisme et du choléra ») n'apporte rien : l'examinateur a lu le texte. Chaque idée relevée doit être reliée à un \cle{procédé linguistique} et à un \cle{effet}. Formule magique : \emph{citation + procédé + effet}. Autre piège : ne jamais écrire \esp{el autor dice que...} en boucle ; variez avec \esp{afirma, subraya, insiste en, recuerda que}.
\end{attention}

\courssub{Les phrases-outils du commentateur}

Voici votre boîte à outils. Ces phrases toutes faites, correctes et élégantes, vous permettent de construire chaque partie du commentaire sans faute de langue.

\begin{center}
\begin{tabularx}{\linewidth}{|X|X|}
\hline
\rowcolor{popgreenL} \textbf{Español} & \textbf{Français} \\
\hline
\esp{El texto trata de la salud y los hábitos de vida.} & Le texte traite de la santé et des habitudes de vie. \\
\hline
\esp{Se trata de un texto de sensibilización dirigido a los jóvenes.} & Il s'agit d'un texte de sensibilisation destiné aux jeunes. \\
\hline
\esp{El texto se divide en tres partes.} & Le texte se divise en trois parties. \\
\hline
\esp{El autor quiere convencer al lector de que debe cambiar sus hábitos.} & L'auteur veut convaincre le lecteur qu'il doit changer ses habitudes. \\
\hline
\esp{Se observa el uso del imperativo para crear proximidad.} & On observe l'emploi de l'impératif pour créer de la proximité. \\
\hline
\esp{El campo léxico de la salud aparece en todo el texto.} & Le champ lexical de la santé apparaît dans tout le texte. \\
\hline
\esp{En conclusión, el texto informa y, sobre todo, invita a actuar.} & En conclusion, le texte informe et, surtout, invite à agir. \\
\hline
\end{tabularx}
\end{center}

Remarquez la tournure impersonnelle \esp{se observa} (« on observe ») : elle est très appréciée dans un commentaire car elle donne un ton objectif et académique, comme \esp{es necesario que} dans les conseils.

\courssub{Commentaire modèle guidé}

Lisons maintenant un commentaire modèle complet rédigé selon nos quatre étapes. Observez comment chaque citation est immédiatement exploitée et comment le plan apparent est explicité.

\begin{textebox}[Comentario modelo — « Vivir sano »]
\textbf{Introducción.} Este texto de sensibilización presenta la salud como un reto para los jóvenes. Probablemente destinado a escolares, trata de los hábitos de vida que permiten prevenir enfermedades como el paludismo o el cólera. El autor no se contenta con informar: busca transformar la conducta del lector.

\textbf{Plan aparente.} El texto se estructura en tres partes claras: una primera parte informativa (definición de salud y riesgos), una segunda parte prescriptiva (consejos concretos de higiene y alimentación) y una tercera parte exhortativa (llamada final a la acción).

\textbf{Análisis.} En la segunda parte, el texto enumera los gestos que salvan vidas: lavarse las manos, beber agua limpia, dormir bajo mosquitera. Mediante imperativos (« Cuida », « Lava ») y estructuras como « es necesario que », el autor no se limita a informar: quiere transformar los hábitos del lector. El imperativo crea un diálogo directo, casi familiar, mientras que la estructura impersonal da al consejo un valor de verdad general. La acumulación de ejemplos concretos muestra que la prevención es sencilla y está al alcance de todos.

\textbf{Conclusión.} En conclusión, el texto combina información y persuasión: su verdadero objetivo es que el joven lector pase del saber al acto. Este mensaje recuerda la importancia de las campañas de prevención en nuestras escuelas, en Camerún y en todo el mundo.
\end{textebox}

\begin{definition}[Glossaire du commentaire]
\esp{reto} = défi ; \esp{escolares} = élèves, scolaires ; \esp{conducta} = comportement ; \esp{gestos} = gestes ; \esp{mosquitera} = moustiquaire ; \esp{al alcance de todos} = à la portée de tous ; \esp{pasar del saber al acto} = passer du savoir à l'acte ; \esp{exhortativo} = exhortatif, incitatif ; \esp{prescriptivo} = prescriptif (qui donne des règles/conseils).
\end{definition}

Observez la mécanique : dans l'analyse, la citation \esp{« Cuida », « Lava »} est nommée (impératif), puis interprétée (dialogue direct, proximité), puis mise en relation avec un second procédé (\esp{es necesario que}). Rien n'est laissé gratuit.

\begin{exoresolu}[Analyser une phrase du texte]
Énoncé : commentez en deux ou trois phrases la phrase suivante tirée de \esp{Vivir sano} : \esp{« Conviene que hagas ejercicio cada día y que comas frutas y verduras. »}
\tcblower
Solution : \esp{El autor utiliza la estructura impersonal « conviene que » seguida del subjuntivo (« hagas », « comas ») para presentar el consejo como una recomendación objetiva y no como una orden.} « L'auteur utilise la structure impersonnelle "il convient que" suivie du subjonctif pour présenter le conseil comme une recommandation objective et non comme un ordre. » \esp{La coordinación de dos acciones cotidianas (« hacer ejercicio », « comer frutas ») muestra que la salud depende de pequeños hábitos sencillos, al alcance de cualquier joven.} « La coordination de deux actions quotidiennes montre que la santé dépend de petites habitudes simples, à la portée de n'importe quel jeune. » On remarque le schéma citation + procédé (subjonctif après \esp{conviene que}) + effet (conseil objectivé, gestes accessibles).
\end{exoresolu}

\courssub{Les critères d'évaluation à l'examen}

Au baccalauréat, votre commentaire est jugé selon trois grands critères. Gardez-les sous les yeux pendant la rédaction :

\begin{itemize}
\item \textbf{La fidélité au texte} (\esp{fidelidad al texto}) : tout ce que vous affirmez doit être prouvé par une citation exacte. N'inventez ni les intentions de l'auteur, ni des procédés absents.
\item \textbf{La qualité de la langue} (\esp{calidad de la lengua}) : espagnol correct, accents respectés, subjonctif bien employé, connecteurs variés (\esp{además, sin embargo, por eso, en conclusión}). Les phrases-outils apprises plus haut sont votre assurance.
\item \textbf{L'organisation logique} (\esp{organización lógica}) : les quatre étapes visibles, des paragraphes nets, une progression claire de l'introduction à la conclusion.
\end{itemize}

\begin{aretenir}[L'essentiel de la section]
Un commentaire = quatre étapes : \esp{introducción, plan aparente, análisis, conclusión}. La règle d'or de l'analyse : \textbf{citation + procédé + effet}, jamais de résumé. Les procédés clés d'un texte de prévention comme \esp{Vivir sano} : impératifs injonctifs, structures impersonnelles à subjonctif, champ lexical de la santé, énumérations, exclamation finale. À l'examen : fidélité au texte, langue soignée, organisation logique.
\end{aretenir}

\begin{exercice}[À vous de commenter]
À partir du texte \esp{Vivir sano} :
\begin{enumerate}
\item Rédigez une introduction complète (auteur, source supposée, thème, destinataires).
\item Choisissez deux citations du texte et analysez-les chacune en deux phrases (procédé + effet).
\item Rédigez une conclusion avec une ouverture sur la prévention des maladies dans votre établissement.
\end{enumerate}
\end{exercice}

\begin{corrige}[Exercice 1]
\textbf{1) Introduction possible :} \esp{El texto « Vivir sano » es un texto de sensibilización sobre la salud y los hábitos de vida. Dirigido probablemente a los jóvenes escolares, trata de la prevención de enfermedades como el paludismo y el cólera. El autor quiere convencer al lector de que la prevención es sencilla y necesaria.}

\textbf{2) Exemples d'analyse :}
\begin{itemize}
\item \esp{Se observa el uso del imperativo « Lava las manos » para crear un contacto directo con el lector.} « On observe l'emploi de l'impératif "Lave-toi les mains" pour créer un contact direct avec le lecteur. »
\item \esp{La estructura « es necesario que » + subjuntivo presenta el consejo como una verdad general.} « La structure "il faut que" + subjonctif présente le conseil comme une vérité générale. »
\end{itemize}

\textbf{3) Conclusion possible :} \esp{En conclusión, el texto informa y persuade al mismo tiempo. Su mensaje invita a reflexionar sobre el papel de las campañas de prevención en los liceos de Camerún.}
\end{corrige}

\coursec{Production guidée : messages de prévention}

Dans la section précédente, nous avons appris à \emph{lire} et à \emph{commenter} un texte sur la santé : repérer le thème, dégager la thèse de l'auteur, relever les procédés. Nous passons maintenant à l'étape la plus concrète de la leçon : \cle{produire} nous-mêmes un message de prévention en espagnol. C'est la compétence d'\emph{expression écrite et orale} du programme de Première A : à partir du lexique de la santé et des deux outils grammaticaux étudiés (l'impératif et le subjonctif après \esp{es necesario que} et \esp{conviene que}), vous allez concevoir une affiche de prévention contre une maladie réelle — paludisme, choléra, VIH/sida — puis la présenter à l'oral.

\courssub{Observer d'abord : à quoi ressemble un message de prévention ?}

Avant d'écrire, observons. Partout au Cameroun, en Espagne ou en Guinée équatoriale, les campagnes de santé publique suivent presque toujours la même architecture, très simple et très efficace. Lisez ce modèle original, rédigé dans le style des affiches du ministère de la Santé publique.

\begin{textebox}[Affiche modèle : campagne contre le paludisme]
\textbf{¡EL PALUDISMO MATA! PROTÉGETE CADA NOCHE}

El paludismo es una enfermedad grave transmitida por los mosquitos. Cada año, muchos niños y mujeres embarazadas mueren en nuestra región. Pero tú puedes actuar.

\begin{itemize}
\item Duerme siempre bajo una mosquitera impregnada.
\item Elimina el agua estancada cerca de tu casa.
\item Consulta al médico en cuanto tengas fiebre.
\item Es necesario que toda la familia respete estas medidas.
\end{itemize}

\textbf{TU SALUD, TU FUTURO. ¡JUNTOS CONTRA EL PALUDISMO!}
\end{textebox}

\textbf{Glossaire :} \esp{mosquitera impregnada} = moustiquaire imprégnée ; \esp{agua estancada} = eau stagnante ; \esp{en cuanto tengas fiebre} = dès que tu as de la fièvre (subjonctif après \esp{en cuanto}) ; \esp{medidas} = mesures ; \esp{mujeres embarazadas} = femmes enceintes.

\textbf{Questions d'observation :}
\begin{enumerate}
\item Combien de parties compte cette affiche ? Nommez-les.
\item Quels modes verbaux utilise-t-on pour les conseils ?
\item Pourquoi le titre est-il en majuscules et très court ?
\end{enumerate}

Vous l'avez remarqué : tout message de prévention efficace repose sur \textbf{trois blocs} : un \cle{titre accrocheur} qui nomme le danger, trois ou quatre \cle{conseils} sous forme de puces (\esp{viñetas}), et un \cle{slogan} final facile à mémoriser.

\begin{popfigure}
\begin{center}
\begin{tikzpicture}
\draw[line width=1.5pt, popblue, rounded corners=6pt] (0,0) rectangle (10,8);
\fill[popblueL] (0.3,6.4) rectangle (9.7,7.7);
\node[align=center, text=popdark, font=\bfseries] at (5,7.05) {TÍTULO ACCROCHEUR\\ \small (nomme le danger : ¡El paludismo mata!)};
\fill[poporangeL] (0.3,2.6) rectangle (9.7,6.1);
\node[align=center, text=popdark] at (5,5.6) {\textbf{CONSEJOS (viñetas)}};
\node[align=left, text=popdark] at (5,4.2) {\textbullet\ Consejo 1 : imperativo (Duerme bajo...)\\[2pt] \textbullet\ Consejo 2 : imperativo (Elimina...)\\[2pt] \textbullet\ Consejo 3 : imperativo (Consulta...)\\[2pt] \textbullet\ Consejo 4 : subjuntivo (Es necesario que...)};
\fill[popgreenL] (0.3,0.3) rectangle (9.7,2.2);
\node[align=center, text=popdark, font=\bfseries] at (5,1.25) {ESLOGAN FINAL\\ \small (court, mémorable : Tu salud, tu futuro)};
\end{tikzpicture}
\legende{Maquette d'une affiche de prévention : titre, conseils en puces, slogan.}
\end{center}
\end{popfigure}

\courssub{La méthode en quatre étapes}

\begin{methode}[Construire un message de prévention]
\begin{enumerate}
\item \textbf{Nommer le danger.} Choisis la maladie cible et écris un titre court qui la désigne clairement : \esp{Contra el cólera}, \esp{¡El paludismo mata!}, \esp{El VIH no perdona}. Le titre doit frapper en moins de six mots.
\item \textbf{Donner 3 ou 4 conseils à l'impératif.} Liste d'abord les gestes de prévention avec le lexique de la leçon, puis rédige chaque conseil à l'impératif affirmatif : \esp{Usa una mosquitera}, \esp{Hierve el agua}, \esp{Lávate las manos}.
\item \textbf{Justifier avec le subjonctif.} Ajoute une phrase d'explication ou d'insistance avec \esp{es necesario que} + subjonctif ou \esp{conviene que} + subjonctif : \esp{Es necesario que todos usen mosquiteras.} Cela montre que le conseil concerne toute la communauté.
\item \textbf{Conclure par un slogan.} Termine par une formule brève et rythmée : \esp{La prevención es la mejor medicina}, \esp{Tu salud, tu futuro}.
\end{enumerate}
\end{methode}

\courssub{Étape 1 : choisir la maladie et lister les gestes de prévention}

Travail de brainstorming individuel ou en binôme. Choisissez \textbf{une seule} maladie et notez quatre gestes de prévention en piochant dans le lexique de la leçon. Voici un tableau d'aide.

\begin{center}
\begin{tabularx}{\linewidth}{|X|X|X|}
\hline
\rowcolor{popblueL} \textbf{Enfermedad} & \textbf{Gestes de prévention (lexique)} & \textbf{Consejo posible} \\
\hline
El paludismo & mosquitera, agua estancada, fiebre, consulta médica & Duerme bajo una mosquitera. \\
\hline
El cólera & higiene, agua hervida, manos limpias, alimentos & Hierve el agua antes de beberla. \\
\hline
El VIH/sida & prueba, protección, fidelidad, información & Hazte la prueba del VIH. \\
\hline
\end{tabularx}
\end{center}

\begin{exemplebox}[Exemples traduits]
\esp{Contra el cólera: hierve el agua y lávate las manos.} = Contre le choléra : fais bouillir l'eau et lave-toi les mains.

\esp{Es necesario que todos usen mosquiteras.} = Il faut que tous utilisent des moustiquaires.

\esp{Conviene que te hagas la prueba del VIH.} = Il convient que tu fasses le test du VIH.
\end{exemplebox}

\courssub{Étape 2 : rédiger les conseils en alternant impératif et subjonctif}

C'est le cœur grammatical de la production. Rappel des deux outils :

\begin{itemize}
\item \textbf{L'impératif affirmatif} s'adresse directement au lecteur : \esp{Usa}, \esp{Hierve}, \esp{Consulta}, \esp{Elimina}. Avec les verbes pronominaux, le pronom se colle après le verbe et porte l'accent : \esp{lávate}, \esp{protégete}, \esp{hazte}.
\item \textbf{Le subjonctif} sert à généraliser le conseil à tout le monde après \esp{es necesario que} ou \esp{conviene que} : \esp{Es necesario que la gente hierva el agua.}
\end{itemize}

\begin{propriete}[L'impératif négatif (rappel)]
Pour interdire ou déconseiller, on utilise l'\textbf{impératif négatif} : \esp{No} + \textbf{subjonctif présent}.
\esp{No bebas agua sin hervir.} (Ne bois pas d'eau non bouillie.)
\esp{No te olvides de la mosquitera.} (N'oublie pas la moustiquaire.)
Les pronoms se placent \textbf{avant} le verbe : \esp{No te laves}, \esp{No se proteja}.
\end{propriete}

\begin{center}
\begin{tabularx}{\linewidth}{|X|X|X|}
\hline
\rowcolor{popblueL} \textbf{Français} & \textbf{Español} & \textbf{Remarque} \\
\hline
Fais bouillir l'eau. & Hierve el agua. & impératif affirmatif, tutoiement (\esp{tú}) \\
\hline
Il faut faire bouillir l'eau. & Es necesario hervir el agua. & infinitif : sens général, sans \esp{que} \\
\hline
Il faut que tu fasses bouillir l'eau. & Es necesario que hiervas el agua. & subjonctif : sujet précisé après \esp{que} \\
\hline
Lave-toi les mains. & Lávate las manos. & pronom soudé + accent tonique \\
\hline
Ne bois pas cette eau. & No bebas esta agua. & impératif négatif = subjonctif \\
\hline
\end{tabularx}
\end{center}

\begin{aretenir}[La règle d'or de l'alternance]
Dans une affiche, utilisez l'\textbf{impératif} pour les gestes concrets (3 conseils sur 4) et réservez \textbf{une phrase au subjonctif} pour engager la communauté entière : \esp{Es necesario que todos participen.} Cette alternance donne du rythme et montre votre maîtrise des deux modes.
\end{aretenir}

\begin{attention}[Erreur fréquente : mélanger les registres]
Ne mélangez jamais le tutoiement et le vouvoiement dans la même affiche. Fautif : \esp{Usa una mosquitera y consulte al médico.} (\esp{usa} = \esp{tú}, \esp{consulte} = \esp{usted}). Pour un public de jeunes, choisissez \esp{tú} et tenez-vous-y du début à la fin : \esp{Usa una mosquitera y consulta al médico.} Vérifiez aussi que tous les pronoms réfléchis suivent : \esp{lávate}, \esp{protégete}, jamais \esp{lávese} au milieu d'une affiche en \esp{tú}.
\end{attention}

\begin{savaistu}[Un seul mot peut sauver des vies]
\esp{Protégete} (« Protège-toi ») est le slogan officiel des campagnes de prévention du VIH dans tout le monde hispanophone. Un seul mot à l'impératif, avec son pronom soudé et son accent, peut sauver des vies : c'est la preuve que la grammaire n'est pas un exercice scolaire, mais un outil de communication réelle.
\end{savaistu}

\courssub{Étape 3 : ajouter un slogan mémorable}

Le slogan est la signature de votre affiche. Un bon slogan est \textbf{court} (moins de huit mots), \textbf{rythmé} et \textbf{positif}. Modèles à adapter :

\begin{itemize}
\item \esp{La prevención es la mejor medicina.} = La prévention est le meilleur médicament.
\item \esp{Tu salud, tu futuro.} = Ta santé, ton avenir.
\item \esp{Juntos contra el cólera.} = Ensemble contre le choléra.
\item \esp{Una mosquitera, una vida salvada.} = Une moustiquaire, une vie sauvée.
\end{itemize}

\courssub{Étape 4 : relecture guidée avec grille}

Avant de rendre votre affiche, relisez-la avec cette grille d'auto-évaluation. Chaque case cochée vaut un point.

\begin{center}
\begin{tabularx}{\linewidth}{|X|X|}
\hline
\rowcolor{popblueL} \textbf{Critère} & \textbf{Question à se poser} \\
\hline
Verbes corrects & Tous les impératifs sont-ils bien formés ? Le subjonctif suit-il bien \esp{es necesario que} / \esp{conviene que} ? \\
\hline
Registre constant & Ai-je gardé \esp{tú} (ou \esp{usted}) partout ? \\
\hline
Lexique varié & Ai-je utilisé au moins 5 mots du lexique de la santé ? \\
\hline
Orthographe & Accents sur \esp{lávate}, \esp{protégete}, \esp{médico}, \esp{hiérve} ? Signes ¡ ! et ¿ ? \\
\hline
Message clair & Un lecteur pressé comprend-il le danger et les gestes en 10 secondes ? \\
\hline
Slogan & Mon slogan est-il court et mémorable ? \\
\hline
\end{tabularx}
\end{center}

\courssub{Version orale : présenter son affiche en une minute}

\begin{experience}[Prise de parole en continu]
Chaque élève présente son affiche à la classe en une minute, sans lire mot à mot. Plan imposé :
\begin{enumerate}
\item \textbf{Saludo :} \esp{Buenos días a todos. Voy a presentar mi cartel contra...}
\item \textbf{El peligro :} \esp{El paludismo es una enfermedad grave porque...}
\item \textbf{Los consejos :} \esp{Primero, usa... Segundo, es necesario que...}
\item \textbf{El eslogan :} \esp{Recuerda: ¡Tu salud, tu futuro!}
\item \textbf{Despedida :} \esp{Gracias por su atención.}
\end{enumerate}
Critères d'évaluation orale : prononciation claire (le \esp{j} / \esp{g} (devant e,i) se prononce « r » guttural /x/ : \esp{juntos} ≈ « roun-tosse », \esp{gente} ≈ « renn-té »), regard vers le public, débit naturel.
\end{experience}

\courssub{Traduction guidée (thème) : du français vers l'espagnol}

Terminons par l'application des règles de la leçon à la traduction de phrases de conseil.

\begin{exoresolu}[Thème : trois phrases de prévention]
Traduisez en espagnol :
\begin{enumerate}
\item Contre le choléra, fais bouillir l'eau et lave-toi les mains.
\item Il faut que tout le monde dorme sous une moustiquaire.
\item Protège-toi : ta santé, c'est ton avenir.
\end{enumerate}
\tcblower
\textbf{Solution détaillée :}
\begin{enumerate}
\item \esp{Contra el cólera, hierve el agua y lávate las manos.} — Impératif de \esp{hervir} (e $\to$ ie : \esp{hierve}) ; verbe pronominal \esp{lavarse} : pronom soudé et accent sur \esp{lávate}.
\item \esp{Es necesario que todo el mundo duerma bajo una mosquitera.} — Après \esp{es necesario que}, subjonctif de \esp{dormir} (o $\to$ ue : \esp{duerma}). Attention : « tout le monde » se dit \esp{todo el mundo}, singulier en espagnol, verbe au singulier.
\item \esp{Protégete: tu salud es tu futuro.} — Impératif pronominal avec accent : \esp{protégete}. Le slogan reprend le modèle \esp{Tu salud, tu futuro}.
\end{enumerate}
Pièges évités : ne pas écrire \esp{hierve} à l'infinitif (\esp{hervir}) ; ne pas oublier l'accent de \esp{lávate} et \esp{protégete} ; ne pas traduire « il faut que » par \esp{es necesario de}.
\end{exoresolu}

\begin{exercice}[À vous de jouer]
Rédigez une affiche complète (titre, 4 conseils, slogan) contre le choléra ou le VIH/sida, en respectant la méthode en quatre étapes et la grille de relecture. Puis traduisez : « Il convient que les jeunes fassent le test du VIH. »
\end{exercice}

\begin{corrige}[Exercice : affiche et traduction]
Exemple d'affiche attendue (choléra) :
\begin{textebox}[Proposition d'affiche : choléra]
\textbf{¡EL CÓLERA NO ESPERA!}

\begin{itemize}
\item Hierve el agua antes de beberla.
\item Lávate las manos con jabón.
\item Lava bien las verduras y frutas.
\item Es necesario que toda la familia respete la higiene.
\end{itemize}

\textbf{LA PREVENCIÓN ES LA MEJOR MEDICINA.}
\end{textebox}

Traduction de la phrase : \esp{Conviene que los jóvenes se hagan la prueba del VIH.} — Subjonctif de \esp{hacerse} (3\textsuperscript{e} personne du pluriel : \esp{hagan} + pronom \esp{se} = \esp{se hagan}). Accent sur \esp{jóvenes} (paroxytonne terminant en -s).
\end{corrige}

\begin{aretenir}[L'essentiel de la production]
Un message de prévention = \textbf{titre} (le danger) + \textbf{3-4 conseils} (impératif, registre constant) + \textbf{une justification} (\esp{es necesario que} / \esp{conviene que} + subjonctif) + \textbf{un slogan} court. Relisez toujours avec la grille : verbes, registre, lexique, clarté. Vous êtes désormais capables, en espagnol, de conseiller et de protéger votre communauté : c'est la compétence finale du module « Environnement, santé et bien-être ».
\end{aretenir}

\newpage

\coursec{Activité d'intégration}

\courssub{Objectifs de l'activité}

À l'issue de cette activité d'intégration, l'élève sera capable de :

\begin{itemize}
\item réutiliser dans un contexte réel le \cle{lexique de la santé} étudié dans la leçon (\esp{salud, enfermedad, síntoma, vacuna, higiene, alimentación equilibrada, ejercicio, prevención}) ;
\item employer correctement l'\cle{impératif} pour donner des conseils directs à un public (\esp{¡Lávate las manos! ¡Duerme bajo una mosquitera!}) ;
\item construire des recommandations impersonnelles avec \esp{es necesario que} et \esp{conviene que} suivis du \cle{subjonctif présent} ;
\item concevoir, rédiger et présenter oralement un \cle{message de prévention} authentique, clair et adapté à un public réel ;
\item travailler en groupe, répartir les tâches et évaluer la production des autres selon des critères précis.
\end{itemize}

Cette activité mobilise les quatre savoir-faire de la leçon : lire et comprendre un texte sur la santé, analyser des habitudes de vie, manipuler les structures de conseil, et élaborer un message de prévention en espagnol. Elle prépare aussi directement aux épreuves de classe et du Baccalauréat : compréhension de l'écrit, expression écrite guidée et prise de parole en continu.

\courssub{Situation}

Ton lycée de Yaoundé organise sa \emph{Semana de la Salud} (semaine de la santé). Cette année, un événement spécial est prévu : une délégation d'élèves et de professeurs venus de \cle{Malabo}, en Guinée équatoriale — le seul pays africain hispanophone —, visite l'établissement. Le proviseur souhaite que les élèves de Première A, qui apprennent l'espagnol, préparent des affiches de prévention en espagnol pour décorer les couloirs et le CDI pendant cette semaine.

Voici le message que la présidente de l'association des élèves a publié sur le panneau d'affichage du lycée :

\begin{textebox}[Mensaje de la Asociación de Estudiantes — Lycée de Yaoundé]
¡Hola a todos!

La próxima semana celebramos la Semana de la Salud de nuestro liceo. Este año es muy especial: recibimos a estudiantes y profesores de Malabo, Guinea Ecuatorial, un país hermano de habla hispana.

Para esta ocasión, necesitamos carteles de prevención en español sobre tres enfermedades frecuentes en nuestra región: el paludismo, el cólera y el VIH/sida. Cada cartel debe tener un título atractivo, consejos claros para los estudiantes y un eslogan final memorable.

Recuerden: prevenir es mejor que curar. Una buena higiene, una alimentación equilibrada, el ejercicio físico y la información correcta salvan vidas.

El mejor cartel será expuesto en el CDI durante todo el trimestre y será presentado a nuestros invitados de Malabo.

¡Participen! ¡Su creatividad puede proteger la salud de todos!

La Asociación de Estudiantes
\end{textebox}

\textbf{Glossaire :} \esp{el liceo} : le lycée ; \esp{recibimos} : nous accueillons ; \esp{de habla hispana} : de langue espagnole ; \esp{el cartel} : l'affiche ; \esp{prevenir es mejor que curar} : mieux vaut prévenir que guérir ; \esp{salvan vidas} : sauvent des vies ; \esp{será expuesto} : sera exposé ; \esp{los invitados} : les invités ; \esp{recuerden} : rappelez-vous (impératif de \esp{ustedes}).

\begin{savaistu}[La Guinée équatoriale, voisine hispanophone du Cameroun]
La Guinée équatoriale, qui partage une frontière avec le Cameroun, est le seul pays d'Afrique subsaharienne où l'espagnol est langue officielle, héritage de la colonisation espagnole. Sa capitale, Malabo, se trouve sur l'île de Bioko. Apprendre l'espagnol au Cameroun, c'est donc aussi pouvoir communiquer avec un pays voisin : la langue de Cervantès est une langue de communication régionale, pas seulement européenne ou américaine.
\end{savaistu}

\courssub{Consignes}

Par \cle{groupes de 4}, vous allez créer une affiche de prévention en espagnol. Suivez les étapes dans l'ordre.

\begin{enumerate}
\item \textbf{Compréhension.} Relisez le message de la Asociación de Estudiantes. Relevez : le nom de l'événement, les trois maladies concernées, les quatre éléments obligatoires de l'affiche et la récompense prévue.

\item \textbf{Choix du sujet.} En groupe, choisissez \emph{une seule} maladie : \esp{el paludismo}, \esp{el cólera} ou \esp{el VIH/sida}. Répartissez les rôles : un coordinateur, un secrétaire (qui rédige), un correcteur linguistique (qui vérifie les accords et les conjugaisons), un présentateur.

\item \textbf{Recherche lexicale.} Complétez une fiche de vocabulaire : écrivez en espagnol au moins \cle{six mots} liés à votre maladie (symptômes, causes, moyens de prévention). Exemple pour le paludisme : \esp{la fiebre, el mosquito, la mosquitera, el dolor de cabeza, el médico, el tratamiento}.

\item \textbf{Rédaction du titre et des conseils.} Rédigez :
\begin{itemize}
\item un \cle{titre accrocheur} (une phrase courte, éventuellement avec une exclamation) ;
\item \cle{deux conseils à l'impératif} (tutoiement, verbes correctement conjugués) ;
\item \cle{deux recommandations} avec \esp{es necesario que} ou \esp{conviene que} + subjonctif présent ;
\item un \cle{slogan final} court et facile à mémoriser.
\end{itemize}

\item \textbf{Vérification linguistique.} Le correcteur du groupe vérifie : les accents (\esp{salud, prevención, alimentación, médico}), l'impératif des verbes irréguliers, le subjonctif après \esp{es necesario que}, les accords en genre et en nombre.

\item \textbf{Présentation orale.} Le présentateur (aidé si besoin par un camarade) présente l'affiche à la classe en \cle{2 minutes}, en présence de la délégation de Guinée équatoriale (jouée par le professeur et deux élèves). Il doit parler sans lire intégralement son texte.

\item \textbf{Vote et évaluation.} La classe vote pour la meilleure campagne selon la grille suivante : correction de la langue (4 points), respect des consignes (4 points), clarté du message de prévention (4 points), qualité de la présentation orale (4 points), créativité de l'affiche (4 points). L'affiche gagnante sera exposée au CDI.
\end{enumerate}

\begin{experience}[Déroulement en classe (2 heures)]
Séance 1 : étapes 1 à 5 (compréhension, choix, rédaction, vérification). Le professeur circule et aide les groupes. Séance 2 : étapes 6 et 7 (présentations orales, vote, remise du prix symbolique). Prévoir papier grand format, feutres de couleur et dictionnaires bilingues.
\end{experience}

\courssub{Ressources à mobiliser}

\textbf{Lexique utile :}

\begin{center}
\begin{tabularx}{\linewidth}{|X|X|}\hline
\rowcolor{popblueL} Español & Français \\ \hline
la salud / la enfermedad & la santé / la maladie \\ \hline
el paludismo / el cólera / el VIH/sida & le paludisme / le choléra / le VIH/sida \\ \hline
el síntoma / el dolor / la fiebre & le symptôme / la douleur / la fièvre \\ \hline
la vacuna / vacunarse & le vaccin / se faire vacciner \\ \hline
la higiene / lavarse las manos & l'hygiène / se laver les mains \\ \hline
la mosquitera / el mosquito & la moustiquaire / le moustique \\ \hline
el agua potable / hervir el agua & l'eau potable / faire bouillir l'eau \\ \hline
una alimentación equilibrada & une alimentation équilibrée \\ \hline
hacer ejercicio / hacer deporte & faire de l'exercice / du sport \\ \hline
consultar al médico / el tratamiento & consulter le médecin / le traitement \\ \hline
protegerse / prevenir / evitar & se protéger / prévenir / éviter \\ \hline
\end{tabularx}
\end{center}

\begin{attention}[Genre et article : deux pièges de ce tableau]
\begin{itemize}
\item \esp{el síntoma} est \textbf{masculin} malgré sa terminaison en -a (comme \esp{el problema, el tema}) : on dit \esp{los síntomas son graves}.
\item \esp{el agua} prend l'article \esp{el} au singulier (pour éviter le hiatus), mais reste \textbf{féminin} : \esp{el agua está \textbf{fría}}, \esp{el agua \textbf{potable}} ; au pluriel : \esp{las aguas}.
\end{itemize}
\end{attention}

\textbf{Grammaire 1 — l'impératif affirmatif (tú) :} on prend la 3\up{e} personne du singulier du présent de l'indicatif : \esp{lava} $\rightarrow$ \esp{¡Lava las manos!} ; \esp{come} $\rightarrow$ \esp{¡Come frutas!} ; \esp{duerme} $\rightarrow$ \esp{¡Duerme bajo una mosquitera!}. Attention aux irréguliers : \esp{hacer} $\rightarrow$ \esp{haz} ; \esp{ir} $\rightarrow$ \esp{ve} ; \esp{poner} $\rightarrow$ \esp{pon} ; \esp{tener} $\rightarrow$ \esp{ten} ; \esp{salir} $\rightarrow$ \esp{sal} ; \esp{ser} $\rightarrow$ \esp{sé} ; \esp{decir} $\rightarrow$ \esp{di} ; \esp{venir} $\rightarrow$ \esp{ven}.

\begin{propriete}[Place des pronoms à l'impératif affirmatif]
À l'impératif \textbf{affirmatif}, le pronom se place \textbf{après} le verbe et s'y \textbf{accole} : \esp{lavar} $\rightarrow$ \esp{¡Lávate las manos!} (« Lave-toi les mains ! »), \esp{protegerse} $\rightarrow$ \esp{¡Protégete!}. On ajoute alors un accent écrit pour conserver la prononciation (\esp{lávate, protégete}). À l'impératif \textbf{négatif}, le pronom revient \textbf{avant} le verbe, qui est au subjonctif : \esp{¡No \textbf{te laves} con agua sucia!} (« Ne te lave pas avec de l'eau sale ! »).
\end{propriete}

\textbf{Grammaire 2 — le subjonctif après une expression impersonnelle :}

\esp{Es necesario que} / \esp{Conviene que} + \textbf{subjonctif présent} : \esp{Es necesario que todos \textbf{beban} agua potable.} « Il est nécessaire que tout le monde boive de l'eau potable. » Rappel de formation : radical de la 1\up{re} personne du présent de l'indicatif + terminaisons inversées.

\begin{center}
\begin{tabular}{|l|l|l|}\hline
\rowcolor{popblueL} & -ar : \esp{lavarse} & -er : \esp{beber} \\ \hline
yo & me lave & beba \\ \hline
tú & te laves & bebas \\ \hline
él/ella/usted & se lave & beba \\ \hline
nosotros & nos lavemos & bebamos \\ \hline
vosotros & os lavéis & bebáis \\ \hline
ellos/ustedes & se laven & beban \\ \hline
\end{tabular}
\end{center}

Les verbes à radical irrégulier à la 1\up{re} personne gardent ce radical : \esp{hacer} (\esp{hago}) $\rightarrow$ \esp{haga} ; \esp{tener} (\esp{tengo}) $\rightarrow$ \esp{tengas} ; \esp{salir} (\esp{salgo}) $\rightarrow$ \esp{salga}. Exemple : \esp{Es necesario que \textbf{tengas} cuidado con el agua que bebes.} « Il faut que tu fasses attention à l'eau que tu bois. »

\begin{attention}[Pièges à éviter]
\begin{itemize}
\item Jamais d'infinitif après \esp{es necesario que} + sujet exprimé : on dit \esp{Es necesario que \textbf{te laves} las manos}, et non un infinitif. L'infinitif n'est possible que si aucun sujet n'est exprimé : \esp{Es necesario lavarse las manos.} « Il faut se laver les mains. »
\item \esp{Conviene que} est impersonnel : pas de sujet personnel avant (\esp{conviene que descanses} « il convient que tu te reposes »).
\item Ne pas traduire « se laver les mains » par un possessif : en espagnol, article défini + verbe pronominal : \esp{lavarse \textbf{las} manos} (et non un équivalent de « ses mains »).
\item Faux ami : \esp{estar constipado} signifie « être enrhumé », pas « être constipé » !
\item Ne pas oublier les signes d'ouverture ¡ et ¿ : \esp{¡Lávate las manos!}, \esp{¿Te has vacunado?}
\end{itemize}
\end{attention}

\courssub{Proposition de production et corrigé commenté}

\begin{exoresolu}[Affiche modèle : campaña contra el paludismo]
Énoncé. Rédige l'affiche complète d'un groupe qui a choisi le paludisme, en respectant toutes les contraintes (titre, deux impératifs, deux subjonctifs, slogan).
\tcblower
\textbf{Production modèle (texte de l'affiche) :}

\esp{¡ALTO AL PALUDISMO!}

\esp{El paludismo es una enfermedad grave transmitida por los mosquitos. Sus síntomas son la fiebre alta, el dolor de cabeza y el cansancio. ¡Pero podemos prevenirlo!}

\esp{• ¡Duerme siempre bajo una mosquitera!}

\esp{• ¡Elimina el agua estancada cerca de tu casa!}

\esp{• Es necesario que consultes al médico en cuanto tengas fiebre.}

\esp{• Conviene que toda la familia participe en la lucha contra los mosquitos.}

\esp{ESLOGAN: « Una mosquitera cada noche, una vida sana cada día. »}

\medskip
\textbf{Traduction :} « Halte au paludisme ! Le paludisme est une maladie grave transmise par les moustiques. Ses symptômes sont la forte fièvre, le mal de tête et la fatigue. Mais nous pouvons le prévenir ! • Dors toujours sous une moustiquaire ! • Élimine les eaux stagnantes près de ta maison ! • Il faut que tu consultes le médecin dès que tu as de la fièvre. • Il convient que toute la famille participe à la lutte contre les moustiques. Slogan : une moustiquaire chaque nuit, une vie saine chaque jour. »

\medskip
\textbf{Commentaire des choix de langue :}

\begin{itemize}
\item \textbf{Titre accrocheur :} \esp{¡Alto al paludismo!} utilise une exclamation courte avec les signes ¡ ! obligatoires en espagnol ; \esp{alto a...} est une formule de campagne authentique (« halte à... »).
\item \textbf{Impératifs :} \esp{Duerme} (dormir, verbe à changement vocalique o $\rightarrow$ ue) et \esp{Elimina} (verbe régulier en -ar) : formes correctes de l'impératif de \esp{tú}. Le conseil est direct, adapté à un public d'élèves.
\item \textbf{Subjonctifs :} \esp{Es necesario que consultes} : subjonctif présent de \esp{consultar} (terminaison -es pour un verbe en -ar). \esp{Conviene que toda la familia participe} : subjonctif de \esp{participar}, sujet à la 3\up{e} personne. La structure impersonnelle est respectée.
\item \textbf{Lexique de la leçon réemployé :} \esp{enfermedad, síntomas, fiebre, dolor de cabeza, prevenir, médico, vida sana}.
\item \textbf{Slogan :} phrase parallèle, rythmée et facile à mémoriser, avec le vocabulaire du bien-être (\esp{vida sana}).
\item \textbf{Points de vigilance respectés :} accents écrits (\esp{médico, estancada, familia}), article défini dans \esp{la lucha contra los mosquitos}, pas de calque du français.
\end{itemize}

\textbf{Présentation orale modèle (début) :} \esp{Buenos días a todos y bienvenidos a nuestros invitados de Malabo. Nuestro grupo ha elegido el paludismo porque es una enfermedad muy frecuente en Camerún y en Guinea Ecuatorial. Nuestro cartel presenta dos consejos y dos recomendaciones...} « Bonjour à tous et bienvenue à nos invités de Malabo. Notre groupe a choisi le paludisme parce que c'est une maladie très fréquente au Cameroun et en Guinée équatoriale. Notre affiche présente deux conseils et deux recommandations... »
\end{exoresolu}

\begin{exercice}[À ton tour : deux autres campagnes]
\begin{enumerate}
\item Rédige deux conseils à l'impératif et une recommandation avec \esp{es necesario que} pour une affiche contre \esp{el cólera} (pense à : \esp{hervir el agua, lavarse las manos, la higiene}).
\item Rédige un slogan final pour une affiche de prévention contre \esp{el VIH/sida}, puis une recommandation avec \esp{conviene que} + subjonctif.
\end{enumerate}
\end{exercice}

\begin{corrige}[Exercice : deux autres campagnes]
\textbf{1. Choléra (éléments de réponse possibles) :} \esp{¡Hierve el agua antes de beberla!} (« Fais bouillir l'eau avant de la boire ! » — impératif de \esp{hervir}, changement e $\rightarrow$ ie, pronom \esp{la} accolé avec accent écrit) ; \esp{¡Lávate las manos con jabón!} (« Lave-toi les mains avec du savon ! ») ; \esp{Es necesario que toda la familia respete las reglas de higiene.} (« Il faut que toute la famille respecte les règles d'hygiène. » — subjonctif \esp{respete}).

\textbf{2. VIH/sida (éléments de réponse possibles) :} slogan : \esp{« Infórmate, protégete, vive. »} (« Informe-toi, protège-toi, vis. ») ; recommandation : \esp{Conviene que todos los jóvenes hablen sin miedo de la prevención.} (« Il convient que tous les jeunes parlent sans peur de la prévention. » — subjonctif \esp{hablen}).
\end{corrige}

\courssub{Bilan de l'activité}

Cette activité d'intégration a permis de réunir, dans une tâche finale unique et motivante, l'ensemble des acquis de la leçon :

\begin{itemize}
\item \textbf{Compréhension de l'écrit :} lecture ciblée d'un message authentique d'association scolaire pour en extraire les informations utiles.
\item \textbf{Lexique :} mobilisation active du vocabulaire de la santé, de la maladie et de la prévention, réinvesti dans un document réel.
\item \textbf{Grammaire :} emploi contrôlé de l'impératif pour le conseil direct et du subjonctif après \esp{es necesario que} et \esp{conviene que} pour la recommandation impersonnelle — deux outils essentiels du discours de prévention.
\item \textbf{Expression écrite et orale :} conception d'une affiche structurée (titre, conseils, recommandations, slogan) et présentation orale devant un public, avec prise de parole en continu, à la manière des épreuves du Baccalauréat.
\item \textbf{Dimension culturelle et citoyenne :} l'espagnol devient une langue de communication réelle avec le monde hispanophone (la Guinée équatoriale, voisine du Cameroun), et les élèves découvrent que la langue apprise en classe peut servir à \cle{protéger la santé} de leur communauté.
\end{itemize}

\begin{aretenir}[L'essentiel]
Pour élaborer un message de prévention efficace en espagnol : un titre qui attire l'attention, des conseils directs à l'\cle{impératif} (\esp{¡Lávate las manos! ¡Hierve el agua!}), des recommandations avec \esp{es necesario que} / \esp{conviene que} + \cle{subjonctif} (\esp{Es necesario que bebas agua potable.}), et un slogan final mémorable. \esp{¡Prevenir es mejor que curar!} « Mieux vaut prévenir que guérir ! »
\end{aretenir}

\newpage

\coursec{Méthodes et exercices résolus}

\courssub{Méthode 1 : Lire et comprendre un texte sur la santé}

\begin{methode}[Réussir les questions de compréhension sur un texte de santé]
\begin{enumerate}
\item \cle{Lire deux fois} le texte : une première fois pour le sens général (qui ? où ? quel problème de santé ?), une seconde fois en soulignant le lexique-clé : \esp{salud}, \esp{enfermedad}, \esp{síntoma}, \esp{prevención}, \esp{vacuna}, \esp{higiene}.
\item \cle{Repérer la ligne} où se trouve la réponse : les questions suivent presque toujours l'ordre du texte.
\item \cle{Répondre en espagnol} avec une phrase complète, en reprenant les mots du texte sans les copier mot à mot (reformuler légèrement).
\item Pour une question \esp{¿Verdadero o falso?}, \cle{justifier} toujours par une citation courte du texte entre guillemets.
\item Pour une question de vocabulaire (\esp{Busca en el texto una palabra que signifique...}), vérifier que le mot trouvé a bien le sens demandé dans le contexte.
\item \cle{Relire} : accord du verbe avec le sujet, accents, ponctuation ¿ ... ?
\end{enumerate}
\textbf{Réflexes à retenir :} \esp{según el texto} = « selon le texte » ; \esp{el autor dice que...} = « l'auteur dit que... » ; pour donner son avis : \esp{en mi opinión}, \esp{pienso que}, \esp{me parece que}.
\end{methode}

\begin{exoresolu}[Compréhension de texte : la prévention du paludisme]
\textbf{Énoncé.} Lis le texte puis réponds en espagnol aux questions.

\emph{En muchos barrios de Yaundé, el paludismo sigue siendo una enfermedad muy frecuente, sobre todo durante la temporada de lluvias. Los síntomas más comunes son la fiebre alta, el dolor de cabeza y el cansancio. Sin embargo, esta enfermedad se puede prevenir con medidas sencillas: dormir bajo una mosquitera, eliminar el agua estancada cerca de las casas y consultar rápidamente al médico en caso de fiebre. Los especialistas recuerdan también que una buena higiene y una alimentación equilibrada fortalecen la salud de toda la familia.}

\begin{enumerate}
\item ¿Cuándo es más frecuente el paludismo en Yaundé?
\item Cita dos síntomas del paludismo mencionados en el texto.
\item ¿Verdadero o falso? Justifica : « El paludismo no se puede prevenir. »
\item Busca en el texto una palabra que signifique « moustiquaire ».
\end{enumerate}
\tcblower
\textbf{Raisonnement.} Les questions suivent l'ordre du texte : la réponse 1 se trouve à la deuxième ligne, la réponse 2 juste après, la question 3 demande de comparer une affirmation avec la phrase \esp{esta enfermedad se puede prevenir}, et la question 4 vise le mot \esp{mosquitera}.

\textbf{Réponses rédigées :}
\begin{enumerate}
\item \esp{El paludismo es más frecuente durante la temporada de lluvias.} (On reprend \esp{sobre todo durante la temporada de lluvias} en reformulant.)
\item \esp{Dos síntomas del paludismo son la fiebre alta y el dolor de cabeza.} (On pouvait aussi citer \esp{el cansancio}.)
\item \esp{Falso. El texto dice que «esta enfermedad se puede prevenir con medidas sencillas».} La justification par une citation courte est obligatoire.
\item \esp{La palabra es «mosquitera».} Attention au genre : \esp{la mosquitera} est féminin.
\end{enumerate}
\textbf{Conclusion.} Une phrase complète, un verbe conjugué correctement, une justification citée : c'est le trio qui rapporte tous les points.
\end{exoresolu}

\courssub{Méthode 2 : Traduire du français vers l'espagnol (thème) sur le thème de la santé}

\begin{methode}[Réussir le thème : phrases de prévention]
\begin{enumerate}
\item \cle{Identifier la structure} de la phrase française : conseil direct (impératif), conseil impersonnel (\esp{es necesario que} + subjonctif), ou simple constat (présent de l'indicatif).
\item \cle{Choisir le bon mode} : impératif pour un ordre direct (\esp{Lava las manos}), subjonctif après \esp{es necesario que}, \esp{conviene que}, \esp{es importante que}.
\item \cle{Traduire le vocabulaire exact} : « se laver les mains » = \esp{lavarse las manos} ; « dormir sous une moustiquaire » = \esp{dormir bajo una mosquitera} ; « une alimentation équilibrée » = \esp{una alimentación equilibrada} ; « faire du sport » = \esp{hacer ejercicio / hacer deporte}.
\item \cle{Éviter les calques} : « attraper une maladie » ne se traduit pas par *\esp{atrapar una enfermedad} mais par \esp{contraer una enfermedad} ou \esp{enfermar}.
\item \cle{Vérifier les accents} : \esp{salud} (sans accent), \esp{prevención}, \esp{higiene}, \esp{médico}, \esp{vacuna} (sans accent).
\item \cle{Relire} la phrase entière : ordre des mots, place du pronom réfléchi, accord.
\end{enumerate}
\end{methode}

\begin{exoresolu}[Thème : cinq phrases de prévention]
\textbf{Énoncé.} Traduis en espagnol :
\begin{enumerate}
\item Il faut se laver les mains avant de manger.
\item Il est nécessaire que les enfants dorment sous une moustiquaire.
\item Buvez de l'eau propre et mangez des légumes.
\item Il convient que la population se vaccine contre le choléra.
\item Le sport et une alimentation équilibrée protègent la santé.
\end{enumerate}
\tcblower
\textbf{Raisonnement phrase par phrase :}
\begin{enumerate}
\item « Il faut » + infinitif = \esp{hay que} + infinitif (structure impersonnelle, pas de subjonctif). Verbe réfléchi : \esp{lavarse}. « Avant de » + infinitif = \esp{antes de} + infinitif. Traduction : \esp{Hay que lavarse las manos antes de comer.}
\item « Il est nécessaire que » = \esp{es necesario que} + \cle{subjonctif}. Subjonctif présent de \esp{dormir} : \esp{duerman} (changement o $\to$ ue). Traduction : \esp{Es necesario que los niños duerman bajo una mosquitera.}
\item Conseil direct au pluriel (vous) = impératif \esp{vosotros} : \esp{bebed}, \esp{comed}. « Eau propre » = \esp{agua limpia} (attention : \esp{agua} est féminin malgré l'article \esp{el} au singulier ; ici sans article). Traduction : \esp{Bebed agua limpia y comed verduras.}
\item « Il convient que » = \esp{conviene que} + subjonctif. Subjonctif de \esp{vacunarse} : \esp{se vacune}. « Contre » = \esp{contra}. Traduction : \esp{Conviene que la población se vacune contra el cólera.}
\item Constat général = présent de l'indicatif. « Protègent » = \esp{protegen} (verbe en -er : \esp{proteger}, 3e personne du pluriel). Traduction : \esp{El deporte y una alimentación equilibrada protegen la salud.}
\end{enumerate}
\textbf{Conclusion.} Le choix du mode (indicatif / impératif / subjonctif) est la première décision ; le vocabulaire exact et les accents font la différence entre une bonne et une excellente copie.
\end{exoresolu}

\courssub{Méthode 3 : Traduire de l'espagnol vers le français (version)}

\begin{methode}[Réussir la version sur un texte de santé]
\begin{enumerate}
\item \cle{Lire tout le passage} avant d'écrire la moindre phrase : repérer le sujet de chaque verbe.
\item \cle{Traduire les temps justes} : présent espagnol = présent français ; \esp{se puede prevenir} = « on peut prévenir » (tournure passive espagnole rendue par « on » en français).
\item \cle{Attention aux faux amis} : \esp{una enfermedad} = « une maladie » (pas « une infirmité ») ; \esp{asistir al médico} = « aller chez le médecin » (pas « assister ») ; \esp{estar sano} = « être en bonne santé ».
\item \cle{Rendre les tournures idiomatiques} : \esp{en caso de fiebre} = « en cas de fièvre » ; \esp{sobre todo} = « surtout » ; \esp{sin embargo} = « cependant ».
\item \cle{Rédiger en français soigné} : la version est notée aussi sur la qualité du français.
\end{enumerate}
\end{methode}

\begin{exoresolu}[Version : un court passage sur l'hygiène]
\textbf{Énoncé.} Traduis en français :

\emph{La higiene es la primera medicina. En las escuelas de Douala, los profesores enseñan a los alumnos que es importante lavarse las manos con jabón, sobre todo antes de las comidas. Además, conviene que cada familia hierva el agua o la trate antes de beberla, porque el agua sucia transmite enfermedades graves como el cólera.}
\tcblower
\textbf{Raisonnement :}
\begin{itemize}
\item \esp{La higiene es la primera medicina} : formule proverbiale, traduction littérale correcte : « L'hygiène est le premier des médicaments » ou « la première médecine ».
\item \esp{enseñan a los alumnos que es importante lavarse} : \esp{enseñar a alguien que...} = « apprendre aux élèves que... » ; \esp{lavarse} = « se laver » (réfléchi).
\item \esp{conviene que cada familia hierva} : \esp{conviene que} + subjonctif = « il convient que chaque famille fasse bouillir » ; \esp{hierva} est le subjonctif de \esp{hervir} (e $\to$ ie), mais en français on traduit simplement par le subjonctif français : « qu'elle fasse bouillir ».
\item \esp{el agua sucia transmite} : \esp{agua} est féminin, donc « l'eau sale transmet » ; \esp{como el cólera} = « comme le choléra ».
\end{itemize}
\textbf{Traduction rédigée :}

« L'hygiène est le premier des médicaments. Dans les écoles de Douala, les professeurs apprennent aux élèves qu'il est important de se laver les mains avec du savon, surtout avant les repas. De plus, il convient que chaque famille fasse bouillir l'eau ou la traite avant de la boire, car l'eau sale transmet des maladies graves comme le choléra. »

\textbf{Conclusion.} Une version réussie est fidèle au sens, naturelle en français, et ne laisse passer aucun faux ami.
\end{exoresolu}

\courssub{Méthode 4 : Transformer une phrase : de l'impératif au subjonctif}

\begin{methode}[Passer du conseil direct au conseil impersonnel]
\begin{enumerate}
\item \cle{Identifier l'impératif} : \esp{Lava}, \esp{Duerme}, \esp{Vacúnate}, \esp{Haced}.
\item \cle{Choisir l'introducteur} : \esp{es necesario que}, \esp{conviene que}, \esp{es importante que}.
\item \cle{Ajouter un sujet} explicite si nécessaire : \esp{los niños}, \esp{tú}, \esp{la gente}.
\item \cle{Conjuguer au subjonctif présent} : radical de la 1re personne du présent (\esp{hago} $\to$ \esp{hag-}) + terminaisons inversées (-ar : e, es, e, emos, éis, en ; -er/-ir : a, as, a, amos, áis, an).
\item \cle{Surveiller les irréguliers} : \esp{dormir} $\to$ \esp{duerma} ; \esp{hacer} $\to$ \esp{haga} ; \esp{ir} $\to$ \esp{vaya} ; \esp{ser} $\to$ \esp{sea} ; \esp{tener} $\to$ \esp{tenga}.
\item \cle{Replacer les pronoms} : à l'impératif affirmatif le pronom est collé derrière (\esp{Vacúnate}) ; avec \esp{es necesario que}, il se place devant le verbe (\esp{que te vacunes}).
\end{enumerate}
\end{methode}

\begin{exoresolu}[Transformation : impératif $\to$ es necesario que + subjonctif]
\textbf{Énoncé.} Transforme chaque conseil direct en conseil impersonnel avec \esp{es necesario que} ou \esp{conviene que}, selon le modèle :

Modèle : \esp{Duerme bajo una mosquitera.} $\to$ \esp{Es necesario que duermas bajo una mosquitera.}
\begin{enumerate}
\item \esp{Lava las manos con jabón.} (tú)
\item \esp{Haced ejercicio todos los días.} (vosotros)
\item \esp{Vacúnate contra el cólera.} (tú)
\item \esp{Consulta al médico en caso de fiebre.} (usted)
\end{enumerate}
\tcblower
\textbf{Raisonnement et solutions :}
\begin{enumerate}
\item \esp{lavar} est régulier : subjonctif \esp{laves}. $\to$ \esp{Es necesario que laves las manos con jabón.}
\item \esp{hacer} est irrégulier : 1re personne \esp{hago} $\to$ radical \esp{hag-} ; 2e personne du pluriel : \esp{hagáis}. $\to$ \esp{Es necesario que hagáis ejercicio todos los días.}
\item Verbe pronominal \esp{vacunarse} : le pronom \esp{te} se déplace devant le verbe ; subjonctif \esp{vacunes}. $\to$ \esp{Conviene que te vacunes contra el cólera.}
\item \esp{consultar} régulier ; avec \esp{usted}, subjonctif 3e personne : \esp{consulte}. $\to$ \esp{Es necesario que consulte al médico en caso de fiebre.}
\end{enumerate}
\textbf{Conclusion.} Trois opérations suffisent : introducteur + sujet, verbe au subjonctif, pronoms déplacés devant le verbe. L'erreur classique est de garder la forme de l'impératif après \esp{que}.
\end{exoresolu}

\courssub{Méthode 5 : Rédiger un message de prévention en espagnol}

\begin{methode}[Construire un message de prévention clair et efficace]
\begin{enumerate}
\item \cle{Définir la cible et la maladie} : élèves d'un lycée, habitants d'un quartier ; paludisme, choléra, VIH/sida.
\item \cle{Structurer en trois parties} : un titre accrocheur (\esp{¡Protégete del paludismo!}) ; le rappel du danger (une ou deux phrases à l'indicatif) ; les conseils (trois à cinq impératifs ou structures avec subjonctif).
\item \cle{Alterner les structures de conseil} pour éviter la monotonie : impératif (\esp{Usa una mosquitera}), \esp{hay que} + infinitif, \esp{es necesario que} + subjonctif, \esp{conviene que} + subjonctif.
\item \cle{Employer le lexique de la leçon} : \esp{salud}, \esp{enfermedad}, \esp{síntoma}, \esp{vacuna}, \esp{higiene}, \esp{alimentación equilibrada}, \esp{ejercicio}, \esp{dolor}.
\item \cle{Terminer par une formule de mobilisation} : \esp{¡Tu salud está en tus manos!} ; \esp{¡Prevenir es mejor que curar!}
\item \cle{Compter les mots} demandés (généralement 60 à 100 mots) et relire : accents, ¿ ¡, accords.
\end{enumerate}
\end{methode}

\begin{exoresolu}[Rédaction : affiche de prévention contre le paludisme]
\textbf{Énoncé.} Tu es membre du club santé de ton lycée à Yaoundé. Rédige en espagnol un message de prévention contre le paludisme (70 à 90 mots) à afficher dans l'établissement. Utilise au moins deux impératifs et une structure avec \esp{es necesario que} ou \esp{conviene que}.
\tcblower
\textbf{Raisonnement.} Plan : titre + danger (indicatif) + trois conseils variés (impératif, \esp{hay que}, \esp{conviene que} + subjonctif) + slogan final. Vérifier : \esp{consultes} (subjonctif de \esp{consultar}), \esp{duerme} et \esp{elimina} (impératifs), accents sur \esp{está}, \esp{prevención} ; \esp{fiebre} ne prend pas d'accent.

\textbf{Message rédigé :}

\esp{¡Protégete del paludismo!}

\esp{El paludismo es una enfermedad grave, muy frecuente en la temporada de lluvias. Sus síntomas son la fiebre, el dolor de cabeza y el cansancio. Para proteger tu salud, sigue estos consejos: duerme siempre bajo una mosquitera y elimina el agua estancada cerca de tu casa. Hay que mantener una buena higiene en el barrio. Además, conviene que consultes al médico en cuanto tengas fiebre. La prevención salva vidas: ¡tu salud está en tus manos!}

(Environ 80 mots.)

\textbf{Vérification finale :} trois impératifs (\esp{sigue}, \esp{duerme}, \esp{elimina}), une structure \esp{hay que} + infinitif, une structure \esp{conviene que} + subjonctif (\esp{consultes}), lexique de la leçon présent (\esp{salud}, \esp{enfermedad}, \esp{síntomas}, \esp{higiene}, \esp{prevención}). Consignes respectées.

\textbf{Conclusion.} Un message de prévention efficace est court, structuré, varié dans ses structures de conseil et sans faute d'accent.
\end{exoresolu}

\begin{attention}[Les 5 erreurs qui coûtent des points au Bac]
\begin{enumerate}
\item \cle{Oublier le subjonctif} après \esp{es necesario que} et \esp{conviene que} : on écrit \esp{Es necesario que duermas} et jamais *\esp{que duermes}. L'indicatif après ces introducteurs est systématiquement sanctionné.
\item \cle{Les accents oubliés} : \esp{prevención}, \esp{alimentación}, \esp{médico}, \esp{está}, mais \esp{salud} et \esp{vacuna} sans accent. Un accent mal placé change parfois le sens ou le temps du verbe.
\item \cle{Les faux amis} : \esp{una enfermedad} = « une maladie » (pas « une infirmité ») ; \esp{asistir} = « assister à / se rendre » (pas « aider » : \esp{ayudar}) ; \esp{sano} = « en bonne santé ».
\item \cle{Les calques du français} : « attraper une maladie » = \esp{contraer una enfermedad} (jamais *\esp{atrapar}) ; « prendre un médicament » = \esp{tomar una medicina} (jamais *\esp{prender}) ; « faire du sport » = \esp{hacer deporte} (jamais *\esp{hacer sport}).
\item \cle{Le mauvais choix de mode ou de temps} : impératif pour un ordre direct, \esp{hay que} + infinitif pour une obligation générale, subjonctif pour un conseil personnalisé avec \esp{que}. Mélanger les trois dans la même phrase sans raison brouille le message et fait perdre des points de cohérence.
\end{enumerate}
\end{attention}

\newpage

\coursec{Exercices}

\courssub{Je vérifie mes connaissances}

\begin{exercice}[QCM : la santé et la prévention]
Pour chaque question, entoure la bonne réponse.

\begin{enumerate}
\item \esp{El paludismo} se transmite por :
\begin{enumerate}
\item el agua contaminada
\item la picadura de un mosquito
\item el aire
\end{enumerate}
\item \esp{Una vacuna} sirve para :
\begin{enumerate}
\item curar una herida
\item prevenir una enfermedad
\item dormir mejor
\end{enumerate}
\item \esp{Es necesario que tú \ldots} :
\begin{enumerate}
\item te lavas las manos
\item te laves las manos
\item te lavar las manos
\end{enumerate}
\item \esp{¡No \ldots\ cerca del agua estancada!} (dormir, tú) :
\begin{enumerate}
\item duermes
\item duermas
\item dormes
\end{enumerate}
\item \esp{Una alimentación equilibrada} significa :
\begin{enumerate}
\item comer mucho azúcar
\item comer de todo en cantidades razonables
\item no comer nada
\end{enumerate}
\end{enumerate}
\end{exercice}

\begin{exercice}[Vrai ou faux ? Justifie ta réponse]
Indique si chaque affirmation est vraie (V) ou fausse (F), puis justifie en une phrase en français.

\begin{enumerate}
\item L'impératif négatif s'exprime avec les formes du présent de subjonctif : \esp{¡No fumes!}
\item Après \esp{es necesario que}, on emploie l'indicatif.
\item \esp{Conviene que} signifie « il convient que, il est bon que ».
\item \esp{La higiene} se prononce « la i-jiène ».
\item Le verbe \esp{dormir} a pour impératif affirmatif à la 2\up{e} personne du singulier \esp{duerme}.
\end{enumerate}
\end{exercice}

\begin{exercice}[Vocabulaire : à chaque mot sa définition]
Relie chaque mot espagnol à sa définition française.

\begin{center}
\begin{tabularx}{\linewidth}{|X|X|}
\hline
\rowcolor{popblueL} \textbf{Español} & \textbf{Français} \\
\hline
1. \esp{el síntoma} & a. substance injectée pour prévenir une maladie \\
\hline
2. \esp{la vacuna} & b. signe qui révèle une maladie (fiebre, dolor\ldots) \\
\hline
3. \esp{el dolor} & c. ensemble des règles de propreté du corps \\
\hline
4. \esp{la higiene} & d. sensation physique désagréable, souffrance \\
\hline
5. \esp{el ejercicio} & e. activité physique qui maintient le corps en forme \\
\hline
\end{tabularx}
\end{center}
\end{exercice}

\begin{exercice}[Conjugaison : des conseils à compléter]
Complète chaque conseil avec la forme correcte du verbe entre parenthèses (impératif ou subjonctif, selon le cas).

\begin{enumerate}
\item ¡\ldots\ (hacer, tú) ejercicio todos los días!
\item ¡No \ldots\ (fumar, vosotros)! Es malo para la salud.
\item Es necesario que nosotros \ldots\ (dormir) bajo una mosquitera.
\item Conviene que tú \ldots\ (lavarse) las manos antes de comer.
\item ¡\ldots\ (vacunarse, ustedes) contra las enfermedades!
\item Es importante que los niños \ldots\ (beber) agua potable.
\end{enumerate}
\end{exercice}

\courssub{Je m'entraîne}

\begin{exercice}[Traduction : du français à l'espagnol]
Traduis les phrases suivantes en espagnol.

\begin{enumerate}
\item Il faut se laver les mains avec du savon.
\item Ne buvez pas l'eau de la rivière !
\item Il est nécessaire que les enfants dorment sous une moustiquaire.
\item Il convient que tu fasses du sport régulièrement.
\item Mangez des fruits et des légumes tous les jours !
\end{enumerate}
\end{exercice}

\begin{exercice}[Réécriture : de l'infinitif au subjonctif]
Transforme chaque conseil général en conseil personnel, selon le modèle.

\textbf{Modèle :} \esp{Es necesario beber agua potable.} $\rightarrow$ \esp{Es necesario que tú bebas agua potable.}

\begin{enumerate}
\item Es necesario dormir ocho horas. (tú)
\item Es importante lavarse las manos. (nosotros)
\item Es necesario vacunarse contra la gripe. (los niños)
\item Conviene hacer ejercicio físico. (usted)
\item Es importante comer verduras frescas. (vosotros)
\end{enumerate}
\end{exercice}

\begin{exercice}[Texte à trous : une campagne de santé]
Complète le texte suivant avec les mots de la liste : \esp{paludismo -- vacunas -- higiene -- síntomas -- mosquitera -- salud}.
\end{exercice}

\begin{textebox}[En el centro de salud de Yaoundé]
En nuestro barrio, el centro de salud organiza una semana de la \ldots\ . Los médicos explican a las familias cómo prevenir el \ldots\ , una enfermedad muy frecuente en Camerún. Los \ldots\ de esta enfermedad son la fiebre, el dolor de cabeza y el cansancio. Para protegerse, es necesario dormir bajo una \ldots\ y mantener una buena \ldots\ en la casa. Además, el centro ofrece \ldots\ gratuitas para los niños. ¡Participa y protege a tu familia!
\end{textebox}

\begin{exercice}[Appariement : problèmes et conseils]
Relie chaque problème (colonne A) au conseil qui convient (colonne B).

\begin{center}
\begin{tabularx}{\linewidth}{|X|X|}
\hline
\rowcolor{popgreenL} \textbf{A. Problema} & \textbf{B. Consejo} \\
\hline
1. Hay mosquitos en el barrio. & a. Hierva el agua antes de beberla. \\
\hline
2. El agua del pozo está sucia. & b. ¡No fume y haga deporte! \\
\hline
3. Quiero protegerme de la gripe. & c. Duerma bajo una mosquitera. \\
\hline
4. Tengo una vida demasiado sedentaria. & d. Vacúnese en el centro de salud. \\
\hline
5. Los niños comen con las manos sucias. & e. Lávense las manos antes de comer. \\
\hline
\end{tabularx}
\end{center}
\end{exercice}

\courssub{Je me prépare au Bac}

\begin{textebox}[La lucha contra el paludismo]
El paludismo es una enfermedad grave que afecta a millones de personas en África. Se transmite por la picadura de un mosquito infectado, generalmente durante la noche. Los síntomas principales son la fiebre alta, el dolor de cabeza, los vómitos y un gran cansancio. Los niños pequeños y las mujeres embarazadas son las personas más vulnerables.

Para prevenir esta enfermedad, los especialistas de la salud dan consejos sencillos. Primero, es necesario que toda la familia duerma bajo una mosquitera impregnada. Segundo, conviene eliminar el agua estancada cerca de las casas, porque allí nacen los mosquitos. Tercero, es importante que los enfermos vayan rápidamente al centro de salud para recibir un tratamiento.

En muchos países, como Camerún, el gobierno y las organizaciones internacionales distribuyen mosquiteras gratuitas y organizan campañas de información en las escuelas y en los mercados. Los jóvenes también pueden participar: hablen con sus familias, compartan los mensajes de prevención en las redes sociales y protejan a su comunidad. La salud es un tesoro: ¡cuidémosla juntos!
\end{textebox}

\emph{Glossaire :} \esp{la picadura} = la piqûre ; \esp{embarazada} = enceinte ; \esp{el agua estancada} = l'eau stagnante ; \esp{la mosquitera impregnada} = la moustiquaire imprégnée (d'insecticide) ; \esp{el tratamiento} = le traitement ; \esp{el tesoro} = le trésor.

\begin{exercice}[Compréhension écrite : le paludisme en Afrique]
Lis le texte \emph{La lucha contra el paludismo}, puis réponds aux questions.

\textbf{Questions de compréhension} (réponds en espagnol, avec des phrases complètes) :

\begin{enumerate}
\item ¿Cómo se transmite el paludismo?
\item ¿Cuáles son los síntomas principales de esta enfermedad? Cita tres.
\item ¿Quiénes son las personas más vulnerables?
\item ¿Por qué hay que eliminar el agua estancada?
\item ¿Qué pueden hacer los jóvenes para participar en la prevención?
\end{enumerate}

\textbf{Vrai ou faux} (justifie avec une phrase du texte) :

\begin{enumerate}
\item El mosquito pica sobre todo durante el día.
\item Las mosquiteras se distribuyen gratuitamente en algunos países.
\end{enumerate}

\textbf{Questions de langue} :

\begin{enumerate}
\item Trouve dans le texte deux phrases avec un subjonctif et recopie-les. Souligne le verbe au subjonctif.
\item Trouve dans le texte deux impératifs et recopie-les.
\item Donne en français l'équivalent de : \esp{mosquitera impregnada}, \esp{agua estancada}, \esp{centro de salud}.
\end{enumerate}
\end{exercice}

\begin{exercice}[Thème et version]
\textbf{A. Thème.} Traduis en espagnol le texte suivant (attention aux impératifs et aux subjonctifs) :

\begin{quote}
« Il faut que les élèves dorment sous une moustiquaire et il convient qu'ils fassent du sport. Ne fumez pas et lavez-vous les mains avant chaque repas ! La santé est très importante pour réussir à l'école. Il est nécessaire que tout le monde boive de l'eau potable. »
\end{quote}

\textbf{B. Version.} Traduis en français le texte \emph{Campaña de vacunación} ci-dessous.
\end{exercice}

\begin{textebox}[Campaña de vacunación]
Queridos ciudadanos: el hospital de Douala organiza una campaña de vacunación gratuita del 10 al 15 de mayo. Es necesario que todos los niños reciban sus vacunas. Conviene que los padres traigan el carnet de salud de sus hijos. ¡Vengan temprano y no tengan miedo: la vacuna no duele! Es importante que protejamos juntos a nuestra comunidad contra las enfermedades.
\end{textebox}

\begin{exercice}[Expression écrite : un message de prévention]
\textbf{Sujet :} Le lycée de Nkongsamba organise une semaine de la santé. Tu es membre du club d'espagnol. Rédige, en \textbf{80 à 100 mots}, un message de prévention contre le \textbf{choléra} destiné aux réseaux sociaux du lycée.

\textbf{Consignes :}
\begin{enumerate}
\item Commence par une accroche qui attire l'attention des élèves (ex. \esp{¡Atención, jóvenes!}).
\item Rappelle brièvement comment se transmet la maladie.
\item Donne au moins \textbf{quatre conseils} en utilisant des impératifs (\esp{¡Lávate\ldots!}, \esp{¡No bebas\ldots!}) et des structures avec subjonctif (\esp{Es necesario que\ldots}, \esp{Conviene que\ldots}).
\item Termine par une phrase d'encouragement.
\item Soigne l'orthographe et les accents.
\end{enumerate}
\end{exercice}

\begin{methode}[Rédiger un message de prévention en cinq étapes]
\begin{enumerate}
\item \textbf{Accroche} : interpelle le lecteur (\esp{¡Atención, jóvenes!}, \esp{¡Cuida tu salud!}).
\item \textbf{Information} : une ou deux phrases sur la maladie et son mode de transmission (\esp{El cólera se transmite por el agua y los alimentos contaminados.}).
\item \textbf{Conseils} : alterne impératifs (\esp{¡Lávate las manos!}, \esp{¡Hierve el agua!}) et subjonctifs (\esp{Es necesario que\ldots}, \esp{Conviene que\ldots}).
\item \textbf{Conclusion} : encourage et appelle à l'action (\esp{¡Juntos podemos vencer al cólera!}).
\item \textbf{Relecture} : vérifie les accents, les points d'exclamation ouvrants (¡) et la concordance des verbes.
\end{enumerate}
\end{methode}

\begin{exemplebox}[Modèle de message de prévention]
\esp{¡Atención, jóvenes del liceo de Nkongsamba! El cólera es una enfermedad peligrosa que se transmite por el agua y los alimentos contaminados. Para protegerte, sigue estos consejos: ¡lávate las manos con jabón antes de cada comida! ¡No bebas agua de la rivière! Es necesario que hiervas el agua antes de beberla. Conviene que laves bien las frutas y las verduras. Si tienes diarrea o vómitos, ve rápidamente al centro de salud. La prevención es la mejor medicina: ¡juntos podemos proteger nuestro liceo y nuestras familias!}

\emph{Traduction :} « Attention, jeunes du lycée de Nkongsamba ! Le choléra est une maladie dangereuse qui se transmet par l'eau et les aliments contaminés. Pour te protéger, suis ces conseils : lave-toi les mains avec du savon avant chaque repas ! Ne bois pas l'eau de la rivière ! Il faut que tu fasses bouillir l'eau avant de la boire. Il convient que tu laves bien les fruits et les légumes. Si tu as des diarrhées ou des vomissements, va vite au centre de santé. La prévention est le meilleur médicament : ensemble, nous pouvons protéger notre lycée et nos familles ! »
\end{exemplebox}

\begin{attention}[Pièges à éviter dans ta rédaction]
\begin{itemize}
\item Ne dis pas \esp{la cólera} : en espagnol, \esp{el cólera} est masculin (et prend un accent sur le \esp{ó}).
\item « Faire bouillir » se dit \esp{hervir} (e $\rightarrow$ ie) : \esp{¡Hierve el agua!}, \esp{es necesario que hiervas el agua}.
\item Après \esp{es necesario que} / \esp{conviene que}, jamais d'indicatif : \esp{es necesario que \textbf{te laves}}, et non \esp{te lavas}.
\item N'oublie pas le point d'exclamation ouvrant : \esp{¡Lávate las manos!}
\item \esp{La salud} (la santé) est féminin ; \esp{el dolor} (la douleur) est masculin : méfie-toi des faux amis de genre.
\end{itemize}
\end{attention}

\newpage

\coursec{Corrigés des exercices}

\begin{corrige}[Exercice 1 — QCM : la santé et la prévention]
\begin{enumerate}
\item \textbf{b} — \esp{la picadura de un mosquito} : le paludisme se transmet par la piqûre d'un moustique infecté (l'anophèle femelle), qui pique surtout la nuit.
\item \textbf{b} — \esp{prevenir una enfermedad} : un vaccin sert à \emph{prévenir} une maladie, pas à la guérir. Verbe associé : \esp{vacunarse} (se faire vacciner).
\item \textbf{b} — \esp{te laves} : après \esp{es necesario que}, on emploie le \textbf{subjonctif présent}. \esp{Lavarse} : que yo me lave, que tú te laves, que él se lave\ldots
\item \textbf{b} — \esp{duermas} : l'impératif négatif emprunte les formes du subjonctif présent : \esp{¡No duermas cerca del agua estancada!} Attention, \esp{duermes} est l'indicatif présent et \esp{dormes} n'existe pas en espagnol.
\item \textbf{b} — une \esp{alimentación equilibrada}, c'est manger de tout en quantités raisonnables : fruits, légumes, céréales, protéines, et peu de sucre.
\end{enumerate}
\end{corrige}

\begin{corrige}[Exercice 2 — Vrai ou faux ? Justifie ta réponse]
\begin{enumerate}
\item \textbf{Vrai.} L'impératif négatif reprend exactement les formes du présent de subjonctif : \esp{¡No fumes!}, \esp{¡No bebas esa agua!}, \esp{¡No comáis con las manos sucias!}
\item \textbf{Faux.} Après \esp{es necesario que} (comme après toutes les expressions impersonnelles de nécessité ou de conseil), on emploie le \textbf{subjonctif} : \esp{Es necesario que duermas bajo una mosquitera.} L'indicatif n'est possible que sans \esp{que} : \esp{Es necesario dormir} (infinitif, conseil général).
\item \textbf{Vrai.} \esp{Conviene que} signifie « il convient que, il est bon que » ; il est suivi du subjonctif : \esp{Conviene que hagas ejercicio.} « Il convient que tu fasses de l'exercice. »
\item \textbf{Faux.} \esp{La higiene} se prononce « la i-\textbf{jié}-né » : le \esp{h} est muet, le \esp{g} devant \esp{e} se prononce « j » (comme le \emph{j} français de « je »), et l'accent tonique tombe sur \esp{-gie-}.
\item \textbf{Vrai.} \esp{Dormir} est un verbe à changement radical (o $\rightarrow$ ue) : impératif affirmatif \esp{¡Duerme!} (tú). Mais attention : à la forme négative, on dit \esp{¡No duermas!}, et à \esp{ustedes} : \esp{¡Duerman!}
\end{enumerate}
\end{corrige}

\begin{corrige}[Exercice 3 — Vocabulaire : à chaque mot sa définition]
1 -- b ; 2 -- a ; 3 -- d ; 4 -- c ; 5 -- e.

\emph{Remarques :}
\begin{itemize}
\item Attention au genre : \esp{el síntoma} et \esp{el dolor} sont masculins ; \esp{la vacuna}, \esp{la higiene}, \esp{la salud} sont féminins.
\item Pluriels : \esp{el dolor} $\rightarrow$ \esp{los dolores} ; \esp{el síntoma} $\rightarrow$ \esp{los síntomas}.
\item Exemple : \esp{La fiebre es un síntoma del paludismo.} « La fièvre est un symptôme du paludisme. »
\end{itemize}
\end{corrige}

\begin{corrige}[Exercice 4 — Conjugaison : des conseils à compléter]
\begin{enumerate}
\item ¡\textbf{Haz} ejercicio todos los días! — impératif irrégulier de \esp{hacer} à \esp{tú} (à mémoriser : \esp{haz}, comme \esp{di}, \esp{ve}, \esp{pon}, \esp{sal}, \esp{ten}, \esp{ven}, \esp{sé}).
\item ¡No \textbf{fuméis}! — impératif négatif = subjonctif présent, 2\up{e} pers. plur. ; verbe en \esp{-ar} : voyelle \esp{-éis}.
\item Es necesario que nosotros \textbf{durmamos} bajo una mosquitera. — subjonctif de \esp{dormir} avec changement o $\rightarrow$ ue.
\item Conviene que tú \textbf{te laves} las manos antes de comer. — subjonctif pronominal : le pronom \esp{te} se place devant le verbe conjugué.
\item ¡\textbf{Vacúnense} contra las enfermedades! — impératif affirmatif pronominal à \esp{ustedes} : le pronom \esp{se} se place après le verbe et s'y soude ; on ajoute un accent sur la voyelle tonique (\esp{vacúnense}).
\item Es importante que los niños \textbf{beban} agua potable. — subjonctif, 3\up{e} pers. plur. de \esp{beber} (voyelle \esp{-a} pour les verbes en \esp{-er/-ir}).
\end{enumerate}
\end{corrige}

\begin{corrige}[Exercice 5 — Traduction : du français à l'espagnol]
\begin{enumerate}
\item \esp{Hay que lavarse las manos con jabón.} (ou \esp{Es necesario lavarse las manos con jabón.}) — infinitif, car le conseil est général, sans sujet précis.
\item \esp{¡No bebáis el agua del río!} — impératif négatif, 2\up{e} pers. plur. ; attention à l'accent de \esp{río} et à la contraction \esp{de + el = del}.
\item \esp{Es necesario que los niños duerman bajo una mosquitera.} — subjonctif de \esp{dormir} : \esp{duerman} (changement o $\rightarrow$ ue).
\item \esp{Conviene que hagas deporte regularmente.} — subjonctif irrégulier de \esp{hacer} : \esp{haga, hagas, haga, hagamos, hagáis, hagan}.
\item \esp{¡Comed frutas y verduras todos los días!} — impératif affirmatif régulier, 2\up{e} pers. plur. de \esp{comer} (infinitif sans \esp{-r} + \esp{-d}).
\end{enumerate}
\emph{Point de vigilance :} « il faut » impersonnel se rend par \esp{hay que} + infinitif ; « il faut que » + sujet se rend par \esp{es necesario que} + subjonctif.
\end{corrige}

\begin{corrige}[Exercice 6 — Réécriture : de l'infinitif au subjonctif]
\begin{enumerate}
\item \esp{Es necesario que tú duermas ocho horas.}
\item \esp{Es importante que nosotros nos lavemos las manos.}
\item \esp{Es necesario que los niños se vacunen contra la gripe.}
\item \esp{Conviene que usted haga ejercicio físico.}
\item \esp{Es importante que vosotros comáis verduras frescas.}
\end{enumerate}
\emph{Rappel de la règle :} quand le conseil a un sujet précis, on ajoute \esp{que} + subjonctif ; quand il est général (valable pour tout le monde), l'infinitif suffit. Comparer : \esp{Es necesario vacunarse.} (général) / \esp{Es necesario que los niños se vacunen.} (sujet précis). Attention aux verbes pronominaux : le pronom réfléchi s'accorde avec le sujet (\esp{nos lavemos}, \esp{se vacunen}).
\end{corrige}

\begin{corrige}[Exercice 7 — Texte à trous : une campagne de santé]
Dans l'ordre : \esp{salud} -- \esp{paludismo} -- \esp{síntomas} -- \esp{mosquitera} -- \esp{higiene} -- \esp{vacunas}.

\emph{Texte complété :} \esp{En nuestro barrio, el centro de salud organiza una semana de la salud. Los médicos explican a las familias cómo prevenir el paludismo, una enfermedad muy frecuente en Camerún. Los síntomas de esta enfermedad son la fiebre, el dolor de cabeza y el cansancio. Para protegerse, es necesario dormir bajo una mosquitera y mantener una buena higiene en la casa. Además, el centro ofrece vacunas gratuitas para los niños. ¡Participa y protege a tu familia!}

\emph{Traduction :} « Dans notre quartier, le centre de santé organise une semaine de la santé. Les médecins expliquent aux familles comment prévenir le paludisme, une maladie très fréquente au Cameroun. Les symptômes de cette maladie sont la fièvre, le mal de tête et la fatigue. Pour se protéger, il faut dormir sous une moustiquaire et maintenir une bonne hygiène dans la maison. De plus, le centre offre des vaccins gratuits pour les enfants. Participe et protège ta famille ! »

\emph{Indices de résolution :} « una semana de la\ldots » appelle un nom féminin abstrait (\esp{salud}) ; « los\ldots\ son la fiebre\ldots » appelle un pluriel (\esp{síntomas}) ; « dormir bajo una\ldots » évoque immédiatement \esp{mosquitera}.
\end{corrige}

\begin{corrige}[Exercice 8 — Appariement : problèmes et conseils]
1 -- c ; 2 -- a ; 3 -- d ; 4 -- b ; 5 -- e.

\emph{Justifications :}
\begin{itemize}
\item 1 -- c : contre les moustiques, on dort sous une moustiquaire.
\item 2 -- a : si l'eau du puits est sale, on la fait bouillir avant de la boire.
\item 3 -- d : contre la grippe, on se fait vacciner au centre de santé.
\item 4 -- b : contre la sédentarité, on fait du sport (et on évite le tabac).
\item 5 -- e : avant de manger, on se lave les mains.
\end{itemize}
\emph{Remarque grammaticale :} les conseils de la colonne B sont à l'impératif de politesse (\esp{usted} / \esp{ustedes}) : \esp{hierva}, \esp{duerma}, \esp{vacúnese}, \esp{lávense}. C'est le registre utilisé dans les campagnes officielles de santé publique. Notez la place du pronom après le verbe à l'impératif affirmatif : \esp{vacúnese}, \esp{lávense} (avec accent écrit).
\end{corrige}

\begin{corrige}[Exercice 9 — Compréhension écrite : le paludisme en Afrique]
\textbf{Questions de compréhension :}
\begin{enumerate}
\item \esp{El paludismo se transmite por la picadura de un mosquito infectado, generalmente durante la noche.}
\item \esp{Los síntomas principales son la fiebre alta, el dolor de cabeza y los vómitos.} (on peut aussi citer \esp{un gran cansancio})
\item \esp{Las personas más vulnerables son los niños pequeños y las mujeres embarazadas.}
\item \esp{Hay que eliminar el agua estancada porque allí nacen los mosquitos.}
\item \esp{Los jóvenes pueden hablar con sus familias, compartir los mensajes de prevención en las redes sociales y proteger a su comunidad.}
\end{enumerate}

\textbf{Vrai ou faux :}
\begin{enumerate}
\item \textbf{Falso.} Justification : \esp{Se transmite por la picadura de un mosquito infectado, generalmente durante la noche.} Le moustique pique surtout la nuit, pas le jour.
\item \textbf{Verdadero.} Justification : \esp{El gobierno y las organizaciones internacionales distribuyen mosquiteras gratuitas.}
\end{enumerate}

\textbf{Questions de langue :}
\begin{enumerate}
\item Subjonctifs : \esp{Es necesario que toda la familia \textbf{duerma} bajo una mosquitera impregnada.} ; \esp{Es importante que los enfermos \textbf{vayan} rápidamente al centro de salud.} (subjonctifs après \esp{es necesario que} et \esp{es importante que} ; \esp{vayan} est le subjonctif irrégulier de \esp{ir}).
\item Impératifs : \esp{hablen}, \esp{compartan}, \esp{protejan} (forme \esp{ustedes}) et \esp{¡cuidémosla!} (forme \esp{nosotros} = « prenons-en soin », avec le pronom \esp{la} soudé après le verbe).
\item \esp{mosquitera impregnada} = moustiquaire imprégnée (d'insecticide) ; \esp{agua estancada} = eau stagnante ; \esp{centro de salud} = centre de santé.
\end{enumerate}

\emph{Barème indicatif (sur 20) :} compréhension 5 $\times$ 2 = 10 pts ; vrai/faux justifié 2 $\times$ 2 = 4 pts ; questions de langue 3 $\times$ 2 = 6 pts. Toute réponse doit être une phrase complète en espagnol ; une réponse copiée sans justification au vrai/faux ne vaut que la moitié des points.
\end{corrige}

\begin{corrige}[Exercice 10 — Thème et version]
\textbf{A. Thème (proposition de traduction) :}

\esp{Es necesario que los alumnos duerman bajo una mosquitera y conviene que hagan deporte. ¡No fuméis y lavaos las manos antes de cada comida! La salud es muy importante para tener éxito en la escuela. Es necesario que todo el mundo beba agua potable.}

\emph{Points de langue :}
\begin{itemize}
\item \esp{duerman} et \esp{hagan} : subjonctifs irréguliers (\esp{dormir} : o $\rightarrow$ ue ; \esp{hacer} : radical \esp{hag-}).
\item \esp{¡No fuméis!} : impératif négatif, 2\up{e} pers. plur.
\item \esp{lavaos} : impératif affirmatif pronominal à \esp{vosotros} — chute du \esp{-d} : \esp{lavad} + \esp{os} $\rightarrow$ \esp{lavaos}.
\item \esp{beba} : subjonctif après \esp{es necesario que} ; \esp{todo el mundo} (« tout le monde ») est suivi du singulier.
\item « réussir à l'école » : \esp{tener éxito en la escuela} (ou \esp{triunfar en la escuela}).
\end{itemize}

\textbf{B. Version (traduction) :}

« Chers concitoyens : l'hôpital de Douala organise une campagne de vaccination gratuite du 10 au 15 mai. Il est nécessaire que tous les enfants reçoivent leurs vaccins. Il convient que les parents apportent le carnet de santé de leurs enfants. Venez tôt et n'ayez pas peur : le vaccin ne fait pas mal ! Il est important que nous protégions ensemble notre communauté contre les maladies. »

\emph{Points de langue :} \esp{reciban}, \esp{traigan}, \esp{protejamos} : subjonctifs après des expressions impersonnelles ; \esp{¡vengan!}, \esp{¡no tengan!} : impératifs de politesse (\esp{ustedes}), formes irrégulières de \esp{venir} et \esp{tener} ; \esp{duele} : 3\up{e} pers. du singulier de \esp{doler} (« faire mal »), verbe à changement radical o $\rightarrow$ ue, construit comme \esp{gustar} : \esp{la vacuna no duele} = « le vaccin ne fait pas mal ».

\emph{Barème indicatif (sur 20) :} thème 10 pts (2 pts par phrase : mode correct 1 pt, vocabulaire et orthographe 1 pt) ; version 10 pts (fidélité du sens 6 pts, qualité du français 4 pts).
\end{corrige}

\begin{corrige}[Exercice 11 — Expression écrite : un message de prévention]
\textbf{Proposition de rédaction modèle (environ 95 mots) :}

\esp{¡Atención, jóvenes del liceo de Nkongsamba! El cólera es una enfermedad peligrosa que se transmite por el agua y los alimentos contaminados. Para protegerte, sigue estos consejos: ¡lávate las manos con jabón antes de cada comida! ¡No bebas agua del río! Es necesario que hiervas el agua antes de beberla. Conviene que laves bien las frutas y las verduras. Si tienes diarrea o vómitos, ve rápidamente al centro de salud. La prevención es la mejor medicina: ¡juntos podemos proteger nuestro liceo y nuestras familias!}

\emph{Traduction :} « Attention, jeunes du lycée de Nkongsamba ! Le choléra est une maladie dangereuse qui se transmet par l'eau et les aliments contaminés. Pour te protéger, suis ces conseils : lave-toi les mains avec du savon avant chaque repas ! Ne bois pas l'eau de la rivière ! Il faut que tu fasses bouillir l'eau avant de la boire. Il convient que tu laves bien les fruits et les légumes. Si tu as des diarrhées ou des vomissements, va vite au centre de santé. La prévention est le meilleur médicament : ensemble, nous pouvons protéger notre lycée et nos familles ! »

\textbf{Barème indicatif (sur 20) :}
\begin{itemize}
\item Respect de la tâche (accroche, transmission, quatre conseils, encouragement) : 6 pts.
\item Correction grammaticale (impératifs et subjonctifs bien formés) : 6 pts.
\item Richesse et exactitude du lexique de la santé : 4 pts.
\item Orthographe, accents, ponctuation espagnole (¡ \ldots !) : 2 pts.
\item Longueur respectée (80--100 mots) et présentation : 2 pts.
\end{itemize}

\textbf{Points de vigilance :}
\begin{itemize}
\item \esp{El cólera} est masculin et prend un accent sur le \esp{ó} ; ne pas écrire \esp{la colera}.
\item « Faire bouillir » : \esp{hervir} (e $\rightarrow$ ie) : \esp{¡Hierve el agua!}, \esp{es necesario que hiervas el agua}.
\item Après \esp{es necesario que} / \esp{conviene que}, toujours le subjonctif : \esp{es necesario que te laves}, jamais \esp{te lavas}.
\item Point d'exclamation ouvrant obligatoire : \esp{¡Lávate las manos!}
\item Impératif pronominal affirmatif : le pronom se soude après le verbe avec accent écrit : \esp{lávate}, \esp{protégete}.
\end{itemize}
\end{corrige}

\newpage

\coursec{Fiche bilan}

\begin{popfigure}
\begin{tikzpicture}[mindmap, grow cyclic, every node/.style={concept, text=white, font=\small\bfseries, align=center}, root concept/.append style={concept color=popblue, font=\large\bfseries}, level 1 concept/.append style={level distance=3.6cm, sibling angle=60}, level 2 concept/.append style={level distance=2.2cm, font=\scriptsize\bfseries}]
\node[root concept]{La salud y la prevención}
  child[concept color=poporange]{ node {Lexique de la santé}
    child[concept color=poporange!70]{ node[align=center]{enfermedad\\síntoma, dolor} }
    child[concept color=poporange!70]{ node[align=center]{vacuna\\higiene} }
  }
  child[concept color=popgreen]{ node {Habitos de vida}
    child[concept color=popgreen!70]{ node[align=center]{alimentación\\equilibrada} }
    child[concept color=popgreen!70]{ node[align=center]{ejercicio\\descanso} }
  }
  child[concept color=poppurple]{ node {Impératif}
    child[concept color=poppurple!70]{ node {tú : -a / -e} }
    child[concept color=poppurple!70]{ node {vosotros : -ad / -ed / -id} }
  }
  child[concept color=popgold]{ node {Subjonctif}
    child[concept color=popgold!70]{ node[align=center]{es necesario\\que + subj.} }
    child[concept color=popgold!70]{ node[align=center]{conviene\\que + subj.} }
  }
  child[concept color=poppink]{ node {Maladies cibles}
    child[concept color=poppink!70]{ node[align=center]{paludismo\\cólera} }
    child[concept color=poppink!70]{ node {VIH / sida} }
  }
  child[concept color=popblue!80]{ node {Message de prévention}
    child[concept color=popblue!60]{ node[align=center]{slogan + conseils\\+ appel à l'action} }
  };
\end{tikzpicture}
\legende{Carte mentale de la leçon : santé, prévention et outils pour conseiller}
\end{popfigure}

\begin{bilan}
\textbf{1. Lexique clé de la santé (à connaître par cœur)}
\begin{itemize}
  \item \esp{la salud} : la santé ; \esp{la enfermedad} : la maladie ; \esp{estar sano / enfermo} : être en bonne santé / malade.
  \item \esp{el paludismo / la malaria} : le paludisme ; \esp{el cólera} : le choléra ; \esp{el VIH / el sida} : le VIH / le sida.
  \item \esp{la vacuna} : le vaccin ; \esp{vacunarse} : se faire vacciner ; \esp{la higiene} : l'hygiène.
  \item \esp{la alimentación equilibrada} : une alimentation équilibrée ; \esp{hacer ejercicio} : faire de l'exercice.
  \item \esp{el dolor} : la douleur ; \esp{el síntoma} : le symptôme ; \esp{la fiebre} : la fièvre ; \esp{el médico} : le médecin.
  \item Verbes utiles : \esp{prevenir} (prévenir), \esp{protegerse} (se protéger), \esp{lavarse las manos} (se laver les mains), \esp{dormir bajo un mosquitero} (dormir sous une moustiquaire), \esp{hervir el agua} (faire bouillir l'eau), \esp{consultar al médico} (consulter le médecin).
\end{itemize}

\textbf{2. Grammaire à connaître par cœur}
\begin{center}
\begin{tabularx}{\linewidth}{|X|X|X|}
\hline
\rowcolor{popblueL} Règle & Formation & Exemple \\
\hline
Impératif affirmatif (tú) & 3\textsuperscript{e} pers. du singulier du présent & \esp{Come frutas.} « Mange des fruits. » \\
\hline
Impératif (vosotros) & infinitif : -r $\rightarrow$ -d & \esp{Bebed agua potable.} « Buvez de l'eau potable. » \\
\hline
Impératif négatif & no + subjonctif présent & \esp{No fumes.} « Ne fume pas. » \\
\hline
Impératifs irréguliers (tú) & decir $\rightarrow$ di ; hacer $\rightarrow$ haz ; ir $\rightarrow$ ve ; poner $\rightarrow$ pon ; salir $\rightarrow$ sal ; ser $\rightarrow$ sé ; tener $\rightarrow$ ten ; venir $\rightarrow$ ven & \esp{Sé prudente.} « Sois prudent. » \\
\hline
\esp{es necesario que} + subjonctif & subjonctif : -ar $\rightarrow$ -e ; -er/-ir $\rightarrow$ -a & \esp{Es necesario que te vacunes.} « Il faut que tu te fasses vacciner. » \\
\hline
\esp{conviene que} + subjonctif & même construction impersonnelle & \esp{Conviene que duermas bajo un mosquitero.} « Il convient que tu dormes sous une moustiquaire. » \\
\hline
\end{tabularx}
\end{center}

\textbf{3. Méthodes express}
\begin{itemize}
  \item \textbf{Analyser des habitudes de vie} : repérer dans le texte les bonnes habitudes (alimentation, hygiène, sport, sommeil) et les mauvaises (tabac, alcool, eau non potable), puis les classer et les commenter.
  \item \textbf{Rédiger un message de prévention} en 4 étapes : 1) un slogan court et frappant ; 2) 3 à 5 conseils à l'impératif ; 3) une phrase avec \esp{es necesario que} ou \esp{conviene que} + subjonctif ; 4) un appel final à l'action (\esp{¡Protege tu salud!}).
  \item \textbf{Traduire un conseil} : « il faut + infinitif » se rend par \esp{hay que + infinitif} ou \esp{es necesario que + subjonctif} ; ne jamais calquer « il faut que tu fais » par un indicatif.
\end{itemize}
\end{bilan}

\begin{aretenir}[Checklist avant le Bac]
\begin{itemize}
  \item[$\square$] Je connais le lexique de la santé : \esp{salud, enfermedad, síntoma, dolor, vacuna, higiene}.
  \item[$\square$] Je sais nommer les trois maladies cibles : \esp{paludismo, cólera, VIH/sida}.
  \item[$\square$] Je forme l'impératif affirmatif de tú et de vosotros sans hésiter.
  \item[$\square$] Je connais les huit impératifs irréguliers : \esp{di, haz, ve, pon, sal, sé, ten, ven}.
  \item[$\square$] Je sais que l'impératif négatif exige le subjonctif : \esp{No bebas esa agua.}
  \item[$\square$] Je forme le subjonctif présent des verbes réguliers (-ar $\rightarrow$ -e ; -er/-ir $\rightarrow$ -a).
  \item[$\square$] J'emploie correctement \esp{es necesario que} et \esp{conviene que} + subjonctif.
  \item[$\square$] Je distingue \esp{hay que + infinitif} (général) et \esp{es necesario que + subjonctif} (personne précise).
  \item[$\square$] Je place correctement les pronoms avec l'impératif : \esp{lávate}, \esp{protégete} (accent !).
  \item[$\square$] Je sais structurer un message de prévention : slogan, conseils, appel à l'action.
  \item[$\square$] Je n'oublie pas les signes ¡ ! et ¿ ? ni les accents dans mes productions.
  \item[$\square$] Je peux analyser les habitudes de vie d'un texte et les juger (favorables / défavorables).
\end{itemize}
\end{aretenir}

\findecours

\end{document}
